\chapter{Complex pluripotential reference sheet}
\index{complex pluripotential theory!reference sheet}\index{nonpluripolar product!complex}
\label{app:complex-reference}

This appendix collects the complex pluripotential formulas used in
\cref{part:complex,part:comparisons}. It is a working reference, not a second development of the
subject: proofs of the local foundations and the singularity comparison
theorems are given in the supplied book \cite{XiaBook}, while
\Cref{chap:complex-prerequisites,chap:rays-test-curves} establish the
interfaces needed here.

Except in the local-domain settings of
\Cref{sec:appH-local-maximality,sec:appH-smooth-integral-comparison,sec:appH-boundary-envelopes,sec:appH-unbounded-domain-comparison,sec:appH-Bergman-rigidity,sec:appH-ball-quotients,sec:appH-bounded-minorants,sec:appH-toric-multiplier-pairs},
the family setting of \Cref{sec:appH-KE-family-normalization},
the local analytic germ setting of \Cref{sec:appH-finite-pullback-Segre},
and the smooth curvature calculations of \Cref{sec:appH-Griffiths-Chern}, \(X\) is a connected compact
K\"ahler manifold of complex dimension \(n\). We use
\[
             \dc=\frac{\rmi}{4\pi}(\bar\partial-\partial),
             \qquad
             \ddc=\frac{\rmi}{2\pi}\partial\bar\partial .
\]
Thus
\[
             \ddc\log|f|^2=[\operatorname{div}(f)]
\]
for a nonzero meromorphic function \(f\). Positive products are
nonpluripolar products. They are not divided by the volume of a cohomology
class. The local Cegrell products in
\Cref{sec:appH-local-maximality} are explicitly full products and can
charge pluripolar sets.
The spectral expressions in \Cref{sec:appH-smooth-integral-comparison}
are classical smooth coordinate operators; their signed polynomial
combinations are not assigned a nonpluripolar interpretation.
The smooth Chern--Weil forms in \Cref{sec:appH-Griffiths-Chern}
are also ordinary differential forms; the top form there may be signed.
The constant-form calculations of \Cref{sec:appH-Nakano-determinants}
use the full exterior spaces of $\mathbf C^n$ as an additional local
setting, with no compactness assumption.

\section{Potentials and the three singularity orders}
\label{sec:appH-singularity-orders}

Let \(\theta\) be a smooth closed real \((1,1)\)-form representing a
pseudoeffective class.

\begin{definition}[Potentials and currents]
\label{def:appH-potentials}
A quasi-psh function \(u:X\to[-\infty,+\infty)\), not identically
\(-\infty\), lies in \(\PSH(X,\theta)\) if
\[
                         \theta_u:=\theta+\ddc u\geq0
\]
as a current. Here quasi-psh means locally the sum of a psh function and
a smooth function, with that pointwise representative. If
\(\theta'=\theta+\ddc f\) with \(f\) smooth, then
\[
       u\in\PSH(X,\theta)
       \quad\Longleftrightarrow\quad
       u-f\in\PSH(X,\theta'),
\]
and the two potentials represent the same current.
\end{definition}

The potential with minimal singularities is
\[
 V_\theta
 =\uscsup\{v\in\PSH(X,\theta)\mid v\leq0\}.
\]
It is maximal only with respect to singularity: every
\(u\in\PSH(X,\theta)\) satisfies \(u\preceq V_\theta\). It need not be
smooth or locally bounded when \(\{\theta\}\) is merely big.

\begin{definition}[Ordinary singularity type]
\label{def:appH-ordinary-type}
For quasi-psh functions \(u,v\), write
\[
 u\preceq v
 \quad\Longleftrightarrow\quad
 u\leq v+C
 \quad\text{for some }C\in\mathbf R.
\]
Thus \(u\preceq v\) means that \(u\) is \emph{more singular} than \(v\).
Write \(u\sim v\) when both \(u\preceq v\) and \(v\preceq u\) hold.
Equivalently, their difference is bounded on the common finite locus and
the two comparison inequalities hold everywhere.  This formulation avoids
subtracting \(-\infty\) from \(-\infty\).
\end{definition}

There are two coarser comparisons.

\begin{definition}[\(P\)- and \(\mathcal I\)-comparison]
\label{def:appH-p-i-orders}
After placing \(u\) and \(v\) in a common sufficiently positive background,
write
\[
\begin{aligned}
 u\preceq_Pv
 &\quad\Longleftrightarrow\quad
 P_\theta[u]\leq P_\theta[v],\\
 u\preceq_{\mathcal I}v
 &\quad\Longleftrightarrow\quad
 \mathcal I(m u)\subseteq\mathcal I(m v)
 \quad\text{for every }m\geq1.
\end{aligned}
\]
The second condition is equivalent to
\[
                         \nu(u,E)\geq\nu(v,E)
\]
for every prime divisor \(E\) over \(X\).
\end{definition}

The direction of every arrow in the comparison hierarchy is
\[
            u\preceq v
            \Longrightarrow
            u\preceq_Pv
            \Longrightarrow
            u\preceq_{\mathcal I}v.
\]
No reverse implication is automatic. In particular, equality of all
divisorial Lelong numbers need not imply bounded difference.

\begin{center}
\begin{tabular}{lll}
\toprule
relation & invariant retained & information forgotten\\
\midrule
\(\sim\) & full singularity modulo \(O(1)\) & bounded perturbations\\
\(\sim_P\) & model envelope and nonpluripolar mass & fine submodel behavior\\
\(\sim_{\mathcal I}\) & all multiplier ideals/divisorial orders &
nondivisorial behavior\\
\bottomrule
\end{tabular}
\end{center}

\begin{example}[A zero-Lelong singularity]
\label{ex:appH-loglog}
On \(0<|z|<e^{-1}\), the function
\[
                         u(z)=-\log\bigl(-\log|z|^2\bigr)
\]
is subharmonic: as a function of \(s=\log|z|^2<0\), its first and second
derivatives are \(-1/s>0\) and \(1/s^2>0\).  Moreover,
\(u/\log|z|^2\to0\) at the origin, so its Lelong number is zero although
it is unbounded below.  Patch it smoothly away from a smaller disk and
multiply by a sufficiently small positive constant to obtain an
\(\omega_{\FS}\)-psh function on \(\mathbf P^1\).  Its only unbounded
point has no atom in its curvature measure; in dimension one the
nonpluripolar product therefore retains the entire cohomological mass.
By \Cref{thm:appH-p-projection}, its normalized model envelope is
\(V_{\omega_{\FS}}=0\).  It is \(P\)-equivalent, hence also
\(\mathcal I\)-equivalent, to zero, but not ordinarily equivalent to zero.
\end{example}

\section{Nonpluripolar products}
\label{sec:appH-npp}

Let \(T_j=\theta_j+\ddc\varphi_j\) be closed positive \((1,1)\)-currents,
\(1\leq j\leq p\), and on a coordinate open set write \(T_j=\ddc u_j\) with
\(u_j\) \emph{psh} --- a local potential, not the global
\(\theta_j\)-psh \(\varphi_j\), for which \(\ddc\varphi_j=T_j-\theta_j\) is
not positive and \(\max\{\varphi_j,-k\}\) is not psh.  Put
\[
 O_k=\bigcap_{j=1}^p\{u_j>-k\}.
\]
Bedford--Taylor theory defines the bounded truncation product.  On the
plurifine open set \(O_k\), locality identifies its restriction with the
product of the finite local potentials; the notation on the left below
means this restriction, not an unrestricted product of unbounded currents:
\[
 \one_{O_k}\bigwedge_{j=1}^p\ddc u_j
 =
 \one_{O_k}\bigwedge_{j=1}^p\ddc\max\{u_j,-k\}.
\]
The right side increases with \(k\). Extending it by zero and passing to
the limit produces
\[
                    \langle T_1\wedge\cdots\wedge T_p\rangle.
\]

\begin{theorem}\label{thm:appH-npp-ledger}
The nonpluripolar product is a closed positive \((p,p)\)-current. It is:
\begin{enumerate}
\item symmetric and multilinear on the cone of positive currents;
\item local for the plurifine topology;
\item equal to the Bedford--Taylor product for locally bounded potentials;
\item zero on every pluripolar set;
\item monotone in mixed mass, after wedging with a fixed K\"ahler form
to complementary degree, when any potential is replaced by a less
singular one in the same cohomology class.
\end{enumerate}
For a proper bimeromorphic morphism \(\pi:Y\to X\) of compact K\"ahler
manifolds, its pullback construction satisfies
\[
 \pi_*\langle\pi^*T_1\wedge\cdots\wedge\pi^*T_p\rangle
       =\langle T_1\wedge\cdots\wedge T_p\rangle.
\]
Consequently the integrals against corresponding smooth closed test forms
of complementary degree agree.  For \(p=n\) these are the total masses.
\end{theorem}

\begin{proof}[From the nonpluripolar-product inputs]
Apply \Cref{thm:complex-npp-ledger} to the indicated currents.
The global local-mass bound needed for the construction is
\cite[Proposition~1.6]{BEGZ10}; closedness is
\cite[Theorem~1.8]{BEGZ10}.  Restrictions by \(\one_{O_k}\) need not be
closed, so closedness of the limit cannot be obtained by applying Stokes
to those restricted truncations.  Positivity, symmetry, plurifine
locality and multilinearity on the cone of positive currents are the
remaining inputs recorded in that internal theorem.  The unrestricted
mixed-mass monotonicity is \cite[Theorem~1.1]{DDNL18mono}, with the fixed
complementary K\"ahler factors included as additional currents.
The modification identity is \eqref{eq:complex-npp-modification}; pairing
it with a complementary smooth closed form and applying the projection
formula gives the stated equality of integrals.
\end{proof}

The nonpluripolar construction is due to \cite{BEGZ10}; the unrestricted
monotonicity assertion in the theorem is due to \cite{DDNL18mono}.  It is a
statement about total masses, not a pointwise inequality between the two
product measures.

\begin{definition}[Monge--Amp\`ere measure and mass]
\label{def:appH-ma-mass}
For \(u\in\PSH(X,\theta)\), set
\[
             \operatorname{MA}_\theta(u):=\langle\theta_u^n\rangle,
             \qquad
             M_\theta(u):=\int_X\operatorname{MA}_\theta(u).
\]
\end{definition}

The measure \(\operatorname{MA}_\theta(u)\) is unnormalized. If a
probability measure is desired and \(M_\theta(u)>0\), it is
\[
                         M_\theta(u)^{-1}
                         \operatorname{MA}_\theta(u).
\]
When the mass is zero, this normalization is undefined.

\paragraph{Safe continuity rules.}\label{prop:appH-safe-continuity}
The following rules may be used without further qualification.
\begin{enumerate}
\item If \(u_i\) is an increasing net in \(\PSH(X,\theta)\), uniformly
bounded above, and \(u=\uscsup_i u_i\), then
\[
                         M_\theta(u)=\lim_iM_\theta(u_i).
\]
\item In a big class, if \(u_j\searrow u\) and \(u_j\sim u\) for every
\(j\), then \(\operatorname{MA}_\theta(u_j)\) converges weakly to
\(\operatorname{MA}_\theta(u)\).  The limit is explicitly included in
the common ordinary singularity type.
\item For an arbitrary decreasing family, neither preservation of mass nor
weak convergence of Monge--Amp\`ere measures is a formal consequence of
monotonicity.
\end{enumerate}

The first item follows from Choquet reduction and the sequential
nonpluripolar continuity theorem; its proof for genuine nets is
\Cref{prop:complex-mass-continuity}.  For the second, decreasing psh
convergence gives capacity convergence and the common type gives equality
of total masses.  These are precisely the hypotheses of
\cite[Theorem~2.3]{DDNL18mono}.  Continuity in a relative finite-energy
class has its own hypotheses; it is not an assertion about an arbitrary
decreasing family.

\section{Relative full mass and finite energy}
\label{sec:appH-relative-energy}

Fix a positive-mass model potential
\[
                         \phi=P_\theta[\phi],
             \qquad M_\theta(\phi)>0.
\]
The reference singularity is part of the data.

\begin{definition}[Relative full-mass classes]
\label{def:appH-relative-full-mass}
Set
\[
\begin{aligned}
 \PSH(X,\theta,\phi)
 &:=\{u\in\PSH(X,\theta)\mid u\preceq\phi\},\\
 \mathcal E(X,\theta,\phi)
 &:=\{u\in\PSH(X,\theta,\phi)\mid
                   M_\theta(u)=M_\theta(\phi)\}.
\end{aligned}
\]
\end{definition}

Thus a member of \(\mathcal E(X,\theta,\phi)\) may be strictly more
singular than \(\phi\), but loses no nonpluripolar mass relative to
\(\phi\).

If \(u\sim\phi\), define the unnormalized relative Aubin--Mabuchi
energy by
\[
 E_\phi(u)
 :=\frac1{n+1}\sum_{j=0}^n
   \int_X (u-\phi)
   \langle\theta_u^j\wedge\theta_\phi^{\,n-j}\rangle.
\]
The difference in the integral is taken on the common finite locus and
extended arbitrarily across the pluripolar set, which the measures do not
charge.  For \(u\preceq\phi\), choose \(C\) with \(u\leq\phi+C\) and put
\[
                         u_k=\max\{u,\phi-k\}.
\]
Then \(\phi-k\leq u_k\leq\phi+C\), so each \(u_k\sim\phi\).  Set
\[
                         E_\phi(u)
 :=\lim_{k\to\infty}E_\phi(u_k)\in[-\infty,+\infty).
\]

\begin{definition}[Finite relative energy]
\label{def:appH-e1}
The relative finite-energy class is
\[
 \mathcal E^1(X,\theta,\phi)
 :=\{u\in\mathcal E(X,\theta,\phi)\mid E_\phi(u)>-\infty\}.
\]
\end{definition}

\begin{proposition}\label{prop:appH-energy-formulas}
Let \(u,v\in\mathcal E^1(X,\theta,\phi)\). The following identities hold,
first for bounded relative difference and then by monotone approximation:
\begin{enumerate}
\item translation:
\[
                         E_\phi(u+c)
                         =E_\phi(u)+cM_\theta(\phi);
\]
\item monotonicity:
\[
                         u\leq v\Longrightarrow
                         E_\phi(u)\leq E_\phi(v);
\]
\item cocycle formula (all the displayed mixed integrals are finite):
\[
 E_\phi(u)-E_\phi(v)
 =\frac1{n+1}\sum_{j=0}^n
   \int_X(u-v)
   \langle\theta_u^j\wedge\theta_v^{\,n-j}\rangle;
\]
\item first variation along an admissible bounded direction \(f\):
\[
 \left.\frac{\rd}{\rd t}\right|_{t=0}E_\phi(u+tf)
 =\int_X f\,\operatorname{MA}_\theta(u).
\]
\end{enumerate}
\end{proposition}

\begin{proof}[From relative integration by parts]
For \(u\sim\phi\) and \(v\sim\phi\), all mixed masses equal
\(M_\theta(\phi)\), by \cite[Propositions~3.1.3 and 3.1.14]{XiaBook}.
Translation leaves the currents unchanged, so its contribution to the
polarized sum is exactly \(cM_\theta(\phi)\).  The relative
nonpluripolar integration-by-parts formula in
\cite[Lemma~3.1.3, equation~(3.32)]{XiaBook} gives the cocycle identity.
Its positive mixed measures prove monotonicity.  Bounded relative
difference does not make \(\phi\) locally bounded, so ordinary bounded
Bedford--Taylor integration by parts alone would not justify this step.

For general \(u,v\in\mathcal E^1\), mixed integrability is
\cite[Corollary~3.1.6]{XiaBook}.  The relative truncation-tail estimates
in \cite[Proposition~3.1.18, equations~(3.40)--(3.43)]{XiaBook} pass the
cocycle identity to the limit; this requires control of the weighted
integrals, not just weak convergence of measures.  Translation follows
from \cite[equation~(3.31)]{XiaBook}, and monotonicity again follows
from the cocycle formula.  An admissible bounded direction means that
\(u+tf\) remains in the relative psh domain for \(|t|\) small.  These
potentials remain finite energy; apply the affine-segment variation
formula \cite[Lemma~3.1.11]{XiaBook} to obtain the fourth assertion.
\end{proof}

If one divides \(E_\phi\) by \(M_\theta(\phi)\), the translation formula
has coefficient one. This book does not silently switch between these two
normalizations.

\begin{definition}[The \(I\)-functional and \(d_1\)]
\label{def:appH-i-d1}
For finite-energy potentials of the same relative full-mass class, set
\[
 I_\phi(u,v)
 :=\int_X(u-v)
     \bigl(\operatorname{MA}_\theta(v)
           -\operatorname{MA}_\theta(u)\bigr)\geq0.
\]
Let
\[
 P_\theta(u,v)
 :=\uscsup\{w\in\PSH(X,\theta)\mid w\leq u,\ w\leq v\}.
\]
Then
\[
 d_{1,\phi}(u,v)
 :=E_\phi(u)+E_\phi(v)-2E_\phi(P_\theta(u,v))
\]
is the relative \(d_1\)-metric on \(\mathcal E^1(X,\theta,\phi)\).
The nonnegativity of \(I_\phi\) follows by adding the two endpoint energy
inequalities in \cite[Proposition~3.1.19]{XiaBook}.
\end{definition}

The relative metric and its completeness are
\cite[Lemma~4.3.1 and Theorem~4.3.2]{XiaBook}; the rooftop remains in
\(\mathcal E^1\) by \cite[Proposition~3.1.21]{XiaBook}.
Convergence in this metric implies \(L^1\)-convergence and energy
convergence, by \cite[Lemma~4.3.5 and Theorem~4.3.3]{XiaBook}.
For a pointwise monotone sequence with limit in the same relative
finite-energy class, an exact criterion is
\[
 u_j\longrightarrow u\text{ in }d_{1,\phi}
 \quad\Longleftrightarrow\quad
 E_\phi(u_j)\longrightarrow E_\phi(u).
\]
Indeed, the two potentials are ordered, so their rooftop is the smaller
one and their distance is the absolute energy difference.  The
singularity-type metric \(d_S\) of \cite{DDNLmetric} is a different metric.

\begin{theorem}\label{thm:appH-comparison}
If \(u,v\in\mathcal E(X,\theta,\phi)\), then
\[
 \int_{\{u<v\}}\operatorname{MA}_\theta(v)
 \leq
 \int_{\{u<v\}}\operatorname{MA}_\theta(u).
\]
The relative domination principle also holds: if
\(\operatorname{MA}_\theta(u)(\{u<v\})=0\), then \(u\geq v\) everywhere.
\end{theorem}

\begin{proof}[From the relative comparison input]
Subtracting one common constant, we may suppose \(u,v\leq\phi\).  Put
\(u_j=\max\{u,\phi-j\}\) and
\(v_j=\max\{v,\phi-j\}\).  Their differences from \(\phi\) are bounded,
so the relative comparison input \cite[Proposition~3.1.15]{XiaBook}
gives, for every \(\delta>0\),
\[
 \int_{\{u_j<v_j-\delta\}}\operatorname{MA}_\theta(v_j)
 \leq
 \int_{\{u_j<v_j-\delta\}}\operatorname{MA}_\theta(u_j).
\tag{H.1}\label{eq:appH-truncated-comparison}
\]
On
\(A_j:=\{u<v-2\delta\}\cap\{u>\phi-j\}\cap\{v>\phi-j\}\),
plurifine locality identifies the left integrand with
\(\operatorname{MA}_\theta(v)\), and
\(A_j\subset\{u_j<v_j-\delta\}\).  On the right, the part of the comparison
set where \(u>\phi-j\) lies in \(\{u<v\}\) and carries
\(\operatorname{MA}_\theta(u)\).  The remaining mass tends to zero: the
truncation measures have total mass \(M_\theta(\phi)\), their restrictions
to \(\{u>\phi-j\}\) increase to
\(\operatorname{MA}_\theta(u)\), and \(u\) has that same full mass.
Letting \(j\to\infty\) in \eqref{eq:appH-truncated-comparison} therefore
gives
\[
 \int_{\{u<v-2\delta\}}\operatorname{MA}_\theta(v)
 \leq
 \int_{\{u<v\}}\operatorname{MA}_\theta(u).
\]
The possible set where a potential is \(-\infty\) is pluripolar and is not
charged.  Finally let \(\delta\downarrow0\) and use monotone convergence.
The domination assertion is the separate positive-mass input
\cite[Theorem~2.4.5 and Remark~2.4.3]{XiaBook}.  Its hypotheses are the
common relative full mass and the vanishing mass of \(\{u<v\}\) stated
here; the comparison inequality alone does not prove domination.
\end{proof}

This theorem is not asserted for arbitrary \(\theta\)-psh functions of
unequal mass. Full mass is the mechanism that eliminates the error carried
by the singular sets.

\section{Rooftops, model envelopes, and ideal envelopes}
\label{sec:appH-envelopes}

For \(u,v\in\PSH(X,\theta)\), the rooftop envelope is
\[
 P_\theta(u,v)
 =\uscsup\{w\in\PSH(X,\theta)\mid w\leq u,\ w\leq v\}.
\]
For \(u\in\PSH(X,\theta)\), the singularity envelope is
\[
 P_\theta[u]
 =\uscsup\{w\in\PSH(X,\theta)\mid
             w\leq0,\ w\preceq u\}.
\]
Equivalently,
\[
 P_\theta[u]
 =\uscsup_{C>0}P_\theta(u+C,V_\theta).
\]

\begin{theorem}\label{thm:appH-p-projection}
For every \(u\in\PSH(X,\theta)\),
\[
 M_\theta(P_\theta[u])=M_\theta(u).
\]
If \(\{\theta\}\) is big, then \(P_\theta[P_\theta[u]]=P_\theta[u]\).
If moreover \(M_\theta(u)>0\), then \(P_\theta[u]\) is the largest
normalized potential that is less singular than \(u\) and has the same
nonpluripolar mass, and idempotence follows from that characterization.
\end{theorem}

\begin{proof}[From the model-envelope inputs]
Apply \Cref{thm:complex-p-projection}.  To recall the mass argument, put
\(a=\sup_Xu\).  For \(C\) sufficiently large, the rooftop
\(w_C=P_\theta(u+C,V_\theta)\) satisfies
\[
                   u-a\leq w_C\leq u+C.
\]
The first inequality holds because \(u-a\leq V_\theta\) and
\(u-a\leq u+C\).  Hence \(w_C\sim u\) globally, and its mass equals
\(M_\theta(u)\).  Increasing mass continuity gives the mass identity
for the envelope.  The positive-mass maximality characterization is
\eqref{eq:complex-p-characterization}, with its proof in that internal
theorem.  It gives idempotence at positive mass.  The zero-mass big-class
case is the additional input \cite[Theorem~1.2]{PSWZ26}; it is not a
consequence of positive-mass domination.
\end{proof}

Mass preservation needs no positivity hypothesis, and in a big class
idempotence holds at zero mass as well, by \cite[Theorem~1.2]{PSWZ26}; see
\Cref{warn:ch51-p-positive-mass} for exactly what that input does and does
not supply.  Positive mass is needed for the maximality characterization,
for the relative class \(\mathcal E^1(X,\theta,P_\theta[u])\), and for the
probability normalization.  \Cref{chap:transcendental-metrics} handles zero
volume through a directed system of positive perturbations for that
reason.

The ideal envelope is
\[
 P_\theta[u]_{\mathcal I}
 :=\uscsup\{w\in\PSH(X,\theta)\mid
           w\leq0,\ w\sim_{\mathcal I}u\}.
\]

\begin{theorem}\label{thm:appH-i-projection}
For every \(u\in\PSH(X,\theta)\),
\[
\begin{aligned}
 u&\sim_{\mathcal I}P_\theta[u]_{\mathcal I},\\
 P_\theta[P_\theta[u]_{\mathcal I}]_{\mathcal I}
   &=P_\theta[u]_{\mathcal I},\\
 P_\theta[u]&\leq P_\theta[u]_{\mathcal I}.
\end{aligned}
\]
If \(u\) has analytic singularities, then
\[
                         u\sim P_\theta[u]
                         \sim P_\theta[u]_{\mathcal I}.
\]
\end{theorem}

\begin{proof}[From the ideal-envelope inputs]
Apply \Cref{thm:complex-i-projection}, which proves all four assertions
without a positive-mass hypothesis.  Its ideal-envelope argument uses
finite maxima, increasing-envelope continuity of divisorial Lelong
numbers, and the valuative criterion to show that the envelope has
exactly the same multiplier ideals as \(u\).  It is then the largest
member of its own competitor family, proving idempotence.

For the inequality between envelopes, competitor-family containment is
insufficient: \(w\preceq u\) need not imply \(w\sim_{\mathcal I}u\).
Instead, the rooftops in the preceding proof have ordinary type \(u\),
and their increasing envelope preserves its divisorial orders, as proved
in Step~3 of \Cref{thm:complex-i-projection}.  Thus
\(P_\theta[u]\sim_{\mathcal I}u\); since it is normalized, it is itself
an ideal-envelope competitor.  This argument does not use big-class
idempotence in a merely pseudoeffective class.

The analytic-singularity assertion is the separate log-resolution and
modelity argument in Step~5 of the same internal theorem.  Equality of
divisorial orders alone would not remove an unbounded zero-Lelong
remainder of an arbitrary quasi-psh competitor.
\end{proof}

The final assertion uses a log resolution: among analytic singularities,
equality of divisorial orders controls the singularity up to \(O(1)\).
It is false as a formal principle for arbitrary quasi-psh singularities.

\begin{center}
\begin{tabular}{lll}
\toprule
operator & defining constraint & principal output\\
\midrule
\(P_\theta(u,v)\) & lie below two obstacles & rooftop/contact envelope\\
\(P_\theta[u]\) & ordinary type bounded by \(u\) & mass/model projection\\
\(P_\theta[u]_{\mathcal I}\) & same multiplier ideals & divisorial projection\\
\bottomrule
\end{tabular}
\end{center}

\section{Multiplier ideals and regularization}
\label{sec:appH-multiplier-ideals}

For a quasi-psh function \(u\),
\[
 \mathcal I(u)(U)
 =\{f\in\mathcal O_X(U)\mid |f|^2\rme^{-u}
                       \in L^1_{\mathrm{loc}}(U)\}.
\]
The exponent is \(\rme^{-u}\). With our normalization,
\[
 \mathcal I\bigl(\lambda\log|z|^2\bigr)
 =\mathcal O_{\mathbf C}
   \bigl(-\lfloor\lambda\rfloor\{0\}\bigr).
\]

\begin{theorem}\label{thm:appH-strong-openness}
For a quasi-psh function \(u\),
\[
 \mathcal I(u)=\bigcup_{\epsilon>0}\mathcal I((1+\epsilon)u),
\]
and, for every prime divisor \(E\) over \(X\),
\[
 \nu(u,E)
 =\lim_{m\to\infty}\frac1m
                   \ord_E\mathcal I(m u).
\]
The collection of all multiplier ideals
\(\mathcal I(\lambda u)\), \(\lambda>0\), is equivalent to the collection
of all divisorial Lelong numbers.
\end{theorem}

\begin{proof}[From strong openness and the valuative criterion]
The first identity is the strong openness theorem
\cite[Theorem~1.1]{GZ15}, in the normalization of
\Cref{thm:complex-strong-openness}.  The second identity uses Demailly's approximation theorem
\cite{Dem92reg}, as recorded in \Cref{thm:complex-strong-openness},
item \cref{itm:ch51-so:lelong}.  The valuative integrability lower bound
alone is not a two-sided approximation estimate.  Conversely, reconstruction
uses the following exact local valuative
criterion in the book's squared-norm normalization.  A holomorphic germ
\(f\) belongs to \(\mathcal I(\lambda u)_x\) if and only if there are a
neighborhood \(U_x\) and \(\eta_x>0\) such that, for every prime divisor
\(E\) over \(U_x\),
\[
 \ord_E(f)\geq(1+\eta_x)\lambda\nu(u,E)-A(E).
\]
This is \cite[Corollary~10.18]{Bou17} after replacing the source's weight by
\(\lambda u/2\) and translating its \(\log|z|\) Lelong normalization to
the book's \(\log|z|^2\) normalization.  Thus the divisorial Lelong numbers determine every multiplier
ideal.  Here \(A(E)\) is the log discrepancy.  The single \(\eta_x\),
uniform over divisors on that neighborhood, is essential; mere
pointwise strict inequalities are not a substitute.  Strong openness
removes the boundary ambiguity and shows that ideals at integral scales
determine those at every positive real scale.
\end{proof}

A quasi-equisingular approximation of \(u\) is a decreasing, or suitably
almost decreasing, family \(u_j\) with analytic singularities whose
multiplier ideals stabilize at every fixed scale in the strong-openness
sense.  Its limiting ideal data recover the \(\mathcal I\)-type;
an individual analytic approximant need not have exactly the multiplier
ideals of \(u\).  See \Cref{prop:ch51-qes-preserves-i} for the precise
scaled-ideal stabilization, which does not assert preservation of ordinary
type.

\begin{remark}[Restriction and trace]
\label{rem:appH-restriction-trace}
The pointwise restriction \(u|_Y\) to a submanifold can be identically
\(-\infty\), and different representatives of the same \(P\)- or
\(\mathcal I\)-type can have different naive restrictions. The trace
operator of \Cref{chap:operations-trace} is therefore defined from
quasi-equisingular approximants on the natural domain
\[
                         \nu(u,Y)=0.
\]
Its value is a \(P\)-class. The resulting \(P\)-class map factors through
the coarser \(\mathcal I\)-class, but this does not identify the two
equivalence relations.
\end{remark}

\section{Weak geodesics and rays}
\label{sec:appH-geodesics}

Let \(A=\{1<|\tau|<e\}\), and write
\[
                         t=\log|\tau|\in(0,1).
\]
A subgeodesic joining \(u_0\) to \(u_1\) is an \(S^1\)-invariant
\(\pi_X^*\theta\)-psh function \(U(x,\tau)=u_t(x)\) satisfying
\[
                         \pi_X^*\theta+\ddc_{x,\tau}U\geq0
\]
and the prescribed boundary inequalities. The weak geodesic is the usc
regularization of the supremum of all such subgeodesics.

\begin{theorem}\label{thm:appH-weak-geodesic}
Let \(\phi\) be a positive-mass model potential and let
\(u_0,u_1\in\mathcal E^1(X,\theta,\phi)\). The weak geodesic
\((u_t)_{0\leq t\leq1}\):
\begin{enumerate}
\item belongs to \(\mathcal E^1(X,\theta,\phi)\) at every time;
\item has affine relative energy,
\[
                         E_\phi(u_t)
          =(1-t)E_\phi(u_0)+tE_\phi(u_1);
\]
\item is a constant-speed \(d_{1,\phi}\)-geodesic;
\item solves the homogeneous complex Monge--Amp\`ere equation on
\(X\times A\) in the appropriate nonpluripolar sense.
\end{enumerate}
\end{theorem}

\begin{proof}[From the relative geodesic inputs]
The Perron construction, full mass of its slices and endpoint
\(L^1\)-limits are \cite[Proposition~4.2.1]{XiaBook}, in the scope
recorded by \Cref{thm:rays-relative-geodesic}.  Affine relative energy
and the homogeneous nonpluripolar equation are
\cite[Proposition~4.3.3, equations~(4.33)--(4.35)]{XiaBook}.
The complexified equation is
\[
              \bigl\langle(\pi_X^*\theta+\ddc U)^{n+1}\bigr\rangle=0.
\]
These are imported relative analytic inputs.  Even the truncated endpoints
\(\max\{u_i,\phi-j\}\) can be unbounded in ordinary local coordinates,
so an unrestricted Bedford--Taylor equation would be unjustified.

The rooftop lower bound \(P_\theta(u_0,u_1)\leq u_t\) and the convex
upper bound show that the slices lie in \(\mathcal E^1\), using
\cite[Propositions~3.1.21 and 3.1.22]{XiaBook}.  For the metric assertion,
use convexity of distance along two relative geodesics,
\cite[Proposition~4.3.5]{XiaBook}.  If
\(D=d_{1,\phi}(u_0,u_1)\), comparison with a constant endpoint gives
\(d_{1,\phi}(u_s,u_1)\leq(1-s)D\).  Apply the same convexity on
\([s,1]\), whose restriction is again a Perron geodesic, to obtain
\(d_{1,\phi}(u_s,u_t)\leq(t-s)D\).  The triangle inequality together
with the two endpoint bounds forces equality.  In particular, the
endpoints are attained in \(d_{1,\phi}\).
\end{proof}

The boundary values are interpreted in the finite-energy topology. For
arbitrary singular endpoints one should not infer pointwise boundary
continuity.

\begin{definition}[Geodesic ray]
\label{def:appH-ray}
A finite-energy geodesic ray is a map
\[
                         [0,\infty)\ni t\longmapsto u_t
                         \in\mathcal E^1(X,\theta,\phi)
\]
whose restriction to every compact interval is the weak geodesic between
its endpoints. It is normalized by a condition such as
\[
                         \sup_Xu_t=Ct.
\]
\end{definition}

For a normalized ray, convexity gives well-defined asymptotic slopes for
many convex functionals. Affinity of \(E_\phi\) gives
\[
                         E_\phi(u_t)
                         =E_\phi(u_0)+t\,E_\phi^{\mathrm{slope}}(u_\bullet).
\]
When a non-Archimedean energy is compared with this slope, the total mass
and the choice between norm and squared-norm potentials must be checked
against \Cref{app:normalizations}.

\section{Test curves and Legendre transform}
\label{sec:appH-test-curves}

A test curve is a concave, decreasing family
\(\tau\mapsto\psi_\tau\in\PSH(X,\theta)\cup\{-\infty\}\) that is
\emph{pointwise upper semicontinuous in \(\tau\)}, with a left
endpoint representing the chosen full-mass class and with
\(\psi_\tau\equiv-\infty\) for large \(\tau\).  Every level is required to
be a fixed point of the appropriate \(P\)- or \(\mathcal I\)-envelope;
positive mass is \emph{not} imposed at interior levels.  It is imposed only
at the left endpoint, by \eqref{eq:rays-test-minus-infinity}, and is then
\emph{deduced} at every interior level in
\Cref{lem:rays-test-positive-levels}.  Reading positivity into the
level condition would make that lemma vacuous.  Upper semicontinuity in
\(\tau\) is part of the definition, exactly as in
\Cref{def:appD-test-curve} and \eqref{eq:rays-test-endpoint}: a
concave nonincreasing family is automatically continuous in \(\tau\) on the
interior of its finiteness interval, but it may drop at the right endpoint,
and the inverse transform below then returns the regularized family rather
than the given one.  On \(\mathbf P^1\) with \(\theta=\omega_{\FS}\), the
family \(\psi_\tau=0\) for \(\tau<0\),
\(\psi_0=P_\theta[\tfrac12u]\) with \(u=\log(|z|^2/(1+|z|^2))\), and
\(\psi_\tau=-\infty\) for \(\tau>0\), satisfies every other condition
listed and has \(u_t\equiv0\), so the inverse transform returns \(0\) at
\(\tau=0\) instead of \(P_\theta[\tfrac12u]\).

The forward Legendre transform is
\[
                         u_t
       =\uscsup_{\tau\in\mathbf R}(\psi_\tau+t\tau),
                         \qquad t>0.
\]
The inverse transform is
\[
                         \psi_\tau
       =\inf_{t>0}(u_t-t\tau).
\]
The infimum in the second formula is already upper semicontinuous under
the hypotheses of the Ross--Witt Nystr\"om correspondence; inserting an
extra usc regularization changes the convention and can alter endpoint
statements.

\begin{theorem}\label{thm:appH-rwn}
Fix the positive-mass model reference \(\phi\).  In the precise test-curve
classes of \Cref{def:rays-test-curve}, impose maximality of the levels,
pointwise upper semicontinuity in the parameter and left endpoint \(\phi\).
The displayed Legendre transforms give mutually inverse order-preserving
maps between maximal relative test curves and geodesic rays starting at
\(\phi\).  They restrict to bounded-relative and finite-energy
correspondences.  In the finite-energy case, the test-curve energy
\eqref{eq:rays-test-energy} equals the affine energy slope of the ray,
with the same raw-mass normalization.
\end{theorem}

\begin{proof}[From the relative Ross--Witt Nystr\"om input]
This is \Cref{thm:rays-rwn}, imported from
\cite[Theorem~9.2.1, equations~(9.19)--(9.21)]{XiaBook}.
That theorem supplies the analytic statements that the forward transform
is geodesic, the inverse has maximal levels and the two energy expressions
agree.  Scalar closed-concave biconjugacy alone proves none of those
analytic assertions.  Order preservation follows directly from the
supremum and infimum formulas.  Monotone-limit statements are used in the
specific classes and with the regularizations of that internal theorem;
an arbitrary pointwise limit need not preserve the endpoint conditions.
\end{proof}

See \cite{RWN14} and \Cref{chap:rays-test-curves}. In a big class,
the maximality operation includes the \(P\)-envelope.  Maximality is
required in the ample case as well; ordinary boundedness of a level is
not a replacement for its envelope fixed-point condition.

\section{Local maximality and full residual mass}
\label{sec:appH-local-maximality}
\index{maximal plurisubharmonic function!local-to-global criterion}
\index{analytic spread!maximality criterion}
\index{Siu decomposition!complex surface residual}

We now work on domains in \(\mathbf C^n\), or on a complex manifold
with specified local potentials. The local Cegrell class
\(\mathcal E_{\mathrm{loc}}\) consists of psh functions that, after
subtracting a local constant on a sufficiently small hyperconvex
neighborhood, belong to Cegrell's \(\mathcal E\)-class
\cite{Ceg04}. Its full Monge--Amp\`ere measure is a Radon measure
and can charge poles. This local class differs from the compact
relative full-mass class of \Cref{sec:appH-relative-energy}.
All formulas below use the book's \(\ddc\). The two recent
Do--Hiep papers use \(D=\rmi\partial\bar\partial/\pi=2\ddc\);
replacing every degree-\(n\) product by its book normalization
multiplies both sides of their comparison formulas by \(2^{-n}\).

\subsection{Analytic singularities and analytic spread}
\label{sec:appH-analytic-spread-maximality}

\begin{definition}[Local and global maximality]
\label{def:appH-local-global-maximal}
A psh function \(u\) on a domain \(\Omega\) is maximal if every
psh \(v\geq u\) that equals \(u\) outside a compact subset of
\(\Omega\) equals \(u\) everywhere. It is locally maximal if
each point has a neighborhood on which this property holds.
Taking \(\max\{u,v\}\) gives the equivalent formulation with
an arbitrary competitor lying below \(u\) outside that compact set.
\end{definition}

Locally analytic singularities mean
\[
 u=a\log\!\left(\sum_{j=1}^r|f_j|^2\right)+O(1),
 \qquad a>0,
 \tag{H.2}\label{eq:appH-local-analytic-type}
\]
for generators of a nonzero coherent ideal \(I\), with a
two-sided locally bounded remainder. Neither that remainder's
continuity nor a global choice of generators is assumed. For
\(R=\mathcal O_{\Omega,z}\) with maximal ideal \(\mathfrak m_z\),
the analytic spread is
\[
 \ell(I_z)=\dim\!\left(
       \bigoplus_{k\geq0}I_z^k/\mathfrak m_z I_z^k\right).
 \tag{H.3}\label{eq:appH-analytic-spread}
\]
It measures the fibre dimension of the ideal blow-up, plus one;
it is invariant under powers and integral closure
\cite[\S\S2--4]{LejeuneTeissier08}.

\begin{theorem}[Do--Hiep: maximality for locally analytic singularities]
\label{thm:appH-analytic-local-global}
Let \(\Omega\subset\mathbf C^n\), \(n\geq2\), be a domain and
let \(u\in\PSH(\Omega)\) have locally analytic singularities.
With \(E=\{u=-\infty\}\), the following are equivalent:
\begin{enumerate}
\item \(u\) is maximal on \(\Omega\);
\item \(u\) is locally maximal on \(\Omega\);
\item the ordinary Bedford--Taylor measure
\((\ddc u)^n\) vanishes on \(\Omega\setminus E\), and
\(\ell(I_z)<n\) for every pole point and its local ideal in
\eqref{eq:appH-local-analytic-type}.
\end{enumerate}
\end{theorem}

\begin{proof}[From the analytic-spread and residual comparison inputs]
These are Theorem~1.1 and Corollary~1.4 of
\cite{DoHiep26MaximalityV1}.
Their proof has two distinct mechanisms. First, full analytic spread
forces a compact competitor near a pole. McAdam's integral-closure
theorem \cite{McAdam80AnalyticSpread} gives a germ
\(g\notin\overline{I^k}\) with
\(z_jg\in\overline{I^k}\) for every coordinate after translating
the pole to zero. The analytic integral-closure criterion then gives
\(|z||g|\leq A|I|^k\). An annular comparison with a suitably
translated multiple of \(\log|g|\) would force
\(g\in\overline{I^k}\) if \(u\) were locally maximal. Thus
local maximality implies the strict spread bound; off \(E\),
the usual locally bounded maximality criterion applies.

Conversely, the normalized ideal blow-up has fibre dimension
\(\ell(I_z)-1\). Under the strict bound, each relevant exceptional
prime has positive-dimensional image, which cannot be contained
in the competitor's compact support in a domain of \(\mathbf C^n\).
The local coefficients in \eqref{eq:appH-local-analytic-type}
can be matched on finitely many charts: vanishing orders make
overlapping pole coefficients rationally proportional, and powers
followed by integral closure glue the ideal data. Classical complex
Siu decomposition and normal-space psh extension then leave
residual potentials; the residual of \(u\) is locally bounded
and that of a compact competitor is locally bounded above.
Compactness of the blow-up fibres gives a local upper bound for
\(v-u\) on the finite locus.

On a relatively compact comparison domain, normalize both
potentials to be negative. The competitor \((1+\epsilon)v\)
lies below \(u\) near its poles because
\((1+\epsilon)v-u\leq\epsilon u+(1+\epsilon)C\).
Its positive comparison set is therefore away from the poles,
where zero Bedford--Taylor mass gives the usual comparison.
Let \(\epsilon\downarrow0\). This is the source's residual
comparison argument; McAdam's theorem, classical Siu theory and
normal-space extension are its external inputs.
\end{proof}

The blow-up in this argument is the normalized ideal blow-up.
A further resolution may introduce divisors over points; the spread
bound is not a bound for every exceptional divisor on every resolution.
For a principal nonunit ideal the spread is one. Consequently,
in dimension at least two, an analytic divisor potential is maximal
as soon as its measure vanishes off the divisor. An isolated pole
has a different spread obstruction.

\subsection{Compact comparison for local potentials}
\label{sec:appH-local-compact-comparison}

Let \(T=\ddc u_j\) and \(S=\ddc v_j\) on a cover of a complex
manifold \(M\). We require \(u_j,v_j\in\mathcal E_{\mathrm{loc}}\)
and common pluriharmonic transitions:
\[
 u_j-u_k=v_j-v_k\quad\text{on overlaps}.
 \tag{H.4}\label{eq:appH-common-local-transitions}
\]
This is stronger than having arbitrary local potentials for two
unrelated currents. Full local products give global Radon measures
\(T^n,S^n\), since the transitions do not change curvature.

\begin{proposition}[Do--Hiep: compact-support weighted comparison]
\label{prop:appH-local-compact-comparison}
Suppose \(u_j\leq v_j\) and \(u_j=v_j\) off a compact set
\(K\subset M\). Set \(h=v_j-u_j\) on the common finite locus,
and set it to zero on the pole sets in the integral below.
For bounded psh test functions \(-1\leq\varphi_j\leq0\)
on a neighborhood of \(K\),
\[
 \frac1{n!}\int_K h^n\,
       \ddc\varphi_1\wedge\cdots\wedge\ddc\varphi_n
 \leq\int_K(-\varphi_1)(T^n-S^n).
 \tag{H.5}\label{eq:appH-local-compact-weighted}
\]
The right side is a finite signed-measure integral on the compact
set; it is not a subtraction of possibly infinite global masses.
If \(M\) admits a global psh
\(\rho\in\mathcal E_{\mathrm{loc}}\) with
\((\ddc\rho)^n\geq f\,dV\), \(f>0\) almost everywhere,
then \(T^n\leq S^n\) implies \(u_j=v_j\) everywhere.
\end{proposition}

\begin{proof}[From local energy integration by parts]
The inequality is \cite[Proposition~3.3]{DoHiep26LocalComparisonV1},
in the book normalization. For bounded \(h\), truncation produces
the positive energy measure of its differential and local
integration by parts. A mixed trace with one smooth factor charges
no pluripolar set, by \cite[Lemma~4.4]{ACCH08PolarMeasuresV3}.
Common transitions make cutoff-gradient terms cancel on overlaps.
Iterating the weighted identity gives the factorial, and the mixed
curvature sum telescopes to \(T^n-S^n\). Replacing \(u_j\) by
\(\max\{u_j,v_j-t\}\) and using full local Cegrell convergence
gives the general compact-support case. The bounded test product
is nonpolar, so assigning zero on the pole sets does not change
its integral.

For domination, use a bounded local rescaling of
\(\exp\rho\) as every test function. Its Monge--Amp\`ere
measure has positive absolutely continuous density almost
everywhere, as established in the source's local truncation
argument. The assumed measure order makes the right side
nonpositive. Thus \(h=0\) almost everywhere, and equality of
the psh representatives follows. This is
\cite[Lemma~4.3]{DoHiep26LocalComparisonV1}; the zero-mass
case is its Corollary~4.4.
\end{proof}

Positive density almost everywhere is weaker than a global strictly
psh potential. For example, on the blow-up of a ball in \(\mathbf C^2\),
\(\rho=|\pi|^2\) has positive density off the exceptional curve,
which has volume zero; the compact curve prevents a global strictly
psh function. The density hypothesis cannot be dropped: on
\(\mathbf P^1\times\mathbf C\), \(|z|^2\) and
\(\max\{|z|^2,1\}\) differ on a compact set and have zero
top-degree measure. These are
\cite[Example~4.6 and the following discussion]{DoHiep26LocalComparisonV1}.

For noncompactly supported pairs, the same source requires boundary
auxiliaries increasing to zero almost everywhere and requires the
maxima formed with those auxiliaries to retain
\eqref{eq:appH-common-local-transitions}. Its weighted comparison
keeps the masses of \(\max\{u,v\}\) and \(u\) at common poles
on opposite sides. Removing them needs a separate polar comparison
input. The compact argument above supplies the surface application
below without importing the general mixed capacity theorem.

\subsection{The residual criterion on complex surfaces}
\label{sec:appH-surface-residual-maximality}

\begin{theorem}[Do--Hiep: full residual mass in dimension two]
\label{thm:appH-surface-residual-maximality}
Let \(\Omega\subset\mathbf C^2\) be a domain and
\(u\in\PSH(\Omega)\). Locally, suppose there are a negative
\(\psi\in\mathcal E\) and an analytic potential
\(\phi=a\log\sum|f_j|^2\), \(a>0\), or \(\phi=0\), such that
\[
 \phi+\psi\leq u,\qquad u+\psi\leq\phi.
 \tag{H.6}\label{eq:appH-surface-energy-sandwich}
\]
Write the classical complex Siu decomposition as
\[
 \ddc u=S_u+T_u,\qquad
 S_u=\sum_C\nu(\ddc u,C)[C].
\]
Then \(S_u\) is locally finite, \(T_u\geq0\) has local
\(\mathcal E\)-potentials,
and its full Cegrell square \(T_u^2\) is a globally defined
Radon measure. The following are equivalent:
\[
 u\text{ maximal}\quad\Longleftrightarrow\quad
 u\text{ locally maximal}\quad\Longleftrightarrow\quad T_u^2=0.
 \tag{H.7}\label{eq:appH-surface-residual-criterion}
\]
\end{theorem}

\begin{proof}[From surface factorization and compact comparison]
This is Theorem~1.5, with Lemma~5.1, of
\cite{DoHiep26LocalComparisonV1}.
Factor the local ideal generators as \(f_j=h_0g_j\).
The residual ideal in a smooth surface has isolated zeros, if any.
The potential \(a\log\sum|g_j|^2\) belongs locally to
\(\mathcal E\): the analytic lower bound at an isolated zero
and Błocki's surface criterion
\cite[Theorems~1.1 and~4.1]{Blocki04C2Published} supply this input.
That criterion also rules out a positive generic divisorial
Lelong number for \(\psi\); its transverse logarithmic gradient
would have infinite \(W^{1,2}\)-energy.
The sandwich therefore fixes
\(S_u=a[\operatorname{div}(h_0)]\) on this neighborhood.
Subtracting this current gives a positive residual current and
a psh representative bounded below by the residual analytic
potential plus \(\psi\), up to a constant. The local Cegrell
class is closed under such upper psh majorants, so this representative
belongs to \(\mathcal E_{\mathrm{loc}}\).
One subtracts currents and chooses their potentials; no value
\(-\infty-(-\infty)\) is taken pointwise. Pluriharmonic differences
make the local full squares agree.

For a compact competitor \(v\geq u\), every divisorial curve
escapes its compact support, and the generic Siu coefficients
agree outside that support. Thus \(S_v=S_u\). The residual
potentials \(g\geq f\) have common transitions, agree off the
compact set and lie in \(\mathcal E_{\mathrm{loc}}\).
If \(T_u^2=0\), apply
\Cref{prop:appH-local-compact-comparison} on a relatively compact
comparison domain, using \(\rho=|z|^2\), to obtain \(g=f\).
This proves global maximality. Global maximality implies local
maximality. Finally, the local maximality criterion and the
surface decomposition give \(T_u^2=0\), as proved in the
same source theorem. These steps use full local Cegrell mass,
including any mass carried by isolated poles.
\end{proof}

\begin{example}[An isolated pole and a divisor pole]
\label{ex:appH-full-residual-poles}
In \(\mathbf C^2\), let
\(u=\log(|z_1|^2+|z_2|^2)\). Its divisorial Siu part is zero,
but \((\ddc u)^2=\delta_0\) in the book normalization.
The nonpluripolar square is zero: off zero the curvature has
rank one, and that product omits the atom. The maximal ideal
has analytic spread two, so \(u\) is not maximal.
For \(u=\log|z_1|^2\), the divisorial part is
\([z_1=0]\), the residual is zero, and \(u\) is maximal.
This divisor potential is outside the surface local operator
domain; the criterion takes the full square of its residual
current rather than an unrestricted square of \(\ddc u\).
\end{example}

The surface assumption and
\eqref{eq:appH-surface-energy-sandwich} supply the residual
operator domain. The higher-dimensional Sobolev identification
fails, as Błocki's Kiselman example shows. These complex results
do not supply the general homogeneous non-Archimedean Siu
division input discussed in \Cref{app:provenance}; the BJ22
Siu/envelope-closure boundary remains separate.

\section{Smooth spectral operators and integral comparison}
\label{sec:appH-smooth-integral-comparison}
\index{Hessian operators!integral comparison}
\index{Newton tensor!Hessian divergence identity}
\index{plurisubharmonic function!on complex planes}

The smooth comparison used in a Hessian estimate and the construction
of a measure on singular potentials are different inputs. We record
the constant-coefficient local mechanism, following the sufficiency
argument in \cite[\S3]{DinewMygaSroka26IntegralV1}, and the published
obstruction in \cite[Theorem~4.1]{Dinew22SadullaevPublished}.
Their exact source scopes are in
\Cref{sec:provenance-Hessian-integral-comparison}.

\subsection{The coordinate operator and its comparison set}
\label{sec:appH-coordinate-comparison}

For a real potential on \(\Omega\subset\mathbf R^n\), write
\(H_u=D^2u\). For a real potential on
\(\Omega\subset\mathbf C^n\), write
\(H_u=(\partial_i\bar\partial_j u)_{i,j=1}^n\), using the
standard complex coordinates. In either case \(\sigma_k\) denotes
the unnormalized elementary symmetric polynomial of the eigenvalues,
and \(\sigma_0=1\). These matrices, rather than an endomorphism
defined by a variable Hermitian background, are the arguments below.

Let \(\mathcal G\) be an open conjugation-invariant domain in the
appropriate symmetric or Hermitian matrix space. A smooth potential
is admissible when \(H_u(x)\in\mathcal G\) everywhere. For a spectral
operator \(F\) on \(\mathcal G\), integral comparison means
\[
 \tag{H.8}\label{eq:appH-smooth-integral-comparison}
 \int_{\{u<v\}}F(H_v)\,dV
       \leq\int_{\{u<v\}}F(H_u)\,dV,
 \qquad \{u<v\}\Subset\Omega,
\]
where \(dV\) is ordinary, unweighted Lebesgue measure. This is an
inequality for every admissible pair, rather than comparison only
between solutions of one equation with the same right-hand side.
Addition of a constant to F preserves
\eqref{eq:appH-smooth-integral-comparison}: both sides gain the
same constant times the volume of the comparison set.

The complex coordinate convention is compatible with the book's
\(\ddc\), but is not a new convention for it. With
\(\beta=\ddc|z|^2\), diagonalization gives
\[
 \tag{H.9}\label{eq:appH-coordinate-Hessian-form}
 (\ddc u)^k\wedge\beta^{n-k}
   =\frac{\sigma_k(H_u)}{\binom nk}\,\beta^n,
 \qquad \beta^n=\frac{n!}{\pi^n}\,dV.
\]
Thus a coefficient in a polynomial in \(H_u\) is not the same
coefficient in a polynomial of wedge-product densities without this
binomial factor. Replacing \(\ddc\) by a scalar multiple also rescales
different degrees differently; a mixed-degree equation must retain
those factors.

\subsection{Newton tensors and the admissible segment}
\label{sec:appH-Newton-segment}

For a smooth real Hessian field, the Newton tensor satisfies
\[
 \tag{H.10}\label{eq:appH-Hessian-Newton-divergence}
 T_{k-1}^{ij}(D^2u)
      =\frac{\partial\sigma_k}{\partial A_{ij}}(D^2u),
 \qquad \sum_j\partial_jT_{k-1}^{ij}(D^2u)=0.
\]
This is the Euclidean Hessian instance of the Codazzi identity in
\cite[Lemma~2.1, course-note pp.~5--6]{GuanCIMECourseNotes}.
The symmetry of the third derivatives of u is essential.
It is false for an arbitrary symmetric matrix field: for
\(A(x)=\operatorname{diag}(1,x_1)\) in dimension two,
\(T_1(A)=\operatorname{diag}(x_1,1)\) has first-row divergence one.
The complex Hessian tensors have the analogous holomorphic and
antiholomorphic divergence identities. Equivalently, their
polarizations are built from the closed forms \(\ddc u\) and
the fixed flat form \(\beta\).

Here is the smooth-pair sufficiency argument of
\cite[(3.2) and Remark~3.2]{DinewMygaSroka26IntegralV1}, with its
Hessian hypothesis stated and the harmless constant included.

\begin{proposition}\label{prop:appH-polynomial-smooth-comparison}
Let
\(F(A)=c_0+\sum_{k=1}^n c_k\sigma_k(A)\).
Let u and v be smooth real admissible potentials and suppose
\(D=\{u<v\}\Subset\Omega\) has a smooth boundary on which
\(d(u-v)\ne0\). Assume every
\(tD^2u+(1-t)D^2v\), \(0\leq t\leq1\), belongs to
\(\mathcal G\), and the linearization of F is positive semidefinite
at each of these matrices. Then
\eqref{eq:appH-smooth-integral-comparison} holds.
\end{proposition}

\begin{proof}[From the Hessian Newton identity]
Put \(h=u-v\) and \(w_t=tu+(1-t)v\). The fundamental theorem
of calculus and \eqref{eq:appH-Hessian-Newton-divergence} give
\[
 \tag{H.11}\label{eq:appH-Newton-boundary-flux}
 \begin{aligned}
 \int_D\bigl(F(D^2u)-F(D^2v)\bigr)\,dV
  &=\int_0^1\int_D F^{ij}(D^2w_t)h_{ij}\,dV\,dt\\
  &=\int_0^1\int_{\partial D}
       F^{ij}(D^2w_t)\eta_i\eta_j\partial_\eta h\,dS\,dt
    \geq0.
 \end{aligned}
\]
Repeated indices are summed here. On the regular boundary of
\(\{h<0\}\), tangential derivatives vanish and the outward
normal derivative is positive. The matrix in the integrand is
positive semidefinite by assumption, giving the sign. The constant
coefficient has zero linearization and cancels from the difference.
\end{proof}

A convex matrix domain \(\mathcal G\) and ellipticity throughout it
provide the segment hypothesis. Neither membership of the two
endpoints alone nor connectedness of a nonconvex domain provides it.
The proof imposes no separate sign on each \(c_k\); the required
sign is that of the full linearization on the chosen segment.
The complex smooth-pair counterpart follows from the Hermitian
Newton identities, as in the cited Remark~3.3, or by applying Stokes
to the polarized forms in \eqref{eq:appH-coordinate-Hessian-form}.
No direct appeal to Sard's theorem for arbitrary C2 differences in
all dimensions is used in this regular-boundary statement.

For example, the complex dimension-four polynomial
\(F=\sigma_3-\sigma_1\) has positive linearization on the branch
\[
 \tag{H.12}\label{eq:appH-critical-polynomial-branch}
 \sum_{i\ne j}\lambda_i>
       \sqrt{2+\sum_{i\ne j}\lambda_i^2},
 \qquad j=1,\ldots,4.
\]
Each condition is a convex inequality; their intersection contains
\((1,1,1,1)\). This is the positive component in
\cite[Example~4.1]{DinewMygaSroka26IntegralV1}. It gives a selected
smooth elliptic branch with the segment property. Taking all
components of the ellipticity locus together would give a different
domain. The same source's real dimension-three Example~4.2 shows
failure of comparison between quadratics on different components.

\subsection{A product of eigenvalue sums has a different scope}
\label{sec:appH-product-eigenvalue-obstruction}

In \(\mathbf C^n\), a smooth function is p-plurisubharmonic when
its restriction to every affine complex p-plane is subharmonic.
Equivalently, the sum of its p smallest complex Hessian eigenvalues
is nonnegative. This notion differs from the usual m-subharmonic
conditions \(\sigma_j(H_u)\geq0\), \(1\leq j\leq m\), on the
closure of the appropriate G\"arding branch. In particular,
p=1 gives psh functions and p=n gives ordinary subharmonic functions.

For \(1\leq p\leq n\), define
\[
 \tag{H.13}\label{eq:appH-product-eigenvalue-sums}
 \mathcal M_p(H_u)=
       \prod_{\substack{I\subset\{1,\ldots,n\}\\|I|=p}}
                          \left(\sum_{i\in I}\lambda_i(H_u)\right).
\]
On the strict positive-sum branch this is a nonnegative elliptic
spectral operator. The branch is defined by positivity of each
factor, not merely positivity of their product. For instance,
\(\lambda_1\lambda_2>0\) alone also includes two negative eigenvalues.
The endpoints of \eqref{eq:appH-product-eigenvalue-sums} are
\(\mathcal M_1=\sigma_n\) and \(\mathcal M_n=\sigma_1\).
Already in dimension three,
\(\mathcal M_2=\sigma_1\sigma_2-\sigma_3\); this is nonlinear
in the elementary symmetric invariants.

The following is the published result of
\cite[Theorem~4.1, pp.~1004--1008]{Dinew22SadullaevPublished}.
Its index m is written p here to keep the complex-plane convention visible.

\begin{theorem}\label{thm:appH-Dinew-product-comparison}
Fix \(\alpha>0\). If \(\mathcal M_p^\alpha\) satisfies the
universal unweighted integral comparison principle for smooth
p-plurisubharmonic functions on bounded complex domains, then
\(\alpha=1\) and \(p=1\) or \(p=n\). Conversely, these endpoint
operators with \(\alpha=1\) satisfy the classical comparison principle.
\end{theorem}

\begin{proof}[From Dinew's two test families]
We recall the two necessary inequalities in the published proof.
First take a large positive diagonal quadratic and add a negative
smooth perturbation \(\chi\) whose value and gradient vanish at
the boundary of the unit ball. Sending n-p diagonal coefficients
to infinity reduces comparison to a power of the trace on the
remaining p coordinates. The linear term in a small multiple of
\(\chi\) integrates to zero. Its quadratic term has sign
\(\alpha(\alpha-1)\), so comparison requires \(\alpha\geq1\)
\cite[Claim~4.3, pp.~1005--1006]{Dinew22SadullaevPublished}.

For the other inequality, set
\(\chi_A(t)=\tfrac12((t+A)^2/(1+A)-1-A)\), \(A\geq0\),
and \(u_A(z)=\chi_A(|z|^2)\). These are psh, vanish on the
unit sphere, and decrease as A increases. Their Hessian eigenvalues
are \(\chi_A'\), repeated n-1 times, and
\(\chi_A'+|z|^2\chi_A''\). Computing the two integrals for
\(u_0\) and the limit of \(u_A\) gives
\[
 \tag{H.14}\label{eq:appH-Dinew-exponent-test}
 \left(1+\frac1p\right)^b\leq1+\frac bp,
 \qquad b=\binom{n-1}{p-1}\alpha.
\]
Strict Bernoulli inequality rules this out for b>1. Thus
\(\alpha\leq\binom{n-1}{p-1}^{-1}\), by
\cite[Claim~4.4, pp.~1006--1007]{Dinew22SadullaevPublished}.
Together with \(\alpha\geq1\), this leaves only p=1 or n
and \(\alpha=1\). The endpoints are the ordinary complex
Monge--Amp\`ere and Laplace comparison principles.
\end{proof}

The statement concerns all admissible smooth pairs. A radial-only
comparison can have the exponent
\(\binom{n-1}{p-1}^{-1}\), as discussed in the same source's
Remark~4.5, without contradicting it. Nor does failure of this
integral principle exclude viscosity comparison for a fixed equation.

\subsection{What comparison leaves to the weak theory}
\label{sec:appH-smooth-comparison-weak-boundary}

For a general smooth F, compact nonnegative perturbations of an
admissible u give a useful necessary test for universal comparison.
Differentiating \(v_t=u+t\phi\), \(\phi\geq0\) compactly supported,
and integrating by parts yields
\[
 \tag{H.15}\label{eq:appH-comparison-adjoint-test}
 \sum_{i,j}\partial_i\partial_jF^{ij}(D^2u)\leq0
\]
in the real setting, in the distributional sense; for smooth
coefficients it is a pointwise inequality. This is the first-variation
test in \cite[(3.3)--(3.6)]{DinewMygaSroka26IntegralV1}.
It is an obstruction test, not a construction of a weak nonlinear
measure. The remaining invariant-coordinate jet expansion requires
more regularity than the merely differentiable F declared in that
edition; the exact audit in
\Cref{sec:provenance-Hessian-integral-comparison} records this boundary.

The literal zero-constant classification formula in that preprint
also needs a normalization: comparison cannot detect an additive
constant. Its homogeneous branch equality needs a domain selection
condition, since comparison survives restriction to a proper
admissible subcone. We use its sufficiency mechanism with the
displayed hypotheses, rather than importing those unrestricted
classification statements.

Similarly, the fixed equation
\(\sigma_n/\sigma_{n-k}=c\), where the denominator is positive,
can be rewritten as \(\sigma_n-c\sigma_{n-k}=0\).
The polynomial and the raw quotient are different operators in
\eqref{eq:appH-smooth-integral-comparison}. A branch comparison
for one does not automatically supply that integral principle for
the other. For singular potentials one still needs a domain of
definition, approximation independence, local mass control and
the relevant continuity and comparison theorems. The variable-background
Hermitian setting of \Cref{sec:ytd-direct-gradient-compactness}
requires its own analytic inputs. None of these local spectral
observations supplies general non-Archimedean Siu division or
envelope closure.

\section{Unbounded boundary obstacles and hypersurface Green functions}
\label{sec:appH-boundary-envelopes}
\index{Perron--Bremermann envelope!unbounded obstacle}
\index{pluricomplex Green function!hypersurface poles}

A boundary trace, an interior obstacle and a prescribed singularity impose
different conditions on a potential. The two October 2026 preprints
\cite{DaoNguyen26UnboundedDirichletV1,DaoNguyenTran26HypersurfaceGreenV1}
make these conditions explicit for local complex domains. We explain their
attributed conclusions with the deformation hypotheses retained. The exact
editions, the older erratum and the limits of the source audit are recorded
in \Cref{sec:provenance-boundary-envelopes}.

Throughout this section \(\Omega\Subset\mathbf C^n\) is a bounded domain.
Local maximality for a locally bounded psh function means vanishing of its
ordinary Bedford--Taylor top measure. For continuous finite boundary data
\(b\), write \(\operatorname{PB}_{b,\Omega}\) for the supremum of psh
functions whose upper boundary limits are at most b. On a B-regular domain
this is the continuous maximal solution with boundary value b.

\subsection{The obstacle fixes the admissible class}
\label{sec:appH-boundary-obstacle-class}

Let \(\phi:\overline\Omega\to(-\infty,+\infty]\) be extended-real
continuous, ordinary real superharmonic in \(\Omega\), and not identically
\(+\infty\) there. Compactness gives a finite lower bound. Define
\[
 G_{\phi,\Omega}(z)
   =\sup\{v(z):v\in\PSH(\Omega),\ v\leq\phi\text{ on }\Omega\},
 \qquad B=\{\phi=+\infty\}\cap\partial\Omega.
 \tag{H.16}\label{eq:appH-boundary-obstacle}
\]
Theorem~1.1 of \cite[actual pp.~14--15, proof]{DaoNguyen26UnboundedDirichletV1}
shows that this supremum is already psh, locally bounded and maximal,
and is bounded from below. At a boundary point with a continuous psh
barrier it tends to the extended value \(\phi(\xi)\).
The submean inequality and the local integrability of the superharmonic
obstacle give uniform local upper bounds for all minorants. Upper
regularization remains below \(\phi\); a local maximum-pasting argument
then proves maximality. Thus superharmonicity is an operative hypothesis.
For the disk obstacle \(\phi(z)=|z|^2\), the interior psh envelope is
\(|z|^2\), whereas the Perron envelope of its constant boundary value is
one. An arbitrary continuous interior obstacle cannot replace the one in
\eqref{eq:appH-boundary-obstacle}.

Uniqueness also requires a class of competitors. Under the deformation
conditions below, suppose B is contained in a complete pluripolar set on
a connected neighborhood of \(\overline\Omega\), and every boundary
point outside B has a continuous psh barrier. Then
\cite[Theorem~1.9; actual pp.~24--25]{DaoNguyen26UnboundedDirichletV1}
identifies \(G_{\phi,\Omega}\) as the unique lower-bounded, locally
bounded maximal psh function u with \(u\leq\phi\) and boundary trace
\(\phi\) outside B. The same obstacle occurs in both the existence and
uniqueness assertions.

The role of the function class is visible in the comparison lemma
\cite[Lemma~2.6(i), actual pp.~8--9]{DieuLongHieu25QuasiBRegularV1}.
In its locally bounded maximal comparison case, it compares a
quasi upper bounded psh u with a maximal psh v whose lower boundary
limits are finite from below. The upper boundary limits of u must lie
below those lower limits off an \(\Omega\)-pluripolar set. Here quasi
upper bounded means admitting an increasing sequence of upper-bounded
psh minorants converging outside an interior pluripolar set. The
boundary exceptional set has a negative psh witness on \(\Omega\)
with upper boundary limit \(-\infty\). For an upper-bounded minorant,
adding a small positive multiple of that witness restores comparison
on a compact interior domain; one then takes that multiple to zero
and passes to the increasing sequence. A neighborhood pluripolar hull
uses psh functions on a neighborhood of the closure and has a different
definition.

For a concrete boundary test, take the unit disk and
\(\phi(z)=\log(2/|1-z|)\). For any \(c>0\),
\[
 u_c(z)=\phi(z)+c\operatorname{Re}\frac{1+z}{1-z}
\]
is lower-bounded harmonic and has the same extended boundary trace as
\(\phi\): the added Poisson kernel tends to zero away from one, while
\(\phi\) tends to \(+\infty\) at one. At zero, \(u_c(0)=\phi(0)+c\),
so it leaves the class \(u\leq\phi\). This is the fixed-majorant
obstruction explained in the same source, rather than a failure of its
uniqueness theorem. The higher-dimensional comparable-growth phenomenon
in \Cref{sec:ytd-boundary-poisson-nonuniqueness} supplies another reason
to state the admissible class before using boundary comparison.

\subsection{A deformation with one-sided obstacle control}
\label{sec:appH-boundary-neighborhood-deformation}

The Dirichlet continuity theorem uses biholomorphic maps \(\Phi_t\)
on one fixed neighborhood of \(\overline\Omega\), continuous in
\((t,z)\), with \(\Phi_0=\mathrm{id}\) and \(\Phi_t\to\mathrm{id}\)
locally uniformly. For each \(z\in\Omega\), the inverse trajectory is
defined for small t and is real analytic near zero. Its boundary cluster is
\[
 S_\Phi=\{\lim_j y_j:t_j\downarrow0,\quad
       y_j\in\overline\Omega\cap\Phi_{t_j}(\partial\Omega)\}.
\]
The closed domain in this definition matters: the identity deformation
has \(S_\Phi=\partial\Omega\). The required one-sided condition is
\[
 \begin{gathered}
 \limsup_j\bigl(\phi(x_j)-\phi(\Phi_{t_j}(x_j))\bigr)\leq0
 \quad\text{if}\\
 x_j\in\overline\Omega,\quad x_j\to\xi\in B,\\
 \Phi_{t_j}(x_j)\in\overline\Omega\setminus\{\phi=+\infty\}.
 \end{gathered}
 \tag{H.17}\label{eq:appH-boundary-one-sided}
\]
In addition, \(S_\Phi\) must avoid \(\{\phi=+\infty\}\), and every
point of \(S_\Phi\) must have a continuous psh barrier. These are
Definitions~1.2--1.3 and conditions D1--D2 of
\cite[actual pp.~4--5]{DaoNguyen26UnboundedDirichletV1}.
There is no evaluation of the obstacle outside \(\overline\Omega\)
and no difference of two infinite values in
\eqref{eq:appH-boundary-one-sided}.

Under these conditions, Theorem~1.4 expresses \(G_{\phi,\Omega}\)
as the supremum of finite psh minorants continuous on the closure,
and proves its interior continuity. Neighborhood barriers at
\(S_\Phi\) give minorants continuous and psh on neighborhoods of the
closure. The proof, actual~15--18, transports the old envelope by
\(\Phi_t^{-1}\), controls its obstacle error, regularizes on a compact
cut away from the moving boundary and uses strict barrier margins to
glue by a maximum. The domains \(\Phi_t(\Omega)\) need not be nested.

If finite continuous \(b_j\) increase to \(\phi\) on the whole
boundary, the same theorem gives
\(\operatorname{PB}_{b_j,\Omega}\uparrow G_{\phi,\Omega}\)
locally uniformly. Pointwise convergence of a nonmonotone boundary
sequence has less force. Lemma~3.3 and Corollary~3.4 give convergence
away from the indicated neighborhood pluripolar hull; Proposition~3.6
requires uniform finite-truncation control, or its equivalent eventual
domination condition, for convergence throughout \(\Omega\).
Example~5.2, actual~28--29, shows why the stronger hypothesis is needed.

\subsection{A scalar divisor has a psh residual}
\label{sec:appH-boundary-divisor-residual}

Suppose \(f\not\equiv0\) is holomorphic and reduced on a neighborhood
of \(\overline\Omega\), and let A be its divisor in \(\Omega\).
After rescaling f assume \(\|f\|_{\overline\Omega}\leq1\), so
\(\phi=-\log|f|\geq0\). The unsquared logarithm follows the two
sources. In the book normalization,
\(\ddc\log|f|=\tfrac12[A]\). The local divisor Green class can
therefore be written as negative psh u with
\(\ddc u\geq\tfrac12[A]\). This is the reduced-divisor Lelong
constraint expressed as a current inequality; a source Lelong
normalization using \(\log|z|\) has twice the book's coefficient.

Classical divisorial subtraction gives a psh extension of
\(G_{A,\Omega}-\log|f|\), locally bounded above across A. The
competitor \(\log|f|\) bounds that residual from below by zero.
Adding \(\log|f|\) to an obstacle minorant and subtracting it from
the Green envelope yield
\[
 \widetilde G_A=G_{-\log|f|,\Omega},\qquad
 G_{A,\Omega}=\log|f|+\widetilde G_A\quad\text{on }\Omega\setminus A.
 \tag{H.18}\label{eq:appH-boundary-Green-residual}
\]
The residual is extended as a psh function; no value
\(-\infty-(-\infty)\) is assigned on A. The classical input is
\cite[Proposition~3.2, printed~1525--1526]{LarussonSigurdsson99GreenPublished};
the obstacle identity is used in
\cite[(1.4), actual~3]{DaoNguyen26UnboundedDirichletV1}.

This scalar identity does not identify every Green function associated
with a coherent ideal. Preserving a zero set is also weaker than
preserving its ideal. For \(F=(z_1^2,z_2)\), the map
\((z_1,z_2)\mapsto(z_1,z_2+t z_1)\) fixes the zero set, but its
second pullback generator is outside \((z_1^2,z_2)\) for \(t\ne0\).
The tuple-obstacle extension in
\cite[Proposition~6.1, actual~31]{DaoNguyen26UnboundedDirichletV1}
instead assumes a holomorphic identity
\(F\circ\Phi_t=M_tF\), with \(M_t\to I\) uniformly near the
closure. It supplies obstacle control with that ideal hypothesis.

\subsection{Canonical truncation needs an inward sector}
\label{sec:appH-boundary-SID}

For \(n\geq2\) and \(A\ne\varnothing\), the Green-function paper
assumes a sectorial inward deformation. This consists of a jointly
holomorphic \(S(z,\tau)\) near \(\overline\Omega\) and a full
parameter disk, with \(S_0=\mathrm{id}\). For
\(0<|\tau|<t_0\), \(|\arg\tau|<\theta<\pi/2\), it requires
\[
 \begin{gathered}
 S_\tau(\overline\Omega)\subset\Omega,\qquad
 \sup_{\overline\Omega}|S_\tau(z)-z|\leq C|\tau|,\\
 \operatorname{dist}(S_\tau(z),\partial\Omega)\geq c|\tau|
 \quad\text{on }V\cap\overline\Omega,
 \end{gathered}
 \tag{H.19}\label{eq:appH-boundary-SID}
\]
where V is a neighborhood of \(\{f=0\}\cap\partial\Omega\).
At each singular boundary divisor point p, the germ
\(\tau\mapsto f(S_\tau(p))\) must not vanish identically.
This finite-order escape condition is imposed only at singular
boundary points. No center is required to be fixed.
These are exactly the features of Definition~1.4 in
\cite[actual pp.~4--5]{DaoNguyenTran26HypersurfaceGreenV1}.

With \(P_\Omega\) denoting the interior psh envelope, put
\(v_j=P_\Omega\min\{-\log|f|,j\}\). Theorem~1.7 gives
\[
 v_j\uparrow\widetilde G_A\quad\text{on }\Omega\setminus A,
 \qquad\left(\sup_jv_j\right)^*=\widetilde G_A
                              \quad\text{on }\Omega.
 \tag{H.20}\label{eq:appH-boundary-canonical-truncation}
\]
This approximation assertion does not require B-regularity.
With B-regularity, the bounded canonical envelopes are continuous
Perron--Bremermann solutions. Their increasing limit is continuous
off A, where it also equals the psh residual. The local upper bound
for the residual then gives \(G_A\to-\infty\) at A. Thus the
conclusion is extended-real continuity of \(G_A\) on all
\(\Omega\). It does not assert continuity of \(\widetilde G_A\)
at every point of A.

The missing inequality in a much earlier attempt is explicit.
\cite[Theorem~3.9, printed~1532--1533]{LarussonSigurdsson99GreenPublished}
claimed a general almost-continuity statement using bounded boundary
truncations. The controlling \cite[printed~1705]{LarussonSigurdsson01GreenErratum}
says that \(\widetilde G_A\leq v^*\) cannot be concluded and that
the theorem remains unproved. Its alternative candidate restricts
\(u-\log|f|\) to be bounded above on the whole domain, rather than
only locally near the divisor. The erratum gives no counterexample
and leaves the other results unaffected. The SID theorem proves the
needed identification under its added geometric hypotheses.

\subsection{The disc argument recovers logarithmic intersection weight}
\label{sec:appH-boundary-disc-weight}

For an analytic disc h holomorphic near the closed disk, with finite
center away from A, the divisor Riesz functional has the form
\[
 H_A(h)=\sum_{a\in\mathbb D}
           \operatorname{ord}_a(f\circ h)\log|a|.
 \tag{H.21}\label{eq:appH-boundary-divisor-disc}
\]
The weights are nonpositive and use intersection orders of the scalar
defining function. The divisor envelope formula in
\cite[Theorem~3.3, printed~1526--1527]{LarussonSigurdsson99GreenPublished}
retains a relatively compact Stein ambient space and extension of
the subspace to it. The general manifold and locally irreducible
space functional results of
\cite[Theorem~1.1 and functional definitions, printed~172--174]{DrinovecForstneric12DiscPublished}
provide a separate interface. A pointwise Lelong-weight functional
is not literally \eqref{eq:appH-boundary-divisor-disc} for an
arbitrary ideal.

The proof of the new approximation theorem, actual~29--33, starts
with a near-extremal disc for the bounded envelope. It replaces
the lost logarithmic contribution near its boundary divisor trace
by finitely many controlled intersections. A finite Weierstrass
chart cover separates interior, regular boundary and singular
boundary regions. In a singular chart, \(\tau\) is the Weierstrass
variable; the roots need not be holomorphic in the base. The
relative estimate of Lemma~5.4(ii) has a constant bounded by
\(\log(u_+/u_-)+m_*\log(1+2/\kappa)\), independent of the
truncation level. Here \(u_-,u_+\) bound the units, \(m_*\) bounds
the finite chart degrees, and \(\kappa\) controls the insertion
circle relative to the inward depth.

The center moves under \(S_\tau\). Lemma~5.5 uses the submean
inequality for the psh pullback along that parameter curve to
choose a nearby center whose residual loses only \(o(1)\).
It does not assume the continuity being proved. For truncation j,
the proof takes a scale proportional to \(\exp(-j/m_*)\) and a
cutoff size \(j^{-1/2}\). Bad arcs have length \(O(1/j)\);
singular relative insertions lose a bounded amount, while the
regular contributions have total error \(O(\log j/j)\).

The root-cell and gluing mechanism follows
\cite[Lemma~3.2 and Theorem~3.1, printed~179--184]{DrinovecForstneric12DiscPublished}.
The target contours, Rouch\'e margins and source arcs are fixed
before the Cartan pair and cutoff, and before the large integer
distributing the roots. The final radial shrink of the glued disc
has its logarithmic cost included. These choices account for the
uniformity needed in \eqref{eq:appH-boundary-canonical-truncation}.
This description records the source mechanism; it does not certify
every external disc-envelope and gluing theorem.

\subsection{Geometry, complementary results and remaining questions}
\label{sec:appH-boundary-Green-directions}

A convexifying biholomorphism must extend to neighborhoods of both
closures. Contraction of the convex domain toward a point outside A
then gives the depth and escape conditions; this is Proposition~2.1
of \cite[actual~11]{DaoNguyenTran26HypersurfaceGreenV1}.
The same paper treats nonconvexifiable and Reinhardt examples with
their own geometric and B-regularity inputs. Absence of boundary
analytic discs is used with the specified convex or Reinhardt geometry;
it is not imported as a sufficient criterion on every bounded domain.

The published homogeneous result
\cite[Theorem~2.2 and Corollary~2.9, printed~184--185 and~192]{Nguyen05HypersurfaceGreenPublished}
treats unions of hyperplanes through a common interior point, with
holomorphic convexity of the closure, radial inwardness near the
boundary trace and good barriers elsewhere. It gives continuity
of the residual on all \(\Omega\) and includes strongly star-shaped
pseudoconvex domains without assuming B-regularity. Radial motion
through the common singular pole preserves the homogeneous divisor
and fails the SID escape test at a singular boundary point.
The older radial result and the new inward-sector result therefore
retain different hypotheses and conclusions.

The explicit questions in
\cite[actual pp.~33--34]{DaoNguyenTran26HypersurfaceGreenV1}
leave four related directions:
\begin{enumerate}
\item For an arbitrary bounded B-regular domain and a divisor reduced
      near its closure, is the divisor Green function extended-real
      continuous throughout the domain? The new proof handles pairs
      admitting the specified inward sector.
\item On an arbitrary bounded \(C^2\) strongly pseudoconvex domain,
      can continuity and the canonical truncation identity
      \eqref{eq:appH-boundary-canonical-truncation} be proved without
      convexifiability or SID? The question persists even for smooth
      strongly pseudoconvex boundaries in the stated generality.
\item Can higher-codimensional coherent ideals admit a stable disc
      insertion and gluing theory carrying the correct ideal weights?
      The scalar Weierstrass root argument and the matrix obstacle
      identity above do not supply that theory.
\item Which further geometric classes admit SID, including uniform
      singular-boundary escape? Real inward motion alone does not
      supply the complex-sector estimates.
\end{enumerate}
These are complex literature questions accompanying the local reference.
They do not change the canonical problem inventory of
\Cref{chap:major-open-problems}. Their comparison classes and scalar divisor
subtraction also provide no general non-Archimedean Siu division or
envelope-closure input. The BJ22 Theorem~6.21 proof boundary remains
as recorded in \Cref{app:provenance}.

\section{Comparison on unbounded domains}
\label{sec:appH-unbounded-domain-comparison}
\index{Perron--Bremermann envelope!unbounded domains}
\index{q-plurisubharmonic function!Hunt--Murray convention}

The domain can instead be unbounded while its boundary datum is bounded.
This changes the comparison problem: finite boundary values give no
automatic control along sequences escaping to infinity. We now work on
an unbounded domain \(D\subset\mathbf C^n\), separately from the bounded
\(\Omega\) of \Cref{sec:appH-boundary-envelopes}.
The recent account \cite{Pawlaschyk26UnboundedPerronV1} gives bounded
Perron solutions with finite-boundary peaks and a negative psh witness
at infinity. We retain its capped conclusions and distinguish them
from a literal uncapped uniqueness claim. Exact versions and proof
qualifications are recorded in
\Cref{sec:provenance-unbounded-domain-comparison}.

\subsection{The eigenvalue comparison branch}
\label{sec:appH-unbounded-q-branch}

For \(0\leq q\leq n-1\), write
\(\operatorname{PSH}^{\mathrm{HM}}_q(D)\) for the Hunt--Murray
q-plurisubharmonic class used in these sources. It consists of usc
functions satisfying the pluriharmonic comparison principle on complex
\((q+1)\)-dimensional affine planes. For a \(C^2\) function, the
complex Hessian has at most q negative eigenvalues
\cite[Theorem~2.8; actual p.~3]{Pawlaschyk26UnboundedPerronV1}.
At q=0 this is ordinary psh. It differs from nonnegative sums of p
eigenvalues, and therefore from the p-plane condition in
\Cref{sec:appH-smooth-integral-comparison}.

A continuous function F is q-Bremermann when F belongs to this q-class
and \(-F\) belongs to its \((n-q-1)\)-class. In the smooth case, if
\(\lambda_1\leq\cdots\leq\lambda_n\) are its Hessian eigenvalues,
the two conditions give \(\lambda_{q+1}=0\). Consequently the determinant
vanishes. This is the meaning of the smooth Monge--Amp\`ere observation
in \cite[Definition~5.2; printed p.~320]{Slodkowski84BremermannPublished}.
It does not by itself define a Bedford--Taylor measure for a general
continuous, non-psh q-Bremermann function. In \(\mathbf C^3\), for example,
\(-2|z_1|^2+|z_3|^2\) has Hessian \(\operatorname{diag}(-2,0,1)\)
and is smooth 1-Bremermann, but has negative trace and is not subharmonic.
For q=0, a locally bounded psh maximal function does have the ordinary
Bedford--Taylor equation \((\ddc F)^n=0\).

The addition theorem is
\cite[Theorem~5.1; printed pp.~319--320]{Slodkowski84BremermannPublished}:
the sum of a q-psh function and an r-psh function lies in the
\((q+r)\)-class. A fixed q-class need not be closed under addition.
For instance, the diagonal Hessians \((-1,0,1)\) and \((0,-1,1)\)
each have one negative eigenvalue, but their sum has two. Adding an
ordinary psh function preserves q; this is the operation used below.
The general addition theorem uses semiconvex approximation, not an
unrestricted claim that convolution preserves q.

\subsection{A finite cap is part of the envelope}
\label{sec:appH-unbounded-capped-envelope}

Let \(f\in C(\partial D)\) be bounded and real valued, and fix a finite
\(M\geq\sup_{\partial D}f\). The continuous-competitor envelope is
\[
\begin{split}
 P_{f,q,D,M}(z)=\sup\{v(z):\;&v\in
      \operatorname{PSH}^{\mathrm{HM}}_q(D)\cap C(\overline D),\\
     &v\leq f\text{ on }\partial D,\quad v\leq M\text{ on }D\}.
\end{split}
\tag{H.22}\label{eq:appH-unbounded-capped-envelope}
\]
The boundary values concern the ordinary finite boundary of D.
The inequality \(v\leq M\) constrains all competitors in the interior
and along escape directions. On a bounded domain it follows from the
maximum principle; on an unbounded domain it requires justification.
The envelope of closure-usc competitors with the same cap is a
separate class until its equality with
\eqref{eq:appH-unbounded-capped-envelope} is proved.

The source's bounded-type condition requires
\[
\begin{gathered}
 \Psi\in\PSH(D)\cap C(D),\qquad \Psi<0,\\
 \text{for every }C\geq0\text{ there is }R>0\text{ such that }
      \Psi<-C\text{ on }D\setminus B_R(0).
\end{gathered}
\tag{H.23}\label{eq:appH-unbounded-infinity-witness}
\]
Continuity here is on the open domain; a continuous extension of \(\Psi\)
to \(\overline D\) is not assumed. The older
\cite[Definition~2.6; printed p.~1389]{NguyenHung08UnboundedPublished}
requires a finite negative psh witness without imposing its continuity.
Finite-boundary regularity alone does not supply either infinity condition.
The polynomial examples in these sources require positive degree:
a constant polynomial cannot provide the proposed growing witness.

An r-peak at \(p\in\partial D\) is an r-psh function continuous on
\(\overline D\), zero at p and negative elsewhere on the closure.
Assume every finite boundary point has such a peak and retain
\(0\leq r\leq q\leq n-r-1\). Under
\eqref{eq:appH-unbounded-infinity-witness},
\cite[Theorem~4.19(1)--(4) and Remark~4.20; actual pp.~10--12]{Pawlaschyk26UnboundedPerronV1}
gives a bounded continuous q-Bremermann solution with boundary value f,
equal to both capped envelopes. Its statement uses \(f\geq0\);
an additive shift treats a general bounded f. All finite caps at least
\(\sup_{\partial D}f\) give the same envelope. This is an attributed
bounded existence conclusion, with the classical q-Dirichlet and
approximation inputs retained.

\subsection{Comparison at infinity uses the growth class}
\label{sec:appH-unbounded-infinity-comparison}

Proposition~4.17 of \cite[actual pp.~9--10]{Pawlaschyk26UnboundedPerronV1}
applies the maximum principle to a bounded-above q-psh competitor \(\varphi\).
For \(\varepsilon>0\), \(\varphi+\varepsilon\Psi\) stays in the same q-class
and lies below its finite boundary bound outside a sufficiently large ball.
Its upper limits at the finite boundary obey the same bound because
\(\Psi<0\). The bounded-domain maximum principle and the limit
\(\varepsilon\downarrow0\) then give comparison on D. A continuous
boundary extension of \(\Psi\) is unnecessary for this upper-limit argument.

If F and G are bounded q-Bremermann functions with the same continuous
boundary value, the addition theorem makes \(F-G\) and \(G-F\)
\((n-1)\)-psh. Both are bounded above. Applying that infinity principle
in both directions gives F=G. This bounded uniqueness is a direct
consequence of the cited comparison inputs. The older
\cite[Theorem~4.4; printed pp.~1398--1400]{NguyenHung08UnboundedPublished}
also explicitly restricts uniqueness to bounded maximal psh solutions.
Its Jensen proof has separate source qualifications; it is not a proof
of uniqueness for arbitrary growth.

Local maximality only compares on bounded relatively compact sets.
It does not supply the outer-ball inequality in this argument. Similarly,
taking a supremum over all finite caps does not identify the unrestricted
\(M=+\infty\) envelope.

\subsection{An exterior-disk check of uncapped uniqueness}
\label{sec:appH-unbounded-growth-counterexample}

The following elementary check keeps the missing condition visible:
\[
\begin{aligned}
 D&=\{z\in\mathbf C:|z|>1\},&\qquad f&=0,\\
 \Psi(z)&=-1-\log|z|,& h_p(z)&=\operatorname{Re}(p/z)-1,
                  \quad |p|=1,\\
 F_c(z)&=c\log|z|,& c&\in\mathbf R.
\end{aligned}
\tag{H.24}\label{eq:appH-unbounded-exterior-family}
\]
The witness is negative harmonic and tends uniformly to \(-\infty\).
Each \(h_p\) is a closure-continuous harmonic peak: inside D,
\(\operatorname{Re}(p/z)\leq|z|^{-1}<1\), and equality on the unit
circle occurs only at p. Every \(F_c\) is harmonic on D, continuous
on its closure and zero on its finite boundary. Thus they are distinct
0-Bremermann solutions with the same trace.

The bounded solution is zero. For \(c=1\) the competing solution is
unbounded above, so the bounded-above maximum principle cannot compare
it to zero. For \(c=-1\) it is bounded above but unbounded below;
for \(0<\varepsilon<1\),
\[
  0+\varepsilon\Psi-F_{-1}
     =(1-\varepsilon)\log|z|-\varepsilon\longrightarrow+\infty.
\]
The outer comparison also fails in this direction. Consequently the
literal uniqueness among all q-Bremermann functions in
\cite[Theorem~4.19(5), Claim~5; actual pp.~10, 12]{Pawlaschyk26UnboundedPerronV1}
is false as written. This diagnoses its omitted growth condition; it
does not refute the bounded existence conclusion used above.

This finite-boundary example complements the singular boundary-trace
example in \Cref{sec:appH-boundary-obstacle-class}. In both cases a
uniqueness statement requires an admissible class. Selecting growth at
infinity for an uncapped envelope, specifying an appropriate boundary hull,
and defining weak equations for general q-branches require additional
inputs. These local complex questions do not change the canonical
inventory in \Cref{chap:major-open-problems} or repair general
non-Archimedean Siu division and envelope closure.


\section{Griffiths positivity and higher Chern forms}
\label{sec:appH-Griffiths-Chern}
\index{Griffiths positivity!top Chern form}
\index{Chern form!pointwise and numerical positivity}

The positive products used earlier in this appendix come from scalar
closed positive $(1,1)$-currents. A higher Chern form of a Hermitian
vector bundle has a different algebraic construction. Even for a
smooth Griffiths-positive metric, its top Chern form can have negative
values. This section explains the attributed construction of
\cite[Theorem~1.1]{Du26GriffithsTopChernV1} and the distinction it makes
between curvature, a characteristic representative and its integral.
All bundles and differential forms in this section are complex and
smooth.

\subsection{A representative and its Chern number}
\label{sec:appH-Chern-representative}

Let $(E,h)$ be a holomorphic vector bundle of rank $r$ on a complex
manifold. Write $R^E=\bar\partial(h^{-1}\partial h)$ for its Chern
curvature in a compatible local matrix convention, and define
\begin{equation*}
 \Omega_h=\frac{\rmi}{2\pi}R^E,\qquad
 \det(I+s\Omega_h)=\sum_{q=0}^r s^q c_q(E,h).
 \tag{H.25}\label{eq:appH-Chern-normalization}
\end{equation*}
For a line bundle, $h(e,e)$ is the squared norm in a holomorphic
frame and $c_1(E,h)=-\ddc\log h(e,e)$. Thus
\eqref{eq:appH-Chern-normalization} agrees with the scalar convention
at the start of this appendix.

Griffiths positivity means
$\langle R^E(u,\bar u)v,v\rangle_h>0$ for every nonzero tangent
vector $u$ and fibre vector $v$. It implies positivity of the trace
form $c_1(E,h)$. The determinant defining $c_r(E,h)$ contains
alternating terms; its pointwise sign requires a separate argument.

The cited construction gives, for every $r\geq9$, a smooth
Griffiths-positive metric on
\[
 E=\mathcal O_{\mathbf P^r}(1)^{\oplus r}
\]
whose top Chern form is negative on a nonempty open subset of
$\mathbf P^r$. Nevertheless, if
$H=c_1(\mathcal O_{\mathbf P^r}(1))$, then
\begin{equation*}
 c(E)=(1+H)^r,\qquad c_r(E)=H^r,\qquad
 \int_{\mathbf P^r}c_r(E,h)=1.
 \tag{H.26}\label{eq:appH-Chern-number}
\end{equation*}
The positive integral and the locally negative representative are
consistent. For example, the direct-sum Fubini--Study metric on the
same bundle has positive top Chern form. The example concerns the
representative determined by a specified metric. It does not change
the Chern class or contradict numerical positivity of ample-bundle
Schur classes.

\subsection{The positive-map calculation}
\label{sec:appH-Chern-positive-map}

For a complex-linear adjoint-preserving map
$\Phi:\operatorname{Mat}_n(\mathbf C)\to
\operatorname{Mat}_n(\mathbf C)$, put
\begin{equation*}
 F_n(\Phi)=\det(\partial_Y)\det\Phi(Y),\qquad
 \partial_Y=\left(\frac{\partial}{\partial y_{ij}}\right)_{i,j=1}^n.
 \tag{H.27}\label{eq:appH-Chern-det-pairing}
\end{equation*}
The $n^2$ entries are independent complex polynomial variables.
If $\Phi$ sends each nonzero positive semidefinite matrix to a
positive definite matrix, the local metric
\[
 h_\Phi(z)=I-2\pi\Phi(zz^*)
\]
is positive definite on a small ball and Griffiths positive there.
Its first jets vanish at zero and its curvature coefficients at zero
are $2\pi\Phi(E_{ij})$. Consequently
\begin{equation*}
 c_n(E,h_\Phi)|_0=F_n(\Phi)\,dV_n,\qquad
 dV_n=\bigwedge_{j=1}^n(\rmi\,dz_j\wedge d\bar z_j).
 \tag{H.28}\label{eq:appH-Chern-local-volume}
\end{equation*}
Each factor is $2\,dx_j\wedge dy_j$, fixing the orientation and
normalization. Rank-one input tests and spectral decomposition give
the positive-map/Griffiths dictionary; see
\cite[Propositions~2.10--2.11]{Finski22CharacteristicPublished}.
The negative quadratic jet above fixes the curvature sign explicitly.

The pairing in \eqref{eq:appH-Chern-det-pairing} is bilinear:
$f(\partial)g$ does not conjugate the coefficients of $f$.
For homogeneous polynomials of equal degree, independent standard
complex Gaussian coordinates give
$f(\partial)g=\mathbf E[f(\bar G)g(G)]$. Since determinant
coefficients are real, this yields
\[
 F_n(\Phi)=\mathbf E\!\left[
             \overline{\det G}\,\det\Phi(G)\right].
\]
Here $\Phi$ acts on all complex matrices; $G$ is not constrained to
be Hermitian.

Du's map uses a nine-by-nine input written in three-by-three blocks,
\[
 \mathcal Y=
 \begin{pmatrix}X&W&U_{13}\\V&Y&U_{23}\\U_{31}&U_{32}&Z\end{pmatrix},
 \qquad S=\begin{pmatrix}X&W\\V&Y\end{pmatrix},
 \qquad J=\begin{pmatrix}0&I_3\\-I_3&0\end{pmatrix}.
\]
The Breuer--Hall map and the required nine-dimensional map are
\begin{equation*}
 \begin{aligned}
 B_6(S)&=\operatorname{tr}(S)I_6-S-JS^T J^*,\\
 \Phi_t(\mathcal Y)
   &=\operatorname{diag}\!\left(
        \operatorname{diag}(Z,Z^T)+tB_6(S),\,X+Y^T\right),
        \qquad t\geq0 .
 \end{aligned}
 \tag{H.29}\label{eq:appH-Chern-nine-map}
\end{equation*}
Superscript $T$ denotes transpose and $*$ denotes conjugate transpose;
the first operation is complex-linear. For $S=ww^*$, the vectors
$w$ and $J\bar w$ are orthogonal and have equal norm. Thus $B_6(ww^*)$
is positive semidefinite, being the norm-squared multiple of the
projection onto their common orthogonal complement. Spectral
decomposition proves the same for every positive semidefinite $S$.
This is the positive-map argument of
\cite[Theorem~4; actual pp.~8--9]{Hall06PositiveMapsV1}.
The principal blocks $Z,X,Y$ of a positive semidefinite $\mathcal Y$
are positive semidefinite, so \eqref{eq:appH-Chern-nine-map} also
preserves that cone. It ignores the four corner blocks $U_{ij}$.

Set $A=X+Y^T$ and $\mathcal R(A)=\operatorname{tr}(A)I_3-A$.
Write $W^T-W=[k]$ and $V^T-V=[l]$, using
\[
 [a]=
 \begin{pmatrix}0&a_3&-a_2\\-a_3&0&a_1\\a_2&-a_1&0\end{pmatrix}.
\]
The Schur-complement calculation in the cited source gives the
polynomial identity
\begin{equation*}
 \det\Phi_t(\mathcal Y)=
 \det A\left[
       \det\bigl(Z+t\mathcal R(A)\bigr)
       +t^2 k^T\bigl(Z+t\mathcal R(A)\bigr)l\right]^2.
 \tag{H.30}\label{eq:appH-Chern-block-determinant}
\end{equation*}
It is first calculated for an invertible middle matrix and then
extends to singular matrices as a polynomial identity.

The ordinary Cayley identity
\cite[Theorem~2.1; printed p.~477]{CaraccioloSokalSportiello13CayleyPublished}
gives, with $D=\det(\partial_A)$ in nine independent entries,
\[
 D(\det A)^2=24\det A,\qquad
 D\bigl((\det A)A_{ij}\bigr)=8A_{ij}.
\]
After differentiating the $Z$ variables in
\eqref{eq:appH-Chern-block-determinant}, the remaining calculation is
\begin{equation*}
 \begin{aligned}
 F_9(\Phi_t)
  &=t^3\left[
      24D^2\bigl(\det A\det\mathcal R(A)\bigr)
       +32D\sum_{i,j=1}^3
          \frac{\partial}{\partial a_{ij}}
             \bigl(\det A\,\mathcal R(A)_{ij}\bigr)\right]\\
  &=-1536\,t^3 .
 \end{aligned}
 \tag{H.31}\label{eq:appH-Chern-negative-coefficient}
\end{equation*}
The first contribution is zero: $\det\mathcal R(A)
=\operatorname{tr}(A)e_2(A)-\det A$ and
$D\det\mathcal R(A)=0$, followed by apolar transfer and Cayley.
For the second contribution, the exact Gaussian identity is
\[
 \mathbf E\!\left[
 |\det G|^2 g_{pq}\overline{g_{rs}}\right]
      =8\delta_{pr}\delta_{qs}.
\]
Its weighted trace-square and Frobenius-square moments are respectively
$24$ and $72$. Their difference is $-48$, and the factor $32$ in
\eqref{eq:appH-Chern-negative-coefficient} gives $-1536$.
This is a finite algebraic calculation, with no probabilistic
approximation of the sign.

\subsection{Strict positivity and global realization}
\label{sec:appH-Chern-global}

At $t=1$ the map is only positive semidefinite preserving. For
$r\geq9$, extend it by the scalar diagonal inputs and perturb:
\begin{equation*}
 \Lambda_{r,\eta}(\mathcal Y)=
 \operatorname{diag}\!\left(
    \Phi_1(\mathcal Y_{[9]}),y_{10,10},\ldots,y_{r,r}\right)
       +\eta\operatorname{tr}(\mathcal Y)I_r .
 \tag{H.32}\label{eq:appH-Chern-strict-perturbation}
\end{equation*}
For $r=9$ there are no appended scalar entries. At $\eta=0$ the
determinant pairing factors and remains $-1536$. At every $\eta>0$
the map sends nonzero positive semidefinite matrices to positive
definite matrices. Since $F_r(\Lambda_{r,\eta})$ is a real polynomial
in $\eta$, it remains negative for sufficiently small positive
$\eta$, separately for each finite rank.

The global step in \cite[Proposition~3.2; actual pp.~12--13]{Du26GriffithsTopChernV1}
retains the prescribed negative jet on a small ball and uses the
direct-sum Fubini--Study metric outside a slightly larger ball.
For clarity, write $s=|z|^2$,
\[
 h_0=(1+s)^{-1}I,\qquad
 h_1=I-2\pi\Lambda_{r,\eta}(zz^*),\qquad H=h_1-h_0.
\]
Strict rank-one positivity and the vanishing first jets give a
uniform lower bound for the negative mixed Hessian of
$(1-t)h_0+th_1$, $0\leq t\leq1$, on a sufficiently small ball.
They also give a uniform positive lower bound for the metric
matrices on that ball.

Take a smooth cutoff
$\chi_L(s)=\theta(1+L^{-1}\log(s/\rho^2))$, equal to one for
$s\leq\rho^2e^{-L}$ and zero for $s\geq\rho^2$. Near $s=0$
it is interpreted by its constant value one. The bounds are
\[
 H=O(s),\qquad \partial H=O(\sqrt s),\qquad
 \partial\chi_L=O\!\left(\frac1{L\sqrt s}\right),\qquad
 \partial\bar\partial\chi_L=O\!\left(\frac1{Ls}\right).
\]
The extra mixed-Hessian terms of $h=h_0+\chi_LH$ are therefore
$O(1/L)$. For large fixed $L$ they preserve the positive lower
curvature bound; the quadratic first-derivative term in
$hR^E(u,\bar u)$ is nonnegative. This proves the required local
patching fact. It does not assert general convexity of
Griffiths-positive metrics. Smooth extension by $h_0$ gives the
global metric on $\mathcal O(1)^{\oplus r}$, and continuity of its
top Chern form gives the negative open set.

\subsection{Low ranks and the comparison boundary}
\label{sec:appH-Chern-low-ranks}

The read literature gives a positive answer for top Chern forms
in ranks two and three. The rank-two matrix case is recorded in
\cite[Proposition~1.4(c)]{Finski22CharacteristicPublished}.
For rank three,
\cite[Theorems~0.3 and~2.11]{Wan26ThirdChernV2}
normalizes a strictly positive map by invertible input and output
scaling. The exact scaling input is
\cite[Theorem~1.2]{Idel16MatrixScalingV1}; positivity improving is
a stronger hypothesis than mere positive semidefinite preservation.

Under the normalization
$\sum_iB_{ii}=3I$ and $\operatorname{tr}(B_{ij})=3\delta_{ij}$,
put $C(\xi)=(\xi^*B_{ij}\xi)_{i,j=1}^3$ for $\xi$ on the unit
sphere in $\mathbf C^3$. This positive definite matrix has trace
three. For the normalized mixed-discriminant scalar $\mathcal D$,
Wan's spherical identity is
\begin{equation*}
 \mathcal D=\int_{S^5}
   \bigl(10\det C(\xi)+27-12e_2(C(\xi))\bigr)\,d\mu(\xi)
       \geq\int_{S^5}\det C(\xi)\,d\mu(\xi)>0 .
 \tag{H.33}\label{eq:appH-Chern-rank-three}
\end{equation*}
The inequality is Schur's eigenvalue inequality with eigenvalue
sum three. This normalized scalar is $F_3/3!$ in the determinant
convention of \eqref{eq:appH-Chern-det-pairing}.
The all-Schur conclusion in the cited Theorem~0.3 also retains
rank three and a complex threefold, with its stated external
signed-Segre and lower-degree positivity inputs.

For $c_3$ of a higher-rank bundle, the pointwise principal
three-by-three blocks of the ambient curvature still satisfy the
rank-one positivity test. Applying the three-dimensional matrix
result to each block and restricting to three-dimensional tangent
planes proves the weak positivity stated in the cited Theorem~0.5.
This is a curvature compression argument. An arbitrary holomorphic
subbundle with its induced metric need not retain Griffiths positivity,
because of its second fundamental form.

Du's final paragraph asks for the least rank admitting a negative
top Chern form. The positive top-degree results just recalled and
the rank-nine construction leave ranks four through eight as the
candidate ranks in this source-scoped question. The present account
does not claim an exhaustive audit of all intervening literature.
The stronger Nakano and dual Nakano hypotheses give positive Schur
forms by \cite[Theorem~1.1]{Finski22CharacteristicPublished};
they must not be replaced by Griffiths positivity in that statement.

The scalar curvature products of
\Cref{chap:cld-metrized-bundles} have their own positivity and
intersection-theoretic construction. The complex example above
supplies a useful comparison of meanings of positivity, rather
than a non-Archimedean higher Chern form. Neither a new singular
product nor a general Siu decomposition follows from it.
Exact editions and proof boundaries are recorded in
\Cref{sec:provenance-Griffiths-Chern}.

\section{Boundary rigidity of Bergman metrics}
\label{sec:appH-Bergman-rigidity}
\index{Bergman metric!boundary rigidity}
\index{Carleman estimate!Bergman rigidity}

The compact cscK problem of \Cref{chap:ytd-bridge} has a useful
noncompact comparison. The Bergman metric of a domain is fixed by
its holomorphic $L^2$ space; one is not varying all K\"ahler metrics
in a cohomology class. Its strongly pseudoconvex boundary asymptotics
can therefore impose a rigidity that an arbitrary cscK metric does
not have. Ebenfelt--Ganguly \cite[Theorems~1.1--1.3]{EbenfeltGanguly26BergmanV1}
prove the following attributed result. If a bounded connected
pseudoconvex domain $\Omega\subset\mathbf C^n$, $n\geq2$, has one
smooth strongly pseudoconvex boundary point, its Bergman metric is
extremal if and only if it has constant scalar curvature, if and
only if it is K\"ahler--Einstein. Neither the whole boundary's
smoothness nor its strong pseudoconvexity is assumed.

Their Theorem~1.5 starts with the extremal or cscK condition only
on a collar of such a boundary point, also for possibly unbounded
domains. The Einstein conclusion is on the locus where the
diagonal Bergman kernel is positive. It does not assert that every
such domain is a ball. For the stronger global smoothly bounded
strongly pseudoconvex hypotheses, the ball characterization is the
separately attributed theorem of
\cite[Theorem~1.1]{HuangXiao21BergmanPublished}.

\subsection{The normalization and the scalar defect}
\label{sec:appH-Bergman-normalization}

Use Euclidean volume to define $K(z)=K_\Omega(z,z)$ and work first
where $K>0$ and the Bergman metric is nondegenerate. Write
\begin{equation*}
 \begin{gathered}
 g_{\alpha\bar\beta}=\partial_\alpha\partial_{\bar\beta}\log K,
 \quad G=\det g,\quad B=G/K,\quad
 C_n=\frac{(n+1)^n\pi^n}{n!},\\
 v=K^{-1/(n+1)},\qquad u=\log(B/C_n),\qquad
 \Delta_g f=g^{\alpha\bar\beta}f_{\alpha\bar\beta}.
 \end{gathered}
 \tag{H.34}\label{eq:appH-Bergman-defect}
\end{equation*}
Here the source K\"ahler form is
$\omega_B=\rmi\partial\bar\partial\log K$ and its Ricci form
is $-\rmi\partial\bar\partial\log G$. The source scalar curvature
is the complex trace, without an additional real factor two.
For a positive-inside defining function $r$, use
\begin{equation*}
 J(f)=(-1)^n\det
 \begin{pmatrix}f&f_{\bar\beta}\\ f_\alpha&f_{\alpha\bar\beta}\end{pmatrix},
 \qquad B=(n+1)^nJ(v),\qquad
 \operatorname{Scal}_g+n=-\Delta_g u.
 \tag{H.35}\label{eq:appH-Bergman-MA}
\end{equation*}
On the unit ball, $r=1-|z|^2$, $J(r)=1$ and
$K=n!/(\pi^n r^{n+1})$; these formulas give $B=C_n$.

To translate to this book, take
$\widehat\omega_B=\ddc\log K=\omega_B/(2\pi)$.
Its Ricci form in the book convention is
$\operatorname{Ric}(\widehat\omega_B)=-\ddc\log G$.
Thus the Einstein conclusion is
$\operatorname{Ric}(\widehat\omega_B)=-\widehat\omega_B$, and
its complex-trace scalar is $-n$. The determinant of the rescaled
metric coefficients is $(2\pi)^{-n}G$. The constant $C_n$ in
\eqref{eq:appH-Bergman-defect} belongs to $g$, so it must also be
rescaled if that determinant convention is changed.

\subsection{Two different critical orders}
\label{sec:appH-Bergman-indicial}

For a Fefferman defining function, $J(r)=1+O(r^{n+1})$ and
the local kernel expansion has the form
$K=\phi r^{-n-1}+\psi\log r$. Smooth localization of the kernel
is supplied by \cite[Theorem~1.1]{HsiaoHuangLi26LocalizationV2}.
The associated full polyhomogeneous expansion is an asymptotic
series in powers of $r$ and $\log r$, with controlled differentiated
remainders at every order. It is an input to the source proof. Finite H\"older boundary regularity
alone cannot replace this all-orders information.

The MA variation and the harmonic scalar defect have distinct
indicial polynomials. Their leading terms, in the conventions above,
are
\begin{equation*}
 \begin{aligned}
 \delta J(r)[a r^m]&=m(n+2-m)a r^{m-1}+O(r^m),\\
 \Delta_g(a r^\lambda)&=
       \frac{\lambda(\lambda-n)}{n+1}a r^\lambda
       +\text{higher-order terms}.
 \end{aligned}
 \tag{H.36}\label{eq:appH-Bergman-indicial}
\end{equation*}
The first formula is the linearized leading variation; $a$ is
a smooth boundary coefficient. Its critical perturbation exponent
is $n+2$. A smooth error in $K$ produces a change of order
$r^{n+2}$ in $v$, and this critical cancellation improves the
corresponding error in $J(v)$. The harmonic roots are instead
zero and $n$. For full polyhomogeneous harmonic $u$ tending to
zero, its first possible positive term is $r^n$. If its order
is strictly greater than $n$, the indicial recursion makes it flat.
Nonconstant leading polynomials in $\log r$ are also excluded
at the simple root $n$.

These signs and factors can be checked without a boundary
parametrix: on the ball the exact complex-Laplacian identity is
\begin{equation*}
 \Delta_g(r^m)=\frac{m}{n+1}r^m(m-n-mr).
 \tag{H.37}\label{eq:appH-Bergman-ball-Laplacian}
\end{equation*}
It fixes the complex-versus-real normalization when comparing
the auxiliary cscK arguments of
\cite{HsiaoHuangLi26LocalizationV2,LiLiu26CSCBergmanV1}.
Their vanishing arguments use the nonzero coefficient beyond the
harmonic root; that root is unaffected by a nonzero overall factor.

Dimension two needs a separate cancellation. As recorded in the
EG proof, let $\eta_1$ be the first logarithmic coefficient of
the local MA expansion. The kernel-invariant input gives
$c_0=2/\pi^2$,
$\phi=c_0+O(r^3)$ and $\psi=-3c_0\eta_1+O(r)$. Put $C=c_0^{-1/3}$. Then
$v=Cr+r^4\alpha+r^4\beta\log r+O_{\rm phg}(r^7)$, with
$\beta=C\eta_1+O(r)$. Graham's checked MA formula
\cite[equation~(4.11)]{Graham87HigherAsymptoticsPublished}
then gives
\begin{equation*}
 \begin{aligned}
 J(r)&=1+4\eta_1r^3+O(r^4),\\
 J(v)&=C^3J(r)-4C^2\beta r^3J(r)+O_{\rm phg}(r^4).
 \end{aligned}
 \tag{H.38}\label{eq:appH-Bergman-dimension-two}
\end{equation*}
The $r^3$ coefficients cancel. Both $u$ and the scalar defect
have order $r^4$, already beyond the harmonic root two.
For $n\geq3$ the initial order is $r^2$; harmonicity first raises
it to $r^n$, and the source's local curvature-invariant argument
then supplies sphericity and the enhanced order $r^{n+2}$.
The weight-two coefficient underlying that step is recorded in
\cite[Proposition~2.1]{HuangXiao21BergmanPublished}; it vanishes
identically in dimension two and cannot substitute for
\eqref{eq:appH-Bergman-dimension-two} there.

The full MA expansion is associated with
\cite[Theorem~8.18]{LeeMelrose82BoundaryPublished} and the
local logarithmic-flatness inputs cited in
\cite{EbenfeltXiao24LogarithmicFlatness}. Importing those inputs
requires their hypotheses. In particular, flatness of the MA
obstruction alone does not imply a spherical boundary:
\cite[Proposition~4.14]{Graham87HigherAsymptoticsPublished}
constructs nonspherical real-analytic examples with vanishing
obstruction. The spherical step in the EG proof uses the stronger
scalar-defect information, rather than that false converse.

\subsection{Why flatness still needs unique continuation}
\label{sec:appH-Bergman-Carleman}

The extremal condition is first reduced to cscK. Its holomorphic
scalar-gradient vector field tends to zero on the strongly
pseudoconvex patch; one-variable boundary uniqueness on transverse
slices, followed by the identity theorem, makes it zero. Boundary
scalar asymptotics then fix the constant as $-n$, so $u$ is harmonic.
The preceding expansion makes $u$ flat on the patch. A flat smooth
function need not vanish, so this is not yet the Einstein conclusion.

EG resolves this last step by a Carleman estimate on an
asymptotically complex hyperbolic collar. Put $x=\sqrt r$,
use the geodesic defining function supplied locally by
\cite[Lemma~3.1, author copy]{GuillarmouSaBarretoACHAuthorCopy},
denoted by $\widetilde x$, and let $s=-\log\widetilde x$.
For the real source metric
$g_{\mathbf R}=2\operatorname{Re}(g_{\alpha\bar\beta}
 dz_\alpha\otimes d\bar z_\beta)$, set
$\widehat g=g_{\mathbf R}/[2(n+1)]$.
Its radial leading term is $ds^2$. Tangential horizontal and
Reeb directions grow respectively as $e^{2s}$ and $e^{4s}$,
and volume as $e^{2ns}$. The complex and real Laplacians satisfy
$\Delta_{\widehat g}=4(n+1)\Delta_g$.

Choose a compact boundary patch with a smooth positive function
$\delta$ tending to zero at its edge. On the bent end
$s>S-\log\delta$, the weight is $\tau=s-S+\log\delta$.
The inequality $e^{-s}\leq e^{-S}\delta$ bounds the lateral cutoff
derivatives in the scaled geometry. Positive tangential Hessian
and uniform coordinate bounds let the source apply the local
full-symbol estimate of
\cite[Theorem~28.2.3]{HormanderIV09Reprint} and globalize it:
\begin{equation*}
 \gamma^3\|e^{\gamma\tau}w\|_{L^2(\widehat g)}^2
 +\gamma\|e^{\gamma\tau}\nabla w\|_{L^2(\widehat g)}^2
 \leq C\|e^{\gamma\tau}(-\Delta_{\widehat g})w\|_{L^2(\widehat g)}^2
 \tag{H.39}\label{eq:appH-Bergman-Carleman}
\end{equation*}
for compactly supported smooth $w$ in the bent end and large
$\gamma$. Uniformity over the scaled charts and commutator
absorption are part of EG's argument; the cited local theorem
alone does not give a noncompact estimate.

Apply this to cutoff multiples of $u$. Flatness gives decay
$O(e^{-Ms})$ for every fixed $M$ and the required derivatives.
At fixed $\gamma$, choose $M>n+\gamma$ to offset the volume growth.
First remove the lateral and radial cutoffs using the resulting
integrable bounds. Then let $\gamma$ increase, comparing the fixed
lower cutoff strip with a deeper region. This forces $u=0$
on a collar. Sending the weight to infinity before controlling
the radial tail would not justify this conclusion.

Finally, put
$H_{\alpha\bar\beta}=K K_{\alpha\bar\beta}-K_\alpha K_{\bar\beta}$.
On that collar $B=C_n$, hence
\begin{equation*}
 \det H=C_n K^{2n+1}.
 \tag{H.40}\label{eq:appH-Bergman-denominators}
\end{equation*}
This is an identity between real-analytic functions on the
connected ambient domain, so it propagates without dividing by
$K$. On $\{K>0\}$, the Gram representation gives a semipositive
Bergman matrix; \eqref{eq:appH-Bergman-denominators} gives a
positive determinant and therefore strict positivity. The
Einstein conclusion follows there. The reasoning uses analytic
kernel identities and strongly pseudoconvex asymptotics specific
to this complex setting. It supplies no non-Archimedean
unique-continuation or Siu/envelope theorem. Exact source scopes
and external inputs are recorded in
\Cref{sec:provenance-Bergman-rigidity}.

\section{Finite ball quotients and their Bergman kernels}
\label{sec:appH-ball-quotients}
\index{Bergman metric!finite ball quotients}
\index{quotient!Bergman kernel}

The preceding section concerns rigidity detected at a boundary
patch. A complementary question asks whether an Einstein Bergman
metric can conceal an isolated quotient singularity. Ganguly--Sahi
\cite[Theorem~1.3]{GangulySahi26BallQuotientsV1} answer this for
finite ball quotients in every complex dimension $N\geq2$:
if $\Gamma\subset U(N)$ is nontrivial and every nonidentity
element acts without fixed points on the unit sphere, the
Bergman metric on the regular part of $\mathbb B^N/\Gamma$
is not K\"ahler--Einstein. Sphere freeness means
$1\notin\operatorname{Spec}(\gamma)$ for $\gamma\ne I$.
The action is free on the punctured ball, while the origin
is fixed by all of $\Gamma$. For $N\geq2$ its nontrivial
quotient has an isolated singularity at the origin.

This removes the abelian restriction of
\cite[Theorem~1.4]{EbenfeltXiaoXu21BallQuotientsAdvance}
and the dimension-two restriction of
\cite[Theorem~1]{GangulySinha24SteinV3}.
The argument in the controlling 2026 version uses rational
kernel identities rather than the classification of finite
sphere-free groups. We explain that argument with its source
normalizations; the comparisons refer to these exact versions.

\subsection{Top forms, averaging and the quotient Jacobian}
\label{sec:appH-ball-quotient-kernel}

On a singular complex space the intrinsic Bergman Hilbert space
consists of square-integrable holomorphic top forms on its
regular locus. On the covering ball write such a form as
$f(z)\,dz_1\wedge\cdots\wedge dz_N$. Its invariance condition is
$f(\gamma z)\det\gamma=f(z)$. Thus its coefficient transforms
with a determinant weight. The Hilbert space of untwisted
invariant scalar functions would give a different kernel.

Normalize the ball coefficient by
$K_{\mathbb B}(z,w)=N!/[\pi^N(1-z^T\bar w)^{N+1}]$.
If $K_\Gamma$ is the coefficient kernel of the orthogonal
projection onto these invariant top forms, and $k_Q$ is the
intrinsic quotient coefficient in local coordinates, then
\begin{equation*}
 \begin{aligned}
 K_\Gamma(z,w)
   &=\frac1{|\Gamma|}\sum_{\gamma\in\Gamma}
         K_{\mathbb B}(\gamma z,w)\det\gamma,\\
 (k_Q\circ\Pi)(z)\,|\operatorname{Jac}\Pi(z)|^2
   &=|\Gamma|K_\Gamma(z,z).
 \end{aligned}
 \tag{H.41}\label{eq:appH-ball-quotient-transform}
\end{equation*}
Here $\Pi$ is locally a biholomorphism away from the origin,
and its Jacobian uses those local quotient coordinates.
The sum transformation in
\cite[Theorem~2.3]{EbenfeltXiaoXu21BallQuotientsAdvance}
and the averaged projection in
\cite[Lemma~5.1]{HuangLiBergmanSteinAuthorCopy}
therefore have different constants. Positive constant factors
leave a logarithmic metric unchanged, but affect its determinant
equation. The Jacobian need not have modulus one: it belongs
to the quotient map, not to an element of the unitary group.
Its logarithmic square is pluriharmonic off the branch locus,
so the metric still pulls back correctly.

For independent column variables $z,w\in\mathbf C^N$, put
\begin{equation*}
 \begin{gathered}
 h_\gamma=1-z^T\gamma w,\qquad
 \varphi_\Gamma=\sum_{\gamma\in\Gamma}
                \frac{\det\gamma}{h_\gamma^{N+1}},\\
 M(u)=\begin{pmatrix}u&u_w^T\\u_z&u_{zw}\end{pmatrix},\qquad
 F_\Gamma=\det M(\varphi_\Gamma)
                  -(N+1)^N\varphi_\Gamma^{N+2}.
 \end{gathered}
 \tag{H.42}\label{eq:appH-ball-quotient-rational}
\end{equation*}
Second-variable conjugation in the covering formula means that
the quotient by $\Gamma$ uses $\varphi_{\bar\Gamma}$ on
$(z,\bar z)$, where $\bar\Gamma=\{\bar\gamma:\gamma\in\Gamma\}$;
see \cite[Proposition~2.1]{GangulySahi26BallQuotientsV1}.
Conjugating the group preserves all its hypotheses. After this
replacement, the Einstein condition requires $F_\Gamma=0$.
The bordered determinant identity is
$\det(\partial\bar\partial\log\varphi)
=\det M(\varphi)/\varphi^{N+1}$ on the positive diagonal.
This is a determinant built from the kernel, distinct from the
positive-defining-function determinant $J(r)$ in
\eqref{eq:appH-Bergman-MA}; its sign convention is not transferred.

The constant $(N+1)^N$ is an input, not an arbitrary choice of
scale. In \cite[Lemma~4.2]{EbenfeltXiaoXu21BallQuotientsAdvance},
the punctured ball is simply connected for $N\geq2$.
Hartogs extension and the identity summand's uniform boundary
asymptotics allow maximum modulus to be applied to a nowhere-zero
holomorphic exponential and its reciprocal. They eliminate
the pluriharmonic factor and force the source Ricci constant
$-1$. Coordinate metric coefficients here are
$g_{i\bar j}=\partial_i\partial_{\bar j}\log k$;
the appendix form is $\ddc\log k$, with its already fixed factor
$1/(2\pi)$. This calculation does not assign a negative Einstein
constant to an arbitrary K\"ahler metric.

An identity on the anti-diagonal $(z,\bar z)$ extends to a
rational identity in the independent variables, because that
anti-diagonal is maximally totally real. It is consequently
enough to find a complex pole obstructing $F_\Gamma\equiv0$.
The chosen complex points can lie outside the physical ball
diagonal; they do not assert a pole of its interior Bergman kernel.

\subsection{The first restriction and a second pole}
\label{sec:appH-ball-quotient-poles}

Write $h=h_I$, $\kappa_I=h^{-N-1}$ and
$\psi=\varphi_\Gamma-\kappa_I$. The identity summand has a
particularly simple inverse
\cite[Lemma~3.1]{GangulySahi26BallQuotientsV1}:
\begin{equation*}
 \begin{gathered}
 M(\kappa_I)^{-1}=(N+2)h^{N+1}A_0-h^{N+2}A,\\
 A_0=\begin{pmatrix}1&0\\0&0\end{pmatrix},\qquad
 A=\begin{pmatrix}N+1&z^T\\w&(wz^T-I)/(N+1)\end{pmatrix},\\
 \det M(\kappa_I)=(N+1)^N h^{-(N+1)(N+2)}.
 \end{gathered}
 \tag{H.43}\label{eq:appH-ball-quotient-inverse}
\end{equation*}
The rank-one identity behind it is
$(I+h^{-1}wz^T)^{-1}=I-wz^T$.

Let $H_\gamma=\{h_\gamma=0\}$ and
$H_\Gamma=H_I\setminus\bigcup_{\gamma\ne I}H_\gamma$.
The polynomial $h$ is irreducible and no nonidentity divisor
contains $H_I$, so $H_\Gamma$ is nonempty. There $\psi$ has
no pole. With $G_\Gamma=\operatorname{tr}(A M(\psi))$, the
determinant expansion and the expansion of
$\varphi_\Gamma^{N+2}$ give
\begin{equation*}
 \frac{h^{N(N+2)}}{(N+1)^N}F_\Gamma
       =-G_\Gamma+h^N(R-S),
 \tag{H.44}\label{eq:appH-ball-quotient-first-pole}
\end{equation*}
where $R,S$ are regular along $H_\Gamma$.
Thus $F_\Gamma\equiv0$ would force $G_\Gamma|_{H_\Gamma}=0$;
see \cite[Proposition~3.3]{GangulySahi26BallQuotientsV1}.
The leading identity-summand contribution cancels before this
restriction. Keeping that cancellation is essential.

Suppose the group contains a nonscalar $\eta$. Source
Lemmas~4.1--4.2 choose $p=(z,w)$ on exactly $H_I$ and
$H_\eta$. The linear algebra uses the invertibility of
$I-\eta$ and unitarity: for $\delta\notin\{I,\eta\}$,
$(I-\eta)^{-1}(I-\delta)$ cannot be scalar. Otherwise the
eigenvalues of the affine combination
$\delta=(1-\mu)I+\mu\eta$ force two circles with three
distinct common points to coincide, giving $\mu=1$ and
$\delta=\eta$. One can then choose a row outside finitely
many left eigenspaces and a vector in its kernel outside
finitely many proper kernels. This requires $N-1\geq1$.
The resulting pair satisfies
$z^Tw=z^T\eta w=1$ and $z^T\delta w\ne1$ for every other
group element.

Choose $u$ with $z^Tu=0$ and $z^T\eta u\ne0$.
The curve $p_\epsilon=(z,w+\epsilon u)$ stays on $H_I$ and,
for small nonzero $\epsilon$, belongs to $H_\Gamma$.
Only the $\eta$ summand creates the second pole, with
\begin{equation*}
 \lim_{\epsilon\to0}h_\eta(p_\epsilon)^{N+3}
                  G_\Gamma(p_\epsilon)
       =\det\eta\,(N+2)(1-z^T\eta^2w)\ne0.
 \tag{H.45}\label{eq:appH-ball-quotient-second-pole}
\end{equation*}
Indeed $\eta^2$ is neither $I$ nor $\eta$: a sphere-free
unitary involution would be $-I$, which is scalar.
The chosen pair therefore makes the last factor nonzero.
This is \cite[Proposition~4.3]{GangulySahi26BallQuotientsV1}.
It contradicts the vanishing forced by
\eqref{eq:appH-ball-quotient-first-pole}.

\subsection{Scalar groups and the dimension-one exception}
\label{sec:appH-ball-quotient-scalar}

If every element is scalar, the finite nontrivial group is
$\{\lambda I:\lambda^\ell=1\}$ for $\ell\geq2$.
Its potential is a function of $s=z^Tw$:
\begin{equation*}
 \begin{gathered}
 f(s)=\sum_{\lambda^\ell=1}\frac{\lambda^N}{(1-\lambda s)^{N+1}}
       =\frac{P(s)}{(1-s^\ell)^{N+1}},\\
 \det M(f)=(f')^{N-1}\bigl[ff'+s(ff''-(f')^2)\bigr].
 \end{gathered}
 \tag{H.46}\label{eq:appH-ball-quotient-scalar}
\end{equation*}
The determinant weight $\lambda^N$ controls the root-of-unity
filter. The genuine poles have order $N+1$, so $P$ is nonzero
at every root of $1-s^\ell$. The expansion at infinity gives
$\deg P=N(\ell-1)$, while the expansion at zero gives
$\operatorname{ord}_0P\leq\ell-1$. For $N\geq2$ there is
therefore a zero $b\ne0$ away from the poles. Write
$f(s)=(s-b)^r u(s)$, with $r\geq1$ and $u(b)\ne0$, and
$Y=f'/f$. Source Appendix~A obtains
\begin{equation*}
 \frac{F_\Gamma}{f^{N+2}}
  =\frac{Y^{N-1}(sY)'}{f}-(N+1)^N
  =-\frac{b r^N}{u(b)}(s-b)^{-(N+r+1)}
         +\text{terms of lower pole order}.
 \tag{H.47}\label{eq:appH-ball-quotient-scalar-pole}
\end{equation*}
The leading coefficient is nonzero. No simple-zero assumption
is needed. As a direct check, for $N=2$, $\ell=2$ one has
$f=2(1+3s^2)/(1-s^2)^3$; at $b=\rmi/\sqrt3$ the pole
in $F_\Gamma/f^4$ has order four and coefficient $-16/81$.
This finite calculation illustrates the source argument.

Dimension one is different. The degree comparison no longer
forces a nonzero zero, and the quotient map $z\mapsto z^\ell$
identifies the centered finite disk quotient with a smooth disk.
Its Bergman metric is Einstein. Degeneration of the branched
pullback at zero does not make the quotient singular.
For $N=1$, $\ell=2$, the rational function
$f=4s/(1-s^2)^2$ has $F_\Gamma\equiv0$, as required.

\subsection{The global Stein consequence and what remains open}
\label{sec:appH-ball-quotient-Stein}

Ganguly--Sahi \cite[Theorem~1.2]{GangulySahi26BallQuotientsV1}
conclude that a normal Stein space of dimension $N\geq2$,
with isolated singularities and compact smooth strongly
pseudoconvex boundary CR equivalent to an algebraic CR manifold
in a complex Euclidean space, has Einstein Bergman metric
on its regular locus precisely when it is biholomorphic to
$\mathbb B^N$. The global reduction first makes the boundary
spherical, then realizes the space as a finite sphere-free
ball quotient. The inputs are
\cite[Theorem~1.1 and Corollary~4.2]{HuangLiBergmanSteinAuthorCopy},
with the algebraic spherical uniformization discussed in
\cite[Theorem~3.1 and Corollary~3.3]{Huang06CRLinksAuthorCopy}.
The finite-quotient obstruction then forces the group to be trivial.

The algebraic CR hypothesis is part of this reduction.
Conjecture~1.1 of the controlling 2026 paper asks for its removal
while retaining the normal isolated-singularity and compact
smooth strongly pseudoconvex boundary setting. Real-analytic
spherical boundaries with infinite fundamental group, recalled
in Huang's survey, obstruct an automatic finite-quotient
argument; they are not asserted to have Einstein Bergman metrics
and do not refute that conjecture. The finite-type extensions
in source Corollaries~4.4--4.5 require separate external
localization and asymptotic inputs. In dimension at least three
the stated hypotheses also include finite D'Angelo $1$-type
and h-extendibility at every weakly pseudoconvex boundary point.

The two successive pole orders above are valuations of
rational functions on a complexified kernel space. They do
not construct a non-Archimedean psh or finite-quotient descent
theory, nor settle the book's singular-local open problems.
The exact source versions, partial dependency scopes and
auxiliary convention corrections are recorded in
\Cref{sec:provenance-ball-quotients}.

\section{Pointwise envelopes with globally bounded minorants}
\label{sec:appH-bounded-minorants}
\index{Perron--Bremermann envelope!globally bounded minorants}
\index{strong majorant!pole set and pointwise recovery}

Boundary regularity and the admissible family play separate roles.
Wikstr\"om's revised paper
\cite{Wikstrom26BoundedMinorantsV2} shows that even a bounded B-regular
domain can lose interior values when every competitor must have a
global upper bound. The obstacle may be finite and smooth throughout
the domain. We explain the exact restricted envelope, its regularization
and the role of poles in a strong majorant. These are attributed complex
and real potential-theoretic results; their version and dependency scopes
are recorded in \Cref{sec:provenance-bounded-minorants}.

For a domain $\Omega$, an obstacle $f$ and a class $\mathcal A$ of
subharmonic or plurisubharmonic functions, write
\begin{equation*}
\begin{aligned}
 V_f^{\mathcal A}(p)
   &=\sup\{u(p):u\in\mathcal A,\ u\leq f,\ \sup_\Omega u<+\infty\},\\
 (V_f^{\mathcal A})^*(p)
   &=\limsup_{\Omega\ni q\to p}V_f^{\mathcal A}(q).
\end{aligned}
\tag{H.48}\label{eq:appH-restricted-envelope}
\end{equation*}
The supremum of an empty family is $-\infty$. Each admissible function
has its own finite upper bound. A common prescribed cap would define
a different family. In this section the identically $-\infty$ function
is allowed in the upper-bounded family; it is excluded when two-sided
boundedness is required.

\subsection{What the real subharmonic theorem recovers}
\label{sec:appH-subharmonic-restricted-envelope}

Theorem~1.1 of \cite[actual pp.~1, 3--4]{Wikstrom26BoundedMinorantsV2}
gives the following result, with the global boundedness condition retained.

\begin{proposition}
\label{prop:appH-subharmonic-restricted-envelope}
Let $\Omega\subset\mathbf R^d$ be a domain, $d\geq2$, and let
$f:\Omega\to[-\infty,+\infty)$ be upper semicontinuous. For
$\mathcal A=\operatorname{SH}(\Omega)$, the envelope in
\eqref{eq:appH-restricted-envelope} equals its upper regularization.
It is subharmonic or identically $-\infty$. The same holds when the
competitors are required to be bounded on both sides.
\end{proposition}

\begin{proof}
The obstacle supplies local upper bounds. Standard subharmonic envelope
theory gives $W=V^*$ subharmonic unless $W\equiv-\infty$, and
upper semicontinuity of $f$ gives $W\leq f$. In the nontrivial case,
Choquet selection and finite maxima give increasing nontrivial
competitors $u_j\to W$ quasi-everywhere. Fix $p$ with $W(p)>-\infty$
and a ball $B$ centered at $p$, compactly contained in $\Omega$.
Its normalized surface harmonic measure $\sigma$ does not charge the
polar exceptional set.

The restrictions of $W,u_j$ to $\partial B$ are integrable.
For completeness, truncating a nontrivial subharmonic function below
at $-m$ and applying Poisson comparison shows that a boundary integral
of $-\infty$ would force that function to be $-\infty$ throughout $B$,
and hence throughout the connected domain. Local upper bounds then
give the asserted integrability. Let $P_B$ denote Poisson extension and
put $h_j=P_B[u_j]-P_B[W]$. This is an ordinary real harmonic function,
$h_j\leq0$, and $h_j(p)=\int_{\partial B}(u_j-W)\,d\sigma\to0$
by monotone convergence applied to the integrable differences.

Poisson comparison gives $W-P_B[W]\leq0$. Therefore
$s_j=W+h_j\leq W$, and the upper boundary limits of $s_j$ are at most
$u_j$: the first term in $s_j=(W-P_B[W])+P_B[u_j]$ is nonpositive,
and the usc boundary data bound the second. The usual subharmonic
pasting lemma joins $\max(u_j,s_j)$ on $B$ to $u_j$ outside $B$.
The result is below $W\leq f$, is globally bounded above, and preserves
the lower bound of $u_j$ when it has one. Its value at $p$ tends to
$W(p)$, so $V(p)\geq W(p)$. Points where $W=-\infty$ are immediate.
\end{proof}

The conclusion is $V=V^*$, not recovery of every obstacle $f$.
For the disk Poisson kernel, every globally bounded-above subharmonic
minorant is at most zero; its restricted envelope is zero and is
already upper semicontinuous. This older example appears in
\cite[Example~1.2, actual p.~1]{NilssonWikstrom20QuasiboundedV1}.
By contrast, if increasing bounded minorants of a subharmonic $W$
converge to $W$ quasi-everywhere, the proposition recovers $W$
everywhere, as in \cite[Corollary~2.1]{Wikstrom26BoundedMinorantsV2}.
An obstacle allowed to take $+\infty$ needs the separate hypothesis
that its regularized envelope is subharmonic. In $\mathbf C^n$, $n\geq2$,
the real harmonic correction above need not be pluriharmonic and need
not preserve plurisubharmonicity.

\subsection{A B-regular Hartogs domain and an exact defect set}
\label{sec:appH-Hartogs-envelope-defect}

Theorem~1.2 and Remark~3.6 of \cite{Wikstrom26BoundedMinorantsV2}
give an explicit example in every complex dimension $n\geq2$.
For $w\in\mathbf D$ and $z\in\mathbf C^{n-1}$, use the Euclidean norm
and set
\begin{equation*}
\begin{aligned}
 P(w)&=\frac{1-|w|^2}{|1-w|^2}
       =\operatorname{Re}\frac{1+w}{1-w},
       \qquad \Phi(w)=P(w)^2+\frac1{1-|w|^2},\\
 \Omega_n&=\{(z,w):w\in\mathbf D,\ \|z\|<e^{-\Phi(w)}\},\\
 H&=\{0\}\times\mathbf D,\qquad W(z,w)=P(w).
\end{aligned}
\tag{H.49}\label{eq:appH-Hartogs-domain}
\end{equation*}
Here $W$ is positive and pluriharmonic, finite at every interior point,
and unbounded as $w\to1$. The fibres collapse at every point of
$\partial\mathbf D$ because $\Phi\to+\infty$ there. The domain is
bounded and B-regular. Its boundary has a smooth strongly pseudoconvex
side and a cuspidal circle; global smoothness is not asserted.
B-regularity means that every finite continuous boundary datum extends
to a psh function continuous on the closure.

\begin{example}
\label{ex:appH-Hartogs-restricted-envelope}
For $\mathcal A=\PSH(\Omega_n)$ the exact envelope is
\begin{equation*}
\begin{aligned}
 u_\varepsilon&=\max\{0,P(w)+\varepsilon\log\|z\|\},
 &0\leq u_\varepsilon&\leq\frac1{4\varepsilon},\\
 V_W^{\mathcal A}(z,w)
   &=\begin{cases}P(w),&z\ne0,\\0,&z=0,\end{cases}
 &(V_W^{\mathcal A})^*&=W.
\end{aligned}
\tag{H.50}\label{eq:appH-Hartogs-envelope}
\end{equation*}
Set $u_\varepsilon(0,w)=0$. In particular $V(0,0)=0<1=V^*(0,0)$,
although $u_{1/j}\nearrow W$ off the pluripolar analytic disc $H$.
The same formula holds with two-sided bounded competitors.
\end{example}

To check the boundary assertion, put $r=\log\|z\|+\Phi(w)$ near
the side. Its Levi form on a tangent vector $(a,b)$ is
\[
 \frac{\|z\|^2\|a\|^2-|\langle a,z\rangle|^2}{2\|z\|^4}
       +\Phi_{w\bar w}|b|^2,
 \qquad
 \frac{\langle a,z\rangle}{2\|z\|^2}+\Phi_w b=0.
\]
Here $\langle a,z\rangle=\sum_i a_i\bar z_i$.
The first term is nonnegative and the second is positive for $b\ne0$;
if $b=0$, tangency makes $a$ orthogonal to $z$, giving positivity
for $a\ne0$. Indeed
$\Phi_{w\bar w}=2/|1-w|^4+(1+|w|^2)/(1-|w|^2)^3>0$.
Smooth strongly pseudoconvex side points have local negative psh peaks,
globalized by max-pasting with a negative constant. At a cuspidal point
$(0,\zeta)$, $|\zeta|=1$, there is the global peak
\begin{equation*}
\begin{aligned}
 b_\zeta(z,w)&=\operatorname{Re}(\bar\zeta w)-1,\\
 \rho(z,w)&=\max\{\log\|z\|+\Phi(w),\ |w|^2-1\}.
\end{aligned}
\tag{H.51}\label{eq:appH-Hartogs-peaks}
\end{equation*}
The first function is pluriharmonic, zero only at the selected point
of $\overline\Omega_n$ and negative elsewhere. A strong peak does not
require a positive Levi form. The second is a continuous negative psh
exhaustion with zero boundary values: near the axis it equals $|w|^2-1$,
and at the collapsed circle this term tends to zero. Thus the domain
is hyperconvex. A boundary Jensen measure tested against a continuous
global peak must be supported at that peak point. The B-regularity
criterion in \cite[Corollary~3.8, printed p.~188]{Wikstrom01JensenPublished}
therefore applies. Taking $\Omega_2$ times extra discs would introduce
boundary analytic discs and would not give this B-regular construction.

For the envelope, restrict any admissible $u$ to $H$. Its restriction
$v(w)=u(0,w)$ is subharmonic or identically $-\infty$, globally bounded
above and at most $P$. Away from $1$, its boundary limsup is at most zero.
The circle submean inequality and reverse Fatou, using the global upper
bound, give $v\leq0$; a disk automorphism gives the estimate at every
center. The zero competitor gives $V=0$ on $H$.
The functions $u_\varepsilon$ in \eqref{eq:appH-Hartogs-envelope} are psh,
including across the axis, and satisfy $u_\varepsilon\leq W$.
The defining inequality gives
$\log\|z\|<-\Phi\leq-P^2$, hence the bound follows from
$t-\varepsilon t^2\leq1/(4\varepsilon)$.
They tend increasingly to $W$ off $H$. Density of that complement and
continuity of $W$ give $V^*=W$. Their bounds grow as $j/4$ for
$\varepsilon=1/j$; no common global cap is being imposed.

\subsection{Where a strong majorant gives pointwise recovery}
\label{sec:appH-strong-majorant-recovery}

A positive plurisuperharmonic $G$ is allowed to take $+\infty$, but is
not identically $+\infty$. Call it a strong majorant from above when,
for every $\varepsilon>0$, $W\leq\varepsilon G+M_\varepsilon$ with
a finite constant $M_\varepsilon$. In the example,
\begin{equation*}
 G=-\log\|z\|>\Phi\geq W^2,
 \qquad W\leq\varepsilon G+\frac1{4\varepsilon}.
\tag{H.52}\label{eq:appH-Hartogs-majorant}
\end{equation*}
The pole set of $G$ is exactly $H$. A superlinear scalar transform
with a nontrivial plurisuperharmonic majorant makes $W$ tame in
\cite[Theorem~3.2, actual pp.~4--5]{NilssonWikstrom20QuasiboundedV1};
here the transform is $W^2$.
The strong-majorant equivalence is also
\cite[Lemma~2.1, printed p.~3961]{Nilsson22EnvelopesPublished}.
For this positive $W$, \eqref{eq:appH-Hartogs-majorant} gives the
absolute-value domination in
\cite[Definition~3.1, actual p.~7]{Nilsson25QuasiboundedSolutionsV1}.
Nontriviality of $G$ does not mean finiteness on the axis.

Proposition~3.4 of \cite[actual pp.~7--8]{Wikstrom26BoundedMinorantsV2}
identifies the exact pointwise condition.

\begin{proposition}
\label{prop:appH-strong-majorant-recovery}
Let $\Omega\subset\mathbf C^n$ be a domain and $W\in\PSH(\Omega)$.
A positive strong majorant $G$ gives $V_W^{\PSH}=W$ wherever $G$ is
finite. If $W$ is finite and pluriharmonic, then $V_W^{\PSH}(p)=W(p)$
if and only if some positive strong majorant is finite at $p$.
\end{proposition}

\begin{proof}
For sufficiency, $W-\varepsilon G$ is psh, bounded above by
$M_\varepsilon$ and below $W$. Its values tend to $W(p)$ when
$G(p)<+\infty$.
For the converse, choose bounded-above competitors with deficit at $p$
less than $2^{-k}$, and take finite maxima to make them increasing.
Write $M_k=\sup_\Omega u_k$. Pluriharmonicity makes $W-u_k$
nonnegative plurisuperharmonic. Therefore
\begin{equation*}
 G=1+\sum_{k=1}^\infty(W-u_k),\qquad
 G(p)<2,\qquad NW\leq\sum_{k=1}^NM_k+G.
\tag{H.53}\label{eq:appH-majorant-deficit-sum}
\end{equation*}
The negative partial sums decrease as psh functions and are finite in
the limit at $p$, so $G$ is nontrivial plurisuperharmonic. Dividing
the last inequality by $N$ and choosing $N\geq1/\varepsilon$ gives
the strong-majorant bound, with a finite constant.
\end{proof}

The converse uses pluriharmonicity and is not asserted for arbitrary psh
$W$. In the example, every strong majorant has a pole at every point
of $H$. Restricting $G$ to that disc gives the identically infinite
function, so it supplies no nontrivial disk majorant of the Poisson kernel.
The raw-envelope failure is consistent with tameness. It is also consistent
with continuity of a maximal Dirichlet solution: $V^*=W$ itself is already
continuous and maximal in the interior.

\subsection{A boundary Jensen bound without exchanging limits}
\label{sec:appH-restricted-Jensen-bound}

On a bounded domain let $\mathcal J_p^b(\partial\Omega)$ consist of
probability measures supported on the boundary such that
$u(p)\leq\int u^*\,d\mu$ for every globally bounded-above psh $u$;
the star here includes the upper boundary trace. On a B-regular domain
this class agrees with the boundary-supported continuous-closure Jensen
class by \cite[Corollary~4.3, printed p.~192]{Wikstrom01JensenPublished}.
For an usc $f:\Omega\to[-\infty,+\infty]$, the one-sided estimate in
\cite[Proposition~3.2]{Wikstrom26BoundedMinorantsV2} is
\begin{equation*}
 V_f^{\PSH}(p)\leq
 \lim_{M\to+\infty}\inf_{\mu\in\mathcal J_p^b(\partial\Omega)}
              \int_{\partial\Omega}(\min\{f,M\})^*\,d\mu.
\tag{H.54}\label{eq:appH-restricted-Jensen-bound}
\end{equation*}
For any fixed competitor, $M$ eventually exceeds its global bound,
so its boundary trace is below the truncated obstacle's trace.
The Jensen inequality, followed by the infimum and then the monotone
limit, gives the estimate. An empty measure class has infimum $+\infty$.
Equality of Jensen classes does not justify exchanging a limit and an
infimum over varying measures.

For $p\in H$, harmonic measure on its disk lifts to the collapsed
boundary circle and is a Jensen measure for the bounded family.
The trace of $\min(W,M)$ is zero there except at $(0,1)$, which has
harmonic measure zero. Its integral is zero for every finite $M$.
Thus one fixed measure detects the axis loss at every truncation.
These local-domain statements neither replace compact obstacle/energy
theorems nor provide a non-Archimedean Siu or envelope-closure argument.

\section{Siu analyticity for positive pluriharmonic currents}
\label{sec:appH-pluriharmonic-Siu}
\index{Siu decomposition!positive pluriharmonic currents}
\index{pluriharmonic current!analytic threshold sets}

The currents used for potentials and nonpluripolar products earlier in
this appendix are positive and closed. There is a larger class:
a positive current is \emph{pluriharmonic} when $\ddc T=0$.
For currents this condition is weaker than $dT=0$. In particular,
one cannot write an arbitrary pluriharmonic $(1,1)$-current locally
as $\ddc u$, since such a representation would make it closed.
Dinh--Nguyen--Nguyen's revised paper
\cite{DinhNguyenNguyen26SiuV2} extends Siu's threshold analyticity
to this larger class on compact K\"ahler manifolds.

For a current of bidimension $(q,q)$, the plane-one convention is
\begin{equation*}
 \nu(T,x)=\lim_{r\downarrow0}r^{-2q}
       \int_{\|z\|<r}T\wedge(\ddc\|z\|^2)^q,
 \qquad E_c(T)=\{x:\nu(T,x)\geq c\},\quad c>0.
\tag{H.55}\label{eq:appH-PH-Lelong}
\end{equation*}
Here $z$ is a centered holomorphic coordinate system. With the source's
operator $2\ddc$, its Lelong formula has the additional factor $2^{-q}$,
so that the same integration current has Lelong number one.
For $q=0$, the integral is a measure mass and its limit is the atom
at $x$. These are ordinary trace integrals, not nonpluripolar products.

The following is Theorem~1.2 of
\cite[actual p.~2; proof pp.~11--12]{DinhNguyenNguyen26SiuV2}.
The exact revision and qualifications of its imported inputs are
recorded in \Cref{sec:provenance-pluriharmonic-Siu}.

\begin{theorem}
\label{thm:appH-pluriharmonic-Siu}
Let $X$ be a compact K\"ahler manifold of complex dimension $n$,
and let $T$ be a positive pluriharmonic current of bidimension
$(q,q)$, $0\leq q\leq n$. Every $E_c(T)$ is an analytic subset
of dimension at most $q$. There is a finite or countable decomposition
\begin{equation*}
 T=\sum_{i\in I}\lambda_i[V_i]+T_0,
 \qquad \lambda_i>0,\quad \dim V_i=q,
 \qquad \dim E_c(T_0)\leq q-1\quad(c>0),
\tag{H.56}\label{eq:appH-PH-Siu-decomposition}
\end{equation*}
where the $V_i$ are distinct irreducible compact analytic subsets and
$T_0$ is positive and pluriharmonic. The sum converges in mass and
weakly as a current. A dimension bound of $-1$ means the set is empty.
\end{theorem}

The coefficients are the restrictions of $T$ to its analytic components
and are chosen once. With a K\"ahler form $\omega$, positivity gives
\[
 \sum_i\lambda_i\int_{V_i}\omega^q\leq\int_XT\wedge\omega^q<\infty.
\]
Thus only countably many nonzero component masses occur, and finite
partial remainders are positive. This fixed-component decomposition
does not rely on a sequence of changing coefficients.
For $q=0$ it is the atomic/diffuse decomposition of a positive measure;
for $q=n$, compact pluriharmonic functions are constant on each
connected component. These cases can be treated directly.

\subsection{Total transforms and a closed current dominating Lelong numbers}
\label{sec:appH-PH-transforms}

The proof first reduces to bidegree $(1,1)$ by the pluriharmonic
Lelong--Skoda transform, preserving the Lelong function. The compact
pluriharmonic extension is Theorem~5.1 of
\cite[actual p.~12]{Vigny19LelongSkodaV2}; its closed-current
Theorem~3.1 alone would not suffice.

For a modification $\pi:\widetilde X\to X$, distinguish the total
transform $\pi^*T$ from the strict transform $\pi^\diamond T$,
defined by extension of the pullback away from the exceptional locus.
For a positive pluriharmonic $(1,1)$-current, the compact setting gives
\begin{equation*}
 \pi^*T=\pi^\diamond T+\sum_j a_j[D_j],\qquad a_j\geq0,
 \qquad \pi_*\pi^*T=T.
\tag{H.57}\label{eq:appH-PH-total-strict}
\end{equation*}
The total transform has a unique Aeppli-class characterization.
The exceptional coefficient functions in a general modification are
weakly psh; compactness forces them to be constant, so the strict
transform is again pluriharmonic. The precise source is
\cite[Theorems~3.9 and~3.11, printed pp.~725--728]{AlessandriniBassanelli03Transforms}.
Without compact exceptional components, the strict transform can be
only plurisuperharmonic.

The compact analytic input can also be anchored directly in
\cite[Theorem~4.1 and Corollary~4.2, printed pp.~965--966]{DS04}:
bounded-mass differences of smooth positive pluriharmonic currents
and finite-mass pluriharmonic strict transforms. The approximation
is a \emph{difference}. It does not assert that $T$ is a weak limit
of individually positive smooth pluriharmonic representatives.
The continuous-potential wedge used in the transform has the separate
hypotheses of \cite[Theorem~5.1, printed pp.~967--969]{DS04}.

Lemmas~3.1--3.2 of \cite{DinhNguyenNguyen26SiuV2} then produce
a positive \emph{closed} $(1,1)$-current $S$ such that
\begin{equation*}
 \{S\}=\{T\},\qquad \nu(S,x)\geq\nu(T,x)\quad\text{for all }x\in X.
\tag{H.58}\label{eq:appH-PH-closed-domination}
\end{equation*}
The braces identify the real Aeppli class with its corresponding
$(1,1)$ cohomology class through the K\"ahler $\ddc$ lemma.
This identification does not make $T$ closed. Nor is the Lelong
inequality a current inequality $S\geq T$.

The cone input combines the balanced/pluriharmonic duality of
\cite[Lemma~3.3 and Remark~3.4, actual pp.~11--12]{FuXiao14ConesV2}
with the compact-K\"ahler cone equality following
\cite[Corollaries~1.2 and~1.4]{Tosatti26MorseV2}.
To impose Lelong domination at a finite set of points, blow up all
of them and use the independent exceptional classes to determine
their coefficients. Positive closed representatives in one fixed
class have bounded mass and form a weakly compact set. The finite
intersection property then enforces the constraints at every point
simultaneously. In dimension one $T$ itself is closed, and this
point-blow-up construction is unnecessary.

\subsection{The full-fibre step in the analyticity induction}
\label{sec:appH-PH-full-fibre}

Here is the proof mechanism of the attributed theorem. The tangent,
support and resolution inputs retain their source hypotheses; this
description is not a new independent certificate for their entire
recursive proof chains.

By \eqref{eq:appH-PH-closed-domination}, $E_c(T)$ is contained in
the proper analytic set $E_c(S)$. Let $\Sigma$ be its smallest analytic
hull. If an irreducible component $A$ of $\Sigma$ is not contained
in $E_c(T)$, choose $x_0\in A_{\rm reg}$ with $\nu(T,x_0)<c$.
Resolve $A$ by a modification, preserving the regular locus, and
write $A'$ for its smooth strict transform and $T'$ for the total
transform. Blow up $A'$ and denote the smooth exceptional divisor
by $D$. If $A'$ is already a divisor, use the identity map and $D=A'$.

The divisor coefficient and the no-divisor-mass strict transform obey
\begin{equation*}
\begin{aligned}
 \pi^*T'&=\pi^\diamond T'+a[D],\\
 a&\leq\nu(T',x')\leq a+\nu(\pi^\diamond T',y)
                  \leq a+\nu(R,y),\quad y\in\pi^{-1}(x'),
\end{aligned}
\tag{H.59}\label{eq:appH-PH-divisor-tangent}
\end{equation*}
where $R$ is a positive pluriharmonic $(1,1)$-current on $D$.
The last inequality comes from Lemma~4.6 of the controlling paper:
the strict transform's tangent along $D$ has minimum horizontal
dimension and is the pullback of $R$; Lelong semicontinuity gives
the bound. Horizontal dimension means the \emph{largest} exponent
with a nonzero horizontal wedge. Whole normal fibres are required
for the positive pluriharmonic Liouville argument.
The qualitative compact tangent inputs are stated in
\cite[Theorem~1.18 and Corollary~1.21]{Nguyen22TangentsV2} and
\cite[Theorem~2.14(2)]{Nguyen25IntersectionV1}.

At the lift of $x_0$, \eqref{eq:appH-PH-divisor-tangent} gives $a<c$.
For every $x'$ over $E_c(T)\cap A_{\rm reg}$, \emph{all} points of
the exceptional fibre lie in $Y=E_{c-a}(R)$. By induction on dimension,
$Y$ is a proper analytic subset of $D$. If the projective fibres have
dimension $r$, use the locus
\begin{equation*}
 F=\{x'\in A':D_{x'}\subset Y\}
   =\{x'\in A':\dim(Y\cap D_{x'})=r\}.
\tag{H.60}\label{eq:appH-PH-full-fibre-locus}
\end{equation*}
Proper holomorphic fibre-dimension theory makes $F$ analytic.
It is proper because $Y\ne D$. It contains the lifted threshold
set, contradicting its analytic density in the chosen component.
When $r=0$ the same locus is simply $Y$ on the divisor.
Taking the image of an arbitrary proper subset of $D$ would not
give this conclusion: a section can surject onto the base.
The full-fibre inclusion is the decisive extra information.

This proves the analyticity step in the source route. Analytic
restriction and the support theorem then give the fixed coefficients
in \eqref{eq:appH-PH-Siu-decomposition}. Compact K\"ahler geometry,
positive currents and complex blow-ups are used throughout; these
inputs cannot be replaced formally by non-Archimedean analogies.

\subsection{Cone applications and weak-limit qualifications}
\label{sec:appH-PH-cone-boundary}

Theorem~1.5 of \cite{DinhNguyenNguyen26SiuV2} also states a
movable-class conclusion for bidimension-one currents with no mass
on proper analytic subsets. Its proof has an additional input:
the mass/cup-product identity cited as Theorem~1.9 of Nguyen's
\emph{Cohomology of density and intersection of currents}, listed
as to appear in 2026. No authentic source file for that theorem was
obtained in this audit. That application is not independently derived
here, and the main analyticity induction above does not use it.

The mixed tensor argument in that application pairs a pluriharmonic
current with a \emph{closed} current. For two nonclosed pluriharmonic
currents on a surface, the original energy formula instead has cross
terms; see \cite[Theorem~2.2, actual pp.~8--9]{DinhNguyenSibony19SurfaceV2}.
For a self-pair the tangent mass is $\{T\}^2-2E(T)$.
Closedness makes the energy vanish, as in that source's Lemma~B.2.
Dropping it silently changes the intersection formula.

The Morse inequality behind the cone input is an attributed result
of \cite[Theorem~1.1]{Tosatti26MorseV2}: for nef real classes
$\alpha,\beta$ on compact K\"ahler $X$,
\begin{equation*}
 \operatorname{vol}(\alpha-\beta)
       \geq\int_X\alpha^n-n\int_X\alpha^{n-1}\beta.
\tag{H.61}\label{eq:appH-Tosatti-Morse}
\end{equation*}
Its fixed-density weak-limit step must be distinguished from a
stronger pointwise limsup assertion. In local coordinates, if positive
$(1,1)$-currents $T_j$ converge weakly to $T$ and all their absolutely
continuous matrix densities have determinant at least the same
nonnegative locally integrable $f$, determinant-root concavity gives
\[
 \det(T*\rho_\varepsilon)^{1/n}
        \geq f^{1/n}*\rho_\varepsilon.
\]
Indeed the singular positive part can only increase the smoothed
matrix, and convolution commutes with the weak limit for fixed
$\varepsilon$. Lebesgue differentiation as $\varepsilon\downarrow0$
gives $\det T_{\rm ac}\geq f$ almost everywhere.
This is the common-lower-bound step used in Tosatti's Proposition~2.3.

It does not give $\det T_{\rm ac}\geq\limsup_j\det(T_j)_{\rm ac}$.
The latter is printed in \cite[Proposition~2.1, printed p.~1048]{Bou02b},
but fails literally. On a flat elliptic curve with a fixed area form
$\omega$, the positive closed currents
$T_j=(1+\sin(2\pi jx))\omega$ converge weakly to $\omega$.
For almost every $x$, irrational rotation gives density limsup $2$,
whereas the limiting density is $1$. A common fixed lower bound
is preserved nevertheless. This auxiliary distinction neither refutes
\eqref{eq:appH-Tosatti-Morse} nor supplies a representative-level
positive subtraction. Further numerical and edition boundaries are
listed in \Cref{sec:provenance-pluriharmonic-Siu}.


\section{Analytic multiplier ideals on affine toric pairs}
\label{sec:appH-toric-multiplier-pairs}
\index{multiplier ideal!affine toric pair}
\index{Newton body!analytic toric weight}

The smooth monomial calculation in \Cref{prop:ch51-howald} has a
singular affine counterpart. Here the canonical divisor changes the
translation of the Newton body, and the character lattice determines
which exponents exist. Seo's \cite[Theorem~6.1, actual pp.~28--29]{Seo26ToricMultiplierV1}
gives the criterion below for arbitrary toric psh weights, including
cones with a torus factor. This section explains that imported result
and its source qualifications. Its ambient space is a normal affine
complex analytic toric space, rather than the compact K\"ahler space
used in the earlier global formulas of this appendix.

\subsection{The pair, its volume and the weight convention}
\label{sec:appH-toric-pair-volume}

Let $N$ be a lattice, $M=\operatorname{Hom}(N,\mathbf Z)$ its actual
dual lattice, and $\sigma\subset N_{\mathbf R}$ a strongly convex
rational polyhedral cone. Put
$X=X_\sigma=(\operatorname{Spec}\mathbf C[\sigma^\vee\cap M])^{\rm an}$.
The cone need not span $N_{\mathbf R}$. For its primitive ray vectors
$v_\rho$, write $D_\rho$ for the invariant prime divisors. Fix a
rational invariant Weil divisor $D=\sum_\rho a_\rho D_\rho$ such that
$K_X+D$ is $\mathbf Q$-Cartier, using $K_X=-\sum_\rho D_\rho$.
There is $u_D\in M_{\mathbf Q}$ satisfying
\begin{equation*}
\begin{aligned}
 \langle u_D,v_\rho\rangle&=1-a_\rho,\\
 q(K_X+D)&=\operatorname{div}(\chi^{-q u_D}),\qquad q u_D\in M.
\end{aligned}
\tag{H.62}\label{eq:appH-toric-canonical-shift}
\end{equation*}
The integer $q>0$ must clear the Cartier index in $M$. Integral
coordinates in a larger covering lattice are insufficient. When
$\sigma$ is not full-dimensional, $u_D$ is determined only modulo
the annihilator of $\sigma$; the criterion below is independent of
that choice.

The weight $\phi$ is invariant under the compact torus $K_N$, is psh
in the local ambient-extension sense on $X$, and is not identically
$-\infty$. The convention in this source is
$\phi_{\mathfrak a}=\tfrac12\log\sum_j|f_j|^2+O(1)$.
Thus the squared-norm potential $\psi$ used in Chapter~51 is $2\phi$.
A local generator $\beta$ of $\omega_X^{[q]}(qD)$ defines an adapted
volume on the regular locus. In local smooth coordinates, if
$\beta=b\,(dz_1\wedge\cdots\wedge dz_n)^{\otimes q}$, its density
is $|b|^{2/q}$ times the usual positive coordinate volume.
Changing the generator multiplies this density by a locally bounded
positive unit. Definition~3.1 of \cite{Seo26ToricMultiplierV1} uses
\begin{equation*}
 \mathcal I(X,D;\phi)_x
 =\{f\in\mathcal M_{X,x}:
       |f|^2\exp(-2\phi)\text{ is locally }dV_\beta\text{-integrable}\}.
\tag{H.63}\label{eq:appH-toric-adapted-module}
\end{equation*}
Here $\mathcal M_X$ denotes meromorphic germs. Integrability is tested
on the regular locus away from the divisorial poles of the density.
Li's \cite[Definition~2.2 and Remark~2.3, actual p.~3]{Li22KollarArchivedV1}
give the checked adapted-volume antecedent; the integer generator
used here retains the $\mathbf Q$-Cartier hypothesis.

The module is coherent, by Seo's Corollary~3.8, and is contained in
$\mathcal O_X$ when $D$ is effective, by Proposition~3.3. Negative
boundary coefficients can allow poles. On $\mathbf C$, with
$D=-[0]$ and $\phi=0$, the density is $|z|^2$ times area, so $1/z$
is integrable and $\mathcal I=\mathcal O([0])$. Calling the general
module an ordinary holomorphic ideal would exclude this elementary
case.

\subsection{The logarithmic Newton domain}
\label{sec:appH-toric-Newton-domain}

On the dense torus define $\operatorname{Log}$ by
$\langle m,\operatorname{Log}t\rangle=\log|\chi^m(t)|$.
Seo's Lemma~5.1 writes $\phi=g\circ\operatorname{Log}$, where
$g:N_{\mathbf R}\to\mathbf R$ is finite, convex and increasing in
cone directions. These are global weights on $X$; a local psh germ
is not assumed to have a global extension. With $\delta_{-\sigma}$
the convex indicator, the Newton domain of Definition~5.2 is
\begin{equation*}
\begin{aligned}
 P(\phi)&=\operatorname{dom}(g+\delta_{-\sigma})^*\\
 &=\left\{p\in M_{\mathbf R}:
       \sup_{x\in-\sigma}\bigl(\langle p,x\rangle-g(x)\bigr)<+\infty\right\}.
\end{aligned}
\tag{H.64}\label{eq:appH-toric-Newton-domain}
\end{equation*}
This convex domain need not be closed. Its closure and its ordinary
interior in $M_{\mathbf R}$ have different threshold tests. If
$g^\infty$ is the recession function of $g$, Seo's Lemma~4.4 and
Proposition~5.3 give
\begin{equation*}
\begin{aligned}
 \overline{P(\phi)}
   &=\{p:\langle p,x\rangle\leq g^\infty(x)
                              \text{ for all }x\in-\sigma\},\\
 \operatorname{int}P(\phi)
   &=\{p:\langle p,x\rangle<g^\infty(x)
                              \text{ for all }0\ne x\in-\sigma\}.
\end{aligned}
\tag{H.65}\label{eq:appH-toric-recession-tests}
\end{equation*}
An infinite recession value makes the corresponding inequality
automatic. Cone monotonicity implies $P(\phi)\subset\sigma^\vee$.
For the zero cone, $P(\phi)=M_{\mathbf R}$ and the inequalities in
the interior test are vacuous.

The needed convex facts have checked primary anchors in
\cite[Theorems~3.21, 3.26(a), 11.1, 11.5 and 11.23(a)]{RockafellarWets09Author}:
recession slopes, uniform linear growth, biconjugacy, the support
function of a conjugate domain, and conjugation of a sum. Since $g$
is finite on the whole vector space, the domain-interior qualification
for the sum holds even if $\sigma$ has smaller dimension. These
anchors are an alternative to the 1970 Rockafellar locators cited
by Seo; they do not certify the unread controlling 1970 edition.

\subsection{The character criterion}
\label{sec:appH-toric-character-criterion}

The following statement is Seo's imported affine-pair result, in the
normalization \eqref{eq:appH-toric-adapted-module}.

\begin{theorem}
\label{thm:appH-toric-analytic-multiplier}
Under the pair and weight hypotheses of
\Cref{sec:appH-toric-pair-volume}, the coherent fractional module
$\mathcal I(X,D;\phi)$ is globally generated by finitely many
meromorphic characters. The characters belonging to it are exactly
\begin{equation*}
 \chi^m\in\mathcal I(X,D;\phi)
 \quad\Longleftrightarrow\quad
 m\in M\ \text{ and }\ m+u_D\in\operatorname{int}P(\phi).
\tag{H.66}\label{eq:appH-toric-character-criterion}
\end{equation*}
Equivalently, it is the $\mathcal O_X$-module generated by these
characters, and a finite subset suffices. When $D\geq0$, this module
is a coherent ideal in $\mathcal O_X$.
\end{theorem}

For $X=\mathbf C^n$ and $D=0$, the actual lattice is $\mathbf Z^n$
and $u_D=(1,\ldots,1)$. The criterion becomes the familiar strict
translation by the all-ones vector. Guenancia's archived v2
\cite[Theorem~1.13, actual pp.~7--8]{Guenancia11ToricArchivedV2}
supplies the smooth analytic antecedent. A local smooth comparison
uses a polydisk and its logarithmic tails, rather than asserting a
global extension of every germ. For algebraic ideal weights, Seo's
Theorem~3.2 identifies the analytic module with the analytification
of the pair multiplier module. The archived algebraic antecedent
\cite{Blickle03ToricArchivedV1} needs the sign qualification in
\Cref{sec:provenance-toric-analytic-multiplier} before its formulas
can be compared with \eqref{eq:appH-toric-canonical-shift}.

\subsection{Coherence and finite monomial generation}
\label{sec:appH-toric-coherence-generation}

The source proof first establishes the analytic sheaf properties.
For a log resolution $\mu:Y\to X$, put
$A=K_Y-q^{-1}\mu^*(q(K_X+D))$. Lemma~3.4 and Proposition~3.5 give
the birational transformation
\begin{equation*}
\begin{aligned}
 \mathcal I(X,D;\phi)
   =\mu_*\bigl(&\mathcal O_Y(\lceil A\rceil)\otimes\\
              &\mathcal I(Y,0;\phi\circ\mu+
                                  \phi_{\lceil A\rceil-A})\bigr).
\end{aligned}
\tag{H.67}\label{eq:appH-toric-birational-module}
\end{equation*}
The remainder $\lceil A\rceil-A$ is effective with coefficients in
$[0,1)$; the weight of a divisor of coefficient $b$ is $b\log|z|$
in this normalization. Pullback volume powers $|z|^{2a}$ fix the
rounding sign. Integer poles are absorbed into
$\mathcal O_Y(\lceil A\rceil)$. For analytic ideal weights, the
analogous rounding is $\lceil a_i-cb_i\rceil$, leaving integrable
powers with exponents in $(-1,0]$. Meromorphic descent and proper
direct-image compatibility are part of the imported resolution route.

Local denominator clearing and Cartier twisting reduce the general
boundary to an effective one. Smooth Nadel coherence and Grauert's
theorem then give coherence. A finite covering of a compact
resolution fibre must be taken before the base neighbourhood is
shrunk. Smooth strong openness, applied to $2\phi$, also gives
Seo's Corollary~3.9: locally,
$\mathcal I(X,D;\phi)=\mathcal I(X,D;(1+\varepsilon)\phi)$
for all sufficiently small $\varepsilon>0$, with $D$ fixed.
The controlling smooth input is
\cite[Theorem~1.1, printed p.~606]{GuanZhou15AnnalsPublished}.
Scaling the boundary as well would be a different assertion.

Sections~5.3--5.4 of \cite{Seo26ToricMultiplierV1} establish finite
generation on the affine analytic space, which is Stein. Compact-torus
invariance first becomes full-torus invariance by holomorphic
one-parameter families and the identity theorem in the quotient of
global section spaces. Fourier projection of a \emph{global} section
$f$ has the form
$\int_{K_N}\chi^{-m}(k)f(kx)\,dk$. On the dense torus a nonzero
projection is a constant multiple of $\chi^m$. Coherence makes these
projections compatible with the quotient module. One does not average
an individual germ while moving its base point.

In the finite jet quotients at the closed orbit, vanishing of all
Fourier coefficients forces the class to vanish. Krull intersection
then recovers the germ from all jets. Cartan A supplies enough global
sections, and Noetherianity yields finitely many monomial generators.
Torus translates of a closed-orbit neighbourhood cover $X$, so the
generators extend to the whole affine space. For a fractional module,
the invariant coherent colon module $(\mathcal O_X:\mathcal I)$
supplies a global character denominator, reducing to the ideal case.
These are the proof mechanisms of Theorem~5.5, Lemma~5.6,
Corollary~5.8 and Theorem~5.9 of the source. Cartan A, Krull
intersection, smooth coherence, resolution and normal psh extension
remain attributed foundations, rather than new proofs in this book.

\subsection{Logarithmic tails and a torus factor}
\label{sec:appH-toric-local-tails}

For a full-dimensional cone in its integration vector space, the
radial character integral has the form
\begin{equation*}
 \int_{-\sigma}
       \exp\bigl(2\langle m+u_D,x\rangle-2g(x)\bigr)\,dx.
\tag{H.68}\label{eq:appH-toric-exponential-integral}
\end{equation*}
The compact angular integral contributes only a positive constant.
Seo's Theorem~4.6 makes finiteness equivalent to the strict recession
test in \eqref{eq:appH-toric-recession-tests}. For sufficiency,
the convex function $g-\langle m+u_D,\cdot\rangle+
\delta_{-\sigma}$ has strictly positive recession in every nonzero
direction. The uniform linear lower bound of
\cite[Theorem~3.26(a)]{RockafellarWets09Author} gives exponential
decay with a common positive rate. For necessity, a failed direction
produces a half-tube of positive volume, not just a measure-zero ray.
Remark~4.7 retains uniformity over compact transverse parameters.

For a non-full cone, $N_\sigma=N\cap\operatorname{span}_{\mathbf R}\sigma$
is saturated. Choose a lattice complement $N'$; then
$X_\sigma=X_{\sigma,N_\sigma}\times T_{N'}$. The unique closed
orbit is $\{z_0\}\times T_{N'}$, generally of positive dimension.
The cone integral is taken in $N_{\sigma,\mathbf R}$, while transverse
torus coordinates range over compact annuli. Fubini and the source's
uniform tail estimate give the same character criterion. An element
annihilating $\sigma$ changes neither Newton membership nor this
local integrability condition. If $\sigma=\{0\}$, the space is a
torus, the weight is locally bounded, and the module is $\mathcal O_X$.

One literal neighbourhood equality in the proof of Theorem~6.1
requires qualification. The logarithm of an open monomial neighbourhood
is not exactly a closed translated cone. In dimension one, the
logarithm of an open disk omits the translated ray's endpoint.
The comparison needed for local integrability uses cofinal inclusions.
In the full cone span choose nonzero semigroup generators $m_j$ and
put $S_L=\{x:\langle m_j,x\rangle<-L\text{ for every }j\}$.
For $b\in\operatorname{int}\sigma$ with every
$\langle m_j,b\rangle>L$, and for fixed
$b_0\in\operatorname{int}\sigma$ followed by
$L>\max_j\langle m_j,b_0\rangle$, respectively,
\begin{equation*}
 -b-\sigma\subset S_L,
 \qquad S_L\subset-b_0-\sigma.
\tag{H.69}\label{eq:appH-toric-cofinal-tails}
\end{equation*}
The two inclusions use different choices of translation. They give
the local comparison with logarithmic monomial polydisks; translation
preserves recession slopes. Compact transverse annuli are retained
in the product case. This elementary explanation qualifies the
printed equality; it is not a separate original repair theorem.

\subsection{A singular surface calculation in its actual lattice}
\label{sec:appH-toric-A1-calculation}

Seo's Example~6.2 uses $X=\{z_1z_2=z_3^2\}$, with covering
coordinates $z_1=s^2$, $z_2=t^2$, $z_3=st$. Thus
$M=\{(a,b)\in\mathbf Z^2:a+b\text{ even}\}$ and
$N=\mathbf Z^2+\mathbf Z(\tfrac12,\tfrac12)$, with the positive
quadrant cone. Let $E=\{z_1=z_3=0\}$, $D=\tfrac12E$ and
$\phi=\max\{\sqrt2\log|z_1|,2\log|z_2|\}$. Then
$u_D=(\tfrac12,1)$, but $2u_D=(1,2)\notin M$.
The minimal index is $4$, and
$4(K_X+D)=-\operatorname{div}(z_1z_2^2)$.
The logarithmic function is
$g(x_1,x_2)=\max\{2\sqrt2x_1,4x_2\}$, with
$P(\phi)=\operatorname{conv}\{(2\sqrt2,0),(0,4)\}
+\mathbf R_{\geq0}^2$. The translated strict test becomes
\begin{equation*}
\begin{aligned}
 &(a,b)\in\mathbf Z_{\geq0}^2,\quad a+b\text{ even},\qquad
   \frac{a+1/2}{2\sqrt2}+\frac{b+1}{4}>1,\\
 &\mathcal I(X,\tfrac12E;\phi)=(z_1,z_3,z_2^2),\qquad
   \mathcal I(X,0;\phi)=(z_1,z_2,z_3).
\end{aligned}
\tag{H.70}\label{eq:appH-toric-A1-ideals}
\end{equation*}
For the first ideal the minimal exponents are $(2,0)$, $(1,1)$
and $(0,4)$. If $a\geq2$, subtract $(2,0)$; if $a=1$, parity
forces odd $b\geq1$, so subtract $(1,1)$. If $a=0$, the even
values $b=0,2$ fail the strict test and every even $b\geq4$ is
divisible by $(0,4)$. This accounts for all allowed exponents,
not merely a bounded search. For $D=0$, the shift is $(1,1)$;
$(0,0)$ fails, whereas $(2,0)$, $(1,1)$ and $(0,2)$ pass and
generate every remaining allowed semigroup point. The irrational
coefficient is handled by exact inequalities in $\sqrt2$. Neither
rounding an exponent nor replacing $M$ by $\mathbf Z^2$ is legitimate.

\subsection{Subtracting an ideal weight and keeping thresholds distinct}
\label{sec:appH-toric-negative-ideal-weight}

Let a nonzero invariant coherent ideal have finite character generators
$\chi^{a_j}$, and put
$\phi_{\mathfrak a}=\tfrac12\log\sum_j|\chi^{a_j}|^2$.
For positive \emph{rational} $c$, Seo's Proposition~5.10 and
Corollary~6.3 allow the weight $\phi-c\phi_{\mathfrak a}$, which
need not be psh. Finite sum comparability gives
\begin{equation*}
\begin{aligned}
 \mathcal I(X,D;\phi-c\phi_{\mathfrak a})
   &=\bigcap_j\mathcal I(X,D-c\operatorname{div}(\chi^{a_j});\phi),\\
 \chi^m\in\mathcal I(X,D;\phi-c\phi_{\mathfrak a})
   &\Longleftrightarrow
       m+u_D+ca_j\in\operatorname{int}P(\phi)\quad\text{for every }j.
\end{aligned}
\tag{H.71}\label{eq:appH-toric-negative-ideal-criterion}
\end{equation*}
Rationality keeps the shifted pair $\mathbf Q$-Cartier; the displayed
pair proof does not establish an irrational extension. Strong openness
in this formulation scales $\phi$ only, with $D$, $c$ and the ideal
fixed. A finite intersection permits a common small perturbation.

Ordinary multiplier and maximal non-lc ideals have different boundary
conventions. With effective $D$, maximal non-lc rounding moves negative
integer discrepancies off their ordinary multiplier threshold. The
checked antecedent is
\cite[Definition~7.1 and Lemmas~7.2, 7.5, actual pp.~10--12]{FujinoSchwedeTakagiAuthorRevision}.
On $\mathbf C$, $D=[0]$ and $\phi=0$ give ordinary multiplier ideal
$(z)$ but maximal non-lc ideal $\mathcal O$. The closed Newton
membership used for the latter cannot replace the strict interior
in \eqref{eq:appH-toric-character-criterion}.

Finally, a locally bounded perturbation of the weight changes neither
the integrability module nor the relevant finite-supremum Newton domain
on relatively compact logarithmic cone neighbourhoods. Seo's
Remark~6.4 allows psh plus a locally bounded measurable function;
that class is broader than the psh-plus-smooth convention for
quasi-psh functions in this book. It is stated explicitly here without
changing the earlier definition. These local convex tests concern
multiplier integrability. They supply neither a compact partial
Okounkov-body mass formula at zero volume nor a general envelope
closure argument. Exact version boundaries are recorded in
\Cref{sec:provenance-toric-analytic-multiplier}.

\section{Canonical volume forms and normalization in families}
\label{sec:appH-KE-family-normalization}
\index{Kahler--Einstein equation@K\"ahler--Einstein equation!family normalization}
\index{plurigenera!family variation and proof boundaries}

A fixed-fibre Monge--Amp\`ere equation and a positively curved metric
on a total space impose different normalization requirements. This
section explains that distinction in the setting of
\cite{LPWZ26KEV1}. That preprint announces invariance of plurigenera
for smooth K\"ahler families in every dimension. Its independent
reference-volume statement and several proof interfaces require the
qualifications below; its announced conclusion is not used here as
a proved input. The background existence and estimate theorems
are attributed separately to their controlling sources.

\subsection{The constant in an exponential equation}
\label{sec:appH-KE-exponential-constant}

First let $Y$ be a connected compact K\"ahler manifold, $\theta$
a smooth representative of a big class, and $dV$ a smooth positive
volume form. The imported existence theorem
\cite[Theorem~6.1, actual pp.~47--48]{BEGZ10BigV2} gives a unique
full-mass solution, with minimal singularities, of
\begin{equation*}
 \MA_\theta(\varphi)=e^\varphi dV,
 \qquad \int_Y e^\varphi dV=\vol(\{\theta\})=:v>0.
\tag{H.72}\label{eq:appH-KE-exponential-equation}
\end{equation*}
The equation is written in this appendix's $\ddc$ convention;
a positive change of curvature normalization rescales both the
background and the volume form before the source theorem is applied.
The final solution has no independently prescribed supremum. By
contrast, a nonexponential equation is unchanged by adding constants,
so uniqueness requires an auxiliary normalization.

Put $W=\int_YdV$, $d\nu=W^{-1}dV$ and
$u=\varphi-\log(v/W)$. Direct substitution gives
\begin{equation*}
 v^{-1}\MA_\theta(u)=e^u d\nu,
 \qquad \int_Ye^u d\nu=1.
\tag{H.73}\label{eq:appH-KE-probability-equation}
\end{equation*}
This is normalization by a positive mass. It is not defined at $v=0$.
Replacing $dV$ by $c\,dV$ for a positive constant $c$ changes the
exponential solution to $\varphi-\log c$ and leaves its current
unchanged. Replacing $dV$ by a nonconstant density requires more care.

\subsection{Changing a canonical reference volume}
\label{sec:appH-KE-reference-volume}

Now let $p:\mathcal X\to\Delta$ be a proper holomorphic submersion
with connected compact fibres of dimension $n$, and with K\"ahler
total space. The total space need not be compact. A smooth positive
relative volume form $dV_t$ induces a metric on $K_{\mathcal X/\Delta}$.
Locally write $dV_t=a(z,t)\,d\lambda_z$ in fibre coordinates.
The squared norm of the canonical frame is proportional to $a^{-1}$,
so its weight is $\log a$, and its curvature is
$\theta_{dV}=\ddc\log a$ on the total space. The coordinate
transition makes this a global curvature form.

For a smooth positive function $r$ on $\mathcal X$, set
$dV'_t=r|_{\mathcal X_t}dV_t$. With a fixed smooth closed twist
$\beta$, the simultaneous change of curvature and potential is
\begin{equation*}
\begin{aligned}
 \theta_{dV'}&=\theta_{dV}+\ddc\log r,
 &\varphi'&=\varphi-\log r,\\
 \theta_{dV'}+\beta+\ddc\varphi'
   &=\theta_{dV}+\beta+\ddc\varphi,
 &e^{\varphi'_t}dV'_t&=e^{\varphi_t}dV_t.
\end{aligned}
\tag{H.74}\label{eq:appH-KE-volume-covariance}
\end{equation*}
These are identities for the same current and measure. They do not
assert a new variation theorem. If only the density is changed while
the reference curvature is fixed, this cancellation is unavailable.

This is a substantive gate in \cite[Theorem~1.3, actual p.~5]{LPWZ26KEV1}.
Its displayed statement sets $\alpha=\theta_\omega+\beta$ but
uses $e^{\varphi_t}\beta_t^n$; the proof on actual~33--34 uses
$J=\log(\mu^*\omega_t^n/\eta_t^n)$ and proves the equation with
$e^{\varphi_t}\omega_t^n$. Independent $\omega$ and $\beta$
do not give the same exponential solution or the same total-space
positivity assertion by a constant shift. For the special choice
$\beta=s\omega$ with constant $s>0$, the factor $s^n$ is constant
and can be absorbed into $\varphi_t-n\log s$. Thus that particular
discrepancy does not itself refute the positive-twist specialization.
No general corrected variation theorem is asserted here.

\subsection{Uniform fibre bounds and one total-space constant}
\label{sec:appH-KE-uniform-fibre-bound}

The uniform mean/sup comparison of
\cite[Proposition~3.3, actual pp.~26--29]{DGG26FamiliesCorrectedV3}
applies to a smooth submersion in its local-chart setting, after
shrinking to a relatively compact base disk. It does not assert
that all fibres are globally isomorphic. With a fixed dominating
K\"ahler background and probability volumes $d\nu_t$, the relevant
comparison has the form
\begin{equation*}
 M_t-C\leq\int_{\mathcal X_t}u_t\,d\nu_t\leq M_t,
 \qquad M_t=\sup_{\mathcal X_t}u_t.
\tag{H.75}\label{eq:appH-KE-mean-sup-comparison}
\end{equation*}
Assume also $\int e^{u_t}d\nu_t=1$. Then $M_t\geq0$, while
Jensen's inequality gives
$0\geq\int u_t\,d\nu_t\geq M_t-C$. Hence
\begin{equation*}
                        0\leq M_t\leq C.
\tag{H.76}\label{eq:appH-KE-supremum-bound}
\end{equation*}
The same elementary calculation is uniform in an additional twist
parameter whenever the dominating background and comparison constant
are uniform. It does not require a positive lower bound on the
varying class volume, because that volume has already been divided
out at each positive parameter.

Suppose the fibre class volumes $v_s>0$ are known to be constant
in $t$, and a closed relative K\"ahler volume has constant total
mass $W$. Then subtracting $\log(v_s/W)$ from a total-space
potential preserves its curvature for each fixed $s$. The further
fibrewise normalization $\widehat u_{s,t}=u_{s,t}-M_{s,t}$ is useful
for an estimate on that fibre. It does not establish positivity of
$u_s-M_{s,p(x)}$ on $\mathcal X$; a function of the base has its
own horizontal curvature.

The DGG comparison retains its irreducible-fibre and family hypotheses.
Its general singular-family Conjecture~3.1 is distinct from the smooth
case used here. The published corrigendum
\cite{DGG26CorrigendumPublished} corrects specific singular KE lower
bounds; it does not withdraw all the smooth-family mean or exponential
integrability estimates.

\subsection{What fixed-background stability controls}
\label{sec:appH-KE-stability-scope}

The exact one-sided stability theorem is
\cite[Theorem~2.6, printed pp.~1195--1197]{GuedjLu25EnvelopeEstimatesPublished}.
On a compact K\"ahler manifold with a smooth semipositive big
background $P$, let $V=\int P^n>0$ and let
$\phi\in\PSH(Y,P)\cap L^\infty$ satisfy $\sup\phi=0$ and
$V^{-1}\MA_P(\phi)=\mu$, where $\mu$ is a probability measure.
Assume every sup-zero $P$-psh function has a uniformly bounded
$L^m(\mu)$ norm for some $m>n$. For a bounded $P$-psh comparator
$\psi$, the source gives
\begin{equation*}
 \sup_Y(\psi-\phi)_+
 \leq T\left(\int_Y(\psi-\phi)_+\,d\mu\right)^\tau,
 \qquad \tau=\tau(n,m)>0.
\tag{H.77}\label{eq:appH-KE-one-sided-stability}
\end{equation*}
The constant retains the source's moment and comparator bounds.
The measure in this formula is the normalized MA measure of $\phi$.
It cannot be replaced without justification by an arbitrary reference
Lebesgue measure. The theorem is not an unrestricted symmetric
stability assertion. The preceding Theorem~2.1 supplies its a priori
bound by envelopes; contact and domination principles remain inputs.

For a fixed bimeromorphic model, comparison with one K\"ahler
background can provide the required uniform moments for small
positive regularizations. This is a fixed-model argument. Sending
the outer canonical twist to zero may change the model, the volume
and the density; the same constant is not thereby established for
that different limit.

\subsection{A conditional canonical-metric route}
\label{sec:appH-KE-canonical-route}

The Section~8 construction in \cite{LPWZ26KEV1} can be read with
its prerequisites displayed. First require fibrewise pseudoeffectivity
of the canonical classes. Then $\theta_t+s\omega_t$ is big for
all $s>0$, where $\theta=\theta_\omega$. Require, in addition,
positively curved total-space solutions of the corresponding
exponential equations, constancy of their fibre class volumes,
and the fixed-fibre envelope lower bounds uniform in $s$.
With the normalization above, the required lower bound is
\begin{equation*}
 u_{s,t}\geq V_{\theta_t+s\omega_t}-C_t
              \geq V_{\theta_t}-C_t,
 \qquad C_t\text{ independent of }s.
\tag{H.78}\label{eq:appH-KE-envelope-lower-bound}
\end{equation*}
The final inequality uses the explicit pseudoeffectivity premise;
$V_{\theta_t}$ is not an admissible nontrivial reference if that
class has no positive current.

For $s_j\downarrow0$, the bounded upper estimates allow the total-space
upper regularizations
\begin{equation*}
 U_N=\left(\sup_{j\geq N}u_{s_j}\right)^*,
 \qquad \Psi=\lim_{N\to\infty}U_N.
\tag{H.79}\label{eq:appH-KE-tail-upper-limit}
\end{equation*}
Each $U_N$ has background $\theta+s_N\omega$. The lower bound
in \eqref{eq:appH-KE-envelope-lower-bound} prevents the restriction
of the decreasing limit from being identically $-\infty$ on any
fibre. The limiting current is positive in the canonical class,
and the two fibrewise singularity comparisons give minimal
singularities. This explains the source's limiting step conditional
on the stated solution, volume and estimate inputs; it does not
prove those inputs or construct a probability measure at zero volume.
Without pseudoeffectivity, even the canonical bundle of a
projective-line fibre has negative degree and admits no positive
canonical current.

Such a metric would equip $(m-1)K$ with weight $(m-1)\Psi$.
For a pluricanonical section $\sigma$ the resulting local integrability
comparison is with $|\sigma|^{2/m}$. It is the weight
$e^{-(m-1)\Psi}$, rather than $e^{-\Psi}$ paired with the same
$(m-1)$st power of the reference metric, that belongs in the L2
extension argument. The suitable weakly pseudoconvex extension
input and the passage from a shrunken disk to the original disk
remain separate proof gates in this preprint.

The MMP antecedents also have different scopes.
\cite{HaconLiRao26AdjointV2} requires a projective central fibre
in its general-dimensional family statements; its threefold
specialization is separate. \cite{HaconXie26MMPV1} later claims
an extension without that projectivity, but its full inductive
proof chain is not certified by this discussion. Modified-big
MMP conclusions do not settle every pseudoeffective terminal-model,
flip-termination or ordinary K\"ahler abundance problem.
The exact source qualifications are in
\Cref{sec:provenance-KE-family-normalization}. Complex canonical
currents supply no repair of the non-Archimedean BJ22 Siu division
or general envelope-closure proof boundaries.

\section{Finite pullbacks and generic Segre invariants}
\label{sec:appH-finite-pullback-Segre}
\index{Segre number!finite pullback}
\index{Le number@L\^e number!generic and prepolar coordinates}

Analytic singularities connect ideal theory with local intersection
cycles. The comparison in \cite{JelonekPerezSantos26SegreV1}
is useful here because it separates three operations: extending
an ideal by a finite map, replacing that ideal by the Jacobian
of a reduced pullback, and choosing generic coordinate slices.
Its local numerical invariants differ from the global Segre
degrees of \cref{sec:appC-Segre-dependence}.

\subsection{The ambient ideal and the indexing}
\label{sec:appH-Segre-local-indexing}

Let $R=\mathbf C\{x_1,\ldots,x_n\}$, $n\geq2$, and
$0\ne Q\subsetneq R$. Generic linear combinations of generators
successively cut the ambient cycle. At each stage, retain
the components supported on $V(Q)$ and continue with the
residual components. The multiplicity at the origin of the
retained codimension-$j$ cycle is $e_j(Q)$, for $1\leq j\leq n$;
all scheme intersection coefficients are included.
\cite[\S2]{GaffneyGassler96SegreV1} supplies this generic-cycle
construction. Achilles--Manaresi's dimension-indexed sequence
and the Segre numbers satisfy
\[
 e_j(Q)=c_{n-j}(Q),\qquad e_n(Q)=c_0(Q)=j(Q).
\tag{H.80}\label{eq:appH-Segre-index-conversion}
\]
Here $j(Q)$ is j-multiplicity. For a maximal-ideal-primary $Q$,
only $c_0$ can be nonzero and it is the Hilbert--Samuel
multiplicity. The vanishing and highest-dimensional association
formula are \cite[Theorem~3]{AchillesManaresi06Author}.
CPRU use a \emph{codimension} index for their sequence, so
their tuple order must be converted before applying
\cite[Corollary~6.5]{CidRuizPoliniUlrich25Author}.

For a reduced holomorphic germ $f$ with a critical point at the
origin, use the ambient gradient
ideal $J(f)=(\partial f/\partial x_1,\ldots,\partial f/\partial x_n)$
in $R$. The generic L\^e numbers have the dictionary
\[
 \lambda^i_{\mathrm{gen}}(f)=e_{n-i}(J(f))=c_i(J(f)),
 \qquad 0\leq i\leq n-1.
\tag{H.81}\label{eq:appH-Segre-Le-dictionary}
\]
This is the generic construction recalled in
\cite[\S2.3]{JelonekPerezSantos26SegreV1}, with
\cite{GaffneyGassler96SegreV1,Massey15LeCyclesV2} as antecedents.
It does not use the gradient ideal in $R/(f)$.
Coordinates for which the successive polar cuts are proper
give defined L\^e numbers; prepolar transversality is one way
to arrange that condition. Generic linear coordinates impose
the further condition that their numbers are the constant
generic values.

\subsection{Finite degree and ideal inclusion}
\label{sec:appH-Segre-finite-degree}

Let $F:(\mathbf C^n,0)\to(\mathbf C^n,0)$ be finite of degree
$d$, giving smooth local rings $A\to B$ and a finite flat
extension. For a nonzero proper ideal $Q\subset A$,
\cite[Theorem~3.4]{JelonekPerezSantos26SegreV1} states
\[
 e_j(Q;A)\leq e_j(QB;B)\leq d\,e_j(Q;A),
 \qquad e_n(QB;B)=d\,e_n(Q;A).
\tag{H.82}\label{eq:appH-Segre-finite-bounds}
\]
The proof chooses compatible generic proper cuts and retains
the multiplicities of all split source components. Exact
degree scaling at the last index does not extend to every
coordinate: a divisor ideal pulled back under
$(u,v)\mapsto(u,v^d)$ can keep its first Segre number equal to one.

Now let $g,h$ be reduced irreducible equations, with critical
points at the origin, for the target hypersurface and the
reduced source support, and write
$g\circ F=u h^r$, with $u$ a unit. Absorb an $r$-th root of
$u$ into the equation $h$. The gradient ideal's integral
closure, and hence its Segre numbers, are invariant under
that unit change. The chain rule gives
\[
 h^{r-1}J(h)\subset J(g)B.
\tag{H.83}\label{eq:appH-Segre-chain-rule}
\]
For $r=1$, set $K=J(h)$ and $I=J(g)B$.
The ideal inclusion $K\subset I$ yields a lexicographic
comparison in the order $(e_1,\ldots,e_n)$, by the cited
CPRU result. Combining it with
\eqref{eq:appH-Segre-finite-bounds} gives
$e(J(h))\geq_{\mathrm{lex}}e(J(g))$.
A strict increase in an early coordinate leaves later
coordinates uncontrolled. Equality of radicals does not
remove that limitation.

If both critical loci are isolated, their Jacobian ideals
are maximal-ideal-primary. The source-attributed conclusions
are stronger and simpler:
\[
 \mu(h)\geq\mu(g)\quad\hbox{for every }r,
 \qquad \mu(h)\geq d\mu(g)\quad\hbox{when }r=1.
\tag{H.84}\label{eq:appH-Segre-isolated-Milnor}
\]
See \cite[Theorem~4.2, Corollary~4.3]{JelonekPerezSantos26SegreV1}
and the earlier isolated comparison \cite{Jelonek23MilnorV1}.
For higher $r$, the product ideal in
\eqref{eq:appH-Segre-chain-rule} has a divisorial contribution;
it is not a constant multiple of the sequence of $J(h)$.
The isolated comparison uses characteristic-zero Milnor-fibre
cohomology, or a localized Brieskorn-module trace route, rather
than that invalid rescaling.

\subsection{A coordinate excess and the unresolved comparison}
\label{sec:appH-Segre-generic-frontier}

For a one-dimensional critical locus, its transverse cycle is
$\Lambda^1_f=\sum_C\mu^\perp_C(f)[C]$. An admissible first
linear coordinate $\ell$ can intersect this cycle with excess
$\delta\geq0$ over its multiplicity at the origin. The
Euler-characteristic formula then gives
\[
 \lambda^1_{f,\ell}=\lambda^1_{\mathrm{gen}}(f)+\delta,
 \qquad
 \lambda^0_{f,\ell}=\lambda^0_{\mathrm{gen}}(f)+\delta.
\tag{H.85}\label{eq:appH-Segre-prepolar-excess}
\]
The convention is that the ambient dimension is $n$;
the alternating sign in the Milnor-fibre Euler characteristic
is $(-1)^{n-1-i}$. Massey's lecture notes use ambient $n+1$,
so their exponent is adjusted accordingly.

For instance, \cite[Example~7.4]{JelonekPerezSantos26SegreV1}
uses $f=(x-z^k)^2+y^3$, $k\geq2$. The generic pair is
$(\lambda^1,\lambda^0)=(2,0)$, while the prepolar first
coordinate $x$ gives $(2k,2k-2)$. Taking the identity map
and two such coordinate systems would invalidate a comparison
for arbitrary prepolar choices; it leaves the generic
comparison untouched.

The finite-degree factor also has a precise boundary. For
$g=x^2+y^3$ in three variables and
$F(u,v,w)=(u,v,w^d)$, the reduced pullback is $h=u^2+v^3$.
Both Segre tuples are $(0,2,0)$, although $d>1$.
The map's degree is entirely in the singular direction.
In contrast, the non-equimultiple example
$g=x^4+y^4-xyz$, $F(x,y,z)=(x,y,z^2)$ in the cited
Example~7.3 has tuples $(0,1,8)$ and $(0,1,24)$.
The transverse Milnor number is one on both sides;
the later coordinate increases. Its chosen first coordinate
has zero excess, which is why those computed L\^e numbers
are generic.

For multiplicity-one pullback and equal critical-locus
dimension $s$, an equimultiple \emph{target} Jacobian gives
componentwise comparison; equimultiplicity of the source is
not required. In singular dimension one, equality of
$c_1(K)$ and $c_1(I)$ forces
$c_0(K)\geq d c_0(J(g))$. These are
\cite[Theorems~5.3 and~5.5]{JelonekPerezSantos26SegreV1}.
For higher multiplicity, comparison of transverse Milnor
numbers requires compatible smooth finite slices, and the
passage to a generic L\^e number requires the additional
pulled-cut multiplicity equality in Corollary~6.2.
An arbitrary nongeneric cut does not supply that equality.

Under the irreducibility assumptions above, the paper's
Conjectures~8.1 and~8.2 ask for
\[
 e_j(J(h))\geq e_j(J(g))\ (1\leq j\leq n),
 \quad\Longleftrightarrow\quad
 \lambda^i_{\mathrm{gen}}(h)\geq
 \lambda^i_{\mathrm{gen}}(g)\ (0\leq i<n).
\tag{H.86}\label{eq:appH-Segre-componentwise-frontier}
\]
They are two formulations of one open problem, including
arbitrary $r$. Conjecture~8.3 asks for the weaker
lexicographic comparison; it is known here for $r=1$ and
for isolated singularities at every $r$.
The strict-first-increase case leaves the bottom coordinate
open even for one-dimensional critical loci. The large-power
isolated-perturbation formula in
\cite[Theorem~3.5]{Massey24MinkowskiV1} concerns one germ;
it does not automatically compare target and source after
pulling a linear coordinate back to a possibly critical
nonlinear function.

These are local complex and algebraic comparison questions.
They neither transfer the ample global criterion to a nef
b-divisor nor supply the critical homogeneous Siu division
input of the non-Archimedean envelope argument.
\Cref{sec:appL-Segre-source-gates} records the exact source
versions and the trace and slicing qualifications.

\section{Determinants of positive curvature and Lefschetz failure}
\label{sec:appH-Nakano-determinants}
\index{Nakano positivity!determinant Lefschetz failure}
\index{dual Nakano positivity!pointwise determinant}
\index{simultaneous diagonalizability!form coordinates}
\index{Hodge--Riemann relations!pointwise curvature determinants}
\index{Lorentzian polynomial!strictness and higher degrees}

\Cref{sec:appH-Griffiths-Chern} distinguishes the positivity of a
curvature tensor from that of its top Chern representative. There is
a further distinction even under the stronger Nakano hypotheses:
a positive characteristic form need not have all Lefschetz and
Hodge--Riemann properties. The pointwise examples and positive
degree-one result below are attributed to
\cite{Chen26NakanoDeterminantsV2}. They also explain why the
Lorentzian conclusions of \Cref{sec:appG-augmented-polymatroid-volume}
must be distinguished from a complete Hodge package.

\subsection{Constant forms and the pointwise cone}
\label{sec:appH-Nakano-pointwise-cone}

Work over $\mathbf C^n$ and write
$V^{p,q}=\Lambda^p(\mathbf C^n)^*\otimes
\Lambda^q\overline{(\mathbf C^n)^*}$ for constant complex forms.
Fix $1\leq k\leq n$ and a positive definite real $(1,1)$-form
$\omega$. For a real $(k,k)$-form $\Omega$ and $p+q=n-k$, put
\begin{equation*}
 L^{p,q}_{\Omega}:V^{p,q}\longrightarrow V^{n-q,n-p},
 \qquad \alpha\longmapsto\alpha\wedge\Omega.
 \tag{H.87}\label{eq:appH-Nakano-Lefschetz-map}
\end{equation*}
Source Definition~2.1 uses the following path formulation.

\begin{definition}
\label{def:appH-Nakano-HR-path}
The form $\Omega$ has the pointwise Hodge--Riemann property
in bidegree $(p,q)$ relative to $\omega$ if there is a continuous
path of real $(k,k)$-forms joining $\Omega$ to $\omega^k$ such that
\[
 L^{p-j,q-j}_{\Omega_t\wedge\omega^{2j}}
 \text{ is an isomorphism for every }t\in[0,1]
 \text{ and }0\leq j\leq\min(p,q).
\]
The pointwise Hodge--Riemann cone consists of forms satisfying
this condition in every complementary bidegree.
\end{definition}

The mixed linear Hodge theorem and constancy of signature along
the path give the corresponding positive primitive Hermitian
forms. A nonzero kernel of \eqref{eq:appH-Nakano-Lefschetz-map}
already prevents membership, independently of the reference
$\omega$. These are statements about constant forms and their
full exterior spaces, rather than a quotient cohomology algebra.

Let $M$ be a Hermitian $k\times k$ matrix of constant $(1,1)$-forms,
written
\[
 M_{a\bar b}=\frac{\rmi}{2}
       \sum_{\mu,\nu=1}^n A_{a\bar b\mu\bar\nu}
                dz_\mu\wedge d\bar z_\nu.
\]
Griffiths positivity tests $A$ on nonzero decomposable tensors
in $\mathbf C^k\otimes\mathbf C^n$. Nakano positivity means that
the full Hermitian coefficient matrix on that tensor space is
positive definite. Dual Nakano positivity tests its partial
transpose in the fibre indices. These two stronger tests are
distinct. The entries of $M$ have even exterior degree, so they
commute and its determinant uses the ordinary permutation formula.

In this section, \emph{simultaneous diagonalizability} means that
one complex linear coordinate system gives
\begin{equation*}
 M=\sum_{i=1}^n B_iV_i,\qquad
 V_i=\frac{\rmi}{2}dz_i\wedge d\bar z_i,\qquad B_i=B_i^*>0.
 \tag{H.88}\label{eq:appH-Nakano-SD}
\end{equation*}
It diagonalizes the differential-form coordinates, not the
coefficient matrices $B_i$. No commutativity of those matrices
is required. In this setting all three curvature positivity
tests are equivalent to $B_i>0$ for every $i$.

\subsection{A kernel under Nakano positivity}
\label{sec:appH-Nakano-unrestricted}

The unrestricted counterexample is
\cite[Theorem~A and \S3]{Chen26NakanoDeterminantsV2}.

\begin{theorem}
\label{thm:appH-Nakano-Lefschetz-failure}
For every $n\geq4$ and $2\leq k\leq n-2$, there is a
Nakano-positive Hermitian $k\times k$ matrix $M$ of constant
$(1,1)$-forms on $\mathbf C^n$ such that
$L^{1,n-k-1}_{\det M}$ is singular. The transposed matrix is
dual Nakano positive and has the same determinant.
\end{theorem}

The two positivity assertions refer to separately obtained
matrices. They do not assert simultaneous Nakano and dual
Nakano positivity. Finski's stronger-positivity theorem supplies
strict weak positivity of the determinant in both cases;
see \cite[Theorem~1.1]{Finski22CharacteristicPublished}.
It supplies no isomorphism of the full exterior Lefschetz map.

The exact construction makes this distinction visible.

\begin{example}
\label{ex:appH-Nakano-explicit-kernel}
First take $n=k+2$, put $\omega_0=\sum_iV_i$, and let
$E_{ij}=(\rmi/2)dz_i\wedge d\bar z_j$. On the first four
coordinates set
\begin{equation*}
 N=\begin{pmatrix}
 V_1+V_3&E_{12}+E_{34}\\
 E_{21}+E_{43}&V_2+V_4
 \end{pmatrix},\qquad
 M_t=\operatorname{diag}(\omega_0I_2+tN,\omega_0I_{k-2}).
 \tag{H.89}\label{eq:appH-Nakano-seed-matrix}
\end{equation*}
The full Nakano matrix is
$I_{kn}+t(w_1w_1^*+w_2w_2^*)$, where
$w_1=e_1\otimes f_1+e_2\otimes f_2$ and
$w_2=e_1\otimes f_3+e_2\otimes f_4$. It is positive
definite for $t\geq0$.

Put $\omega_U=V_1+\cdots+V_4$,
$\omega_W=\sum_{i=5}^{k+2}V_i$, $\Delta=\det N$ and
$\gamma=E_{13}-E_{31}+E_{24}-E_{42}\ne0$.
The source's exterior identities give
\begin{equation*}
 \begin{aligned}
 \gamma\wedge\Delta&=-\tfrac12\gamma\wedge\omega_U^2,
 &\gamma\wedge\omega_U^3&=0,\\
 \gamma\wedge\det M_t
   &=\left[\binom{k}{2}+(k-1)t-\tfrac12t^2\right]
       \gamma\wedge\omega_U^2\wedge\omega_W^{k-2}.
 \end{aligned}
 \tag{H.90}\label{eq:appH-Nakano-kernel-polynomial}
\end{equation*}
Thus $t_k=k-1+\sqrt{(k-1)(2k-1)}$ gives a nonzero
$(1,1)$-kernel. At $k=2$, interpret $\omega_W^0=1$.
The partial-transpose spectrum of the same tensor is
$1+t$ with multiplicity six, $1-t$ with multiplicity two,
and $1$ with multiplicity $kn-8$. Since $t_k>1$, this
displayed matrix is not dual Nakano positive.
\end{example}

For $n>k+2$, the source's Lemma~2.3 extends the coefficient
tensor by positive identity blocks and wedges $\gamma$ with
$d\bar z_{k+3}\wedge\cdots\wedge d\bar z_n$.
Every new determinant term contains some $V_i$ in those
directions and annihilates this form. Its bidegree is exactly
$(1,n-k-1)$. This explains the ambient-dimension range,
rather than extrapolating it from finitely many computations.

\subsection{What coordinate diagonalization does prove}
\label{sec:appH-Nakano-SD-degree-one}

Under \eqref{eq:appH-Nakano-SD}, square-zero relations
$V_i^2=0$ give
\begin{equation*}
 \det M=\sum_{|I|=k}d_I V_I,\qquad
 d_I=k!D_k(B_i:i\in I)>0,
 \qquad D_k(B,\ldots,B)=\det B.
 \tag{H.91}\label{eq:appH-Nakano-mixed-coefficients}
\end{equation*}
Here $D_k$ is the symmetric polarization with the $1/k!$
normalization. Positivity of all displayed coefficients alone
does not settle the Lefschetz maps.

The precise positive result, and its higher-degree boundary,
are \cite[Theorems~B--C]{Chen26NakanoDeterminantsV2}.

\begin{theorem}
\label{thm:appH-Nakano-SD-boundary}
Let $n\geq4$ and $2\leq k\leq n-2$, and suppose that
$M$ satisfies \eqref{eq:appH-Nakano-SD}.
For every positive definite reference form $\omega$, its
determinant has the pointwise Hodge--Riemann property in every
bidegree $(p,q)$ with $p+q=n-k$ and $\min(p,q)\leq1$.
In particular it has that property in all complementary
bidegrees when $n-k=2$ or $3$.

For every $n\geq6$ and $2\leq k\leq n-4$, there also exist
matrices satisfying \eqref{eq:appH-Nakano-SD} for which
$L^{2,n-k-2}_{\det M}$ is singular. Such matrices are
simultaneously Nakano and dual Nakano positive.
\end{theorem}

The positive assertion uses a strict equality theorem, not
merely a bound of one positive eigenvalue. The source's
Proposition~4.1 decomposes a Lefschetz map into blocks. For
$J\subseteq\{1,\ldots,n\}$ of cardinality $k+2r$, the
block indexed by $r$-subsets of $J$ has entries
\begin{equation*}
 (H_r^J)_{R,S}=\begin{cases}
 d_{J\setminus(R\cup S)},&R\cap S=\varnothing,\\
 0,&R\cap S\ne\varnothing.
 \end{cases}
 \tag{H.92}\label{eq:appH-Nakano-block-matrix}
\end{equation*}
The pure blocks use $d_I>0$. For $r=1$, Proposition~4.4
gives signature $(1,k+1)$.

To obtain that strict signature, Lemma~4.2 writes all
$B_i=\sum_\rho a_{\rho i}u_\rho u_\rho^*$ with a common
finite spanning frame and \emph{all} $a_{\rho i}>0$.
Cauchy--Binet gives a nonzero multiaffine stable polynomial
$s(y)=\det\sum_\rho y_\rho u_\rho u_\rho^*$ with
nonnegative coefficients. Its multiaffine reversal is stable
as well, so \cite[Proposition~2.2]{BrandenHuh20Lorentzian}
makes $s$ dually Lorentzian; its factorial reversal
normalization is the identity.

With $m=k+2$, set
$A_\rho=\operatorname{diag}(a_{\rho i}:i\in J)>0$
and $X(x)=\operatorname{diag}(x_i:i\in J)$.
The source identifies
\[
 m!D_m\bigl(X(x),X(y),s(A_\rho)\bigr)=x^TH_1^Jy,
\]
where the mixed discriminant is extended linearly in the
polynomial $s$. The strict equality characterization of
\cite[Theorem~8.2]{RossSuessWannerer25DuallyPublished}
and positivity on $I_m$ make its orthogonal hyperplane
negative definite. This excludes a kernel. Lemma~2.2 of Chen connects
the determinant to $\omega^k$ through nonsingular blocks;
the reference form need not be diagonal in the coordinates
of \eqref{eq:appH-Nakano-SD}. The classical mixed linear
Hodge theorem and the cone-Lorentzian foundations of the
cited equality result remain imported inputs.

For comparison, the Lorentzian quadratic
\[
 2x_1x_2+2x_3x_4+(x_1+x_2)(x_3+x_4)
\]
has Hessian eigenvalues $4,0,-2,-2$. Its one-positive-eigenvalue
condition does not imply nondegeneracy. Nor does the strict
$r=1$ argument prove anything about every $r\geq2$ block.

\subsection{A simultaneously positive higher-degree kernel}
\label{sec:appH-Nakano-SD-higher-kernel}

The rank-two, six-dimensional construction in \S5 of Chen
is an exact realization of a degeneration already present in
\cite[Example~9.2, physical p.~37]{RossToma21SchurHRV3}.

\begin{example}
\label{ex:appH-Nakano-SD-six-dimensional}
On $\mathbf C^6$ put $W_j=V_{2j-1}+V_{2j}$ and
\[
 S_1=\begin{pmatrix}1&0\\0&-1\end{pmatrix},\quad
 S_2=\tfrac12\begin{pmatrix}-1&\sqrt3\\\sqrt3&1\end{pmatrix},\quad
 S_3=\tfrac12\begin{pmatrix}-1&-\sqrt3\\-\sqrt3&1\end{pmatrix}.
\]
Each $S_j$ is Hermitian with eigenvalues $1,-1$; they do
not all commute. Set
\begin{equation*}
 \begin{aligned}
 \widehat M_t&=\sum_{j=1}^3(I_2+\sqrt t\,S_j)W_j,
 \quad 0\leq t\leq\tfrac12,\\
 \Omega_t&=\det\widehat M_t=(1-t)\Sigma+(2+t)\Pi,
 \end{aligned}
 \tag{H.93}\label{eq:appH-Nakano-SD-seed}
\end{equation*}
where $\Sigma=\sum_jW_j^2$ and
$\Pi=\sum_{i<j}W_iW_j$. Every coefficient matrix is
positive definite, with eigenvalues $1\pm\sqrt t$.

The relations $W_j^3=0$ yield, for
$u=(2+t)/(2(1-t))$,
\begin{equation*}
 \Omega_t\wedge(\Pi-u\Sigma)
       =(1-t)(1+4u-2u^2)\Sigma\Pi.
 \tag{H.94}\label{eq:appH-Nakano-SD-explicit-kernel}
\end{equation*}
At $t_*=\sqrt6-2\in(0,1/2)$ and
$u_*=1+\sqrt6/2$, the nonzero $(2,2)$-form
$\Pi-u_*\Sigma$ is a kernel vector.
The fifteen-by-fifteen block $H_2^{\{1,\ldots,6\}}$
has determinant
\begin{equation*}
 2^{15}3^6(1-t)^3(2t+1)^3(t^2+t+1)^2(t^2+4t-2).
 \tag{H.95}\label{eq:appH-Nakano-SD-determinant-factor}
\end{equation*}
Its endpoint signs at $0$ and $1/2$ are opposite.
\end{example}

For higher rank $k=m+2$, the source first studies the
normalized product $\Omega_t\omega_0^m/m!$ in dimension
$m+6$. Its block $K_m$ is indexed by two-subsets; let
$(L_m)_{R,i}=1$ when $i\in R$ and zero otherwise.
After putting old pairs first, the exact counting identities are
\begin{equation*}
 \begin{aligned}
 K_{m+1}&=\begin{pmatrix}
 K_m&K_mL_m/(m+1)\\
 L_m^TK_m/(m+1)&0
 \end{pmatrix},\\
 L_m^TK_mL_m&=(m+1)(m+2)R_m.
 \end{aligned}
 \tag{H.96}\label{eq:appH-Nakano-SD-extension-block}
\end{equation*}
Here $R_m$ has zero diagonal. Its off-diagonal entries are
$12$ within an original coordinate pair, $12+3t$ between
original pairs, $4(5+t)$ between old and new indices,
and $6(5+t)$ between new indices.
It has signature $(1,m+5)$ on the whole parameter interval.
A determinant-one congruence changes the displayed block to
$\operatorname{diag}(K_m,-(m+2)R_m/(m+1))$, so both
endpoint signs reverse at each rank extension.

The unperturbed extra coefficient directions in the first
two-by-two block are zero. The source therefore adds
$\varepsilon\omega_{\mathrm{new}}I_2$ there, with
$\varepsilon>0$, before claiming strict positivity.
For small fixed $\varepsilon$ the two endpoint signs persist;
continuity gives a singular intermediate block. Lemma~2.3
then extends the ambient dimension. This perturbation and
the general counting argument are essential parts of the
source proof, not conclusions from a bounded rank loop.

Ross--Toma's example uses the rank-three split bundle
$E=\bigoplus_{j=1}^3\pi_j^*\mathcal O_{\mathbf P^2}(1)$
on $(\mathbf P^2)^3$ and its \emph{real-twisted} ample
bundle $E\langle\tau\det E\rangle$, $\tau>0$.
Its second Chern class is
$\Pi+(2\tau+3\tau^2)(\Sigma+2\Pi)$ in the cube-zero
hyperplane-class algebra. At the seed parameter,
$(1-t_*)/(3t_*)=1/\sqrt6$, so taking
$2\tau+3\tau^2=1/\sqrt6$ matches the class up to a
positive scalar. This attribution does not identify Chen's
rank-two constant tensor with a global bundle curvature.

\subsection{Boundary questions and comparison with October results}
\label{sec:appH-Nakano-boundary-questions}

When $n-k=0$ or $1$, membership in the pointwise cone is
equivalent to strict weak positivity of the determinant.
Finski's theorem therefore settles these cases under Nakano
or dual Nakano positivity. Under Griffiths positivity,
\cite[Question~6.1]{Chen26NakanoDeterminantsV2} asks that
positivity for $n\geq k\geq4$; plane restriction reduces
it to $n=k$.

The controlling Chen version is dated 28 September 2026.
It must be read together with the later
\cite[Theorem~1.1]{Du26GriffithsTopChernV1}, dated
5 October 2026. The negative top Chern jets in
\Cref{sec:appH-Chern-global} give a negative answer for
every $k\geq9$. Adding positive identity coefficients in
additional tangent directions retains Griffiths positivity
and the negative restriction to the original $k$-plane.
Thus these ranks are not still open in the two boundary
dimensions either. The previously recalled positive ranks
two and three leave four through eight as the remaining
source-scoped candidate ranks; no exhaustive current-open
certification is asserted.

Question~6.2 of Chen asks for a Hermitian two-by-two
constant-form matrix on $\mathbf C^4$ that is simultaneously
Nakano and dual Nakano positive and has singular
$L^{1,1}_{\det M}$. The first example above fails the
dual test. The six-dimensional example concerns $(2,2)$,
not $(1,1)$. The positive diagonalization theorem excludes
all matrices satisfying \eqref{eq:appH-Nakano-SD} from this
four-dimensional question. It remains a source-era local
question in this account, separate from the canonical
non-Archimedean problem inventory.

All these statements retain their pointwise or explicitly
real-twisted cohomological categories. Positive Schur forms,
numerical positivity, degree-one signatures and higher-degree
Lefschetz maps are different properties. No new
non-Archimedean characteristic product, general Siu division
or envelope closure follows. Exact sources, actual reading
and computational scopes are recorded in
\Cref{sec:appL-Nakano-determinant-sources}.


\section{Unit tests and failure modes}
\label{sec:appH-unit-tests}

The following checks detect most convention errors.

\subsection*{Constants}
\label{sec:appH-constants}

If \(u\) is finite energy and \(c\in\mathbf R\), then
\[
\begin{aligned}
 \operatorname{MA}_\theta(u+c)&=\operatorname{MA}_\theta(u),\\
 M_\theta(u+c)&=M_\theta(u),\\
 E_\phi(u+c)&=E_\phi(u)+cM_\theta(\phi).
\end{aligned}
\]
Any formula giving a different energy coefficient has silently normalized
the mass.

\subsection*{One logarithmic pole}
\label{sec:appH-one-pole}

On a disk,
\[
                         u=\lambda\log|z|^2+O(1)
\]
with \(\lambda\geq0\) has Lelong number \(\lambda\), and
\[
                         \mathcal I(u)
          =\mathcal O(-\lfloor\lambda\rfloor\{0\}).
\]
This simultaneously tests the factor in \(\ddc\), the exponent in the
multiplier ideal, and the inclusion direction in
\(\preceq_{\mathcal I}\).

\subsection*{Minimal singularities}
\label{sec:appH-minimal-type}

For \(u=V_\theta+c\),
\[
 P_\theta[u]=V_\theta,\qquad
 M_\theta(u)=M_\theta(V_\theta).
\]
When \(\{\theta\}\) is big, the latter mass is the volume
\(\operatorname{vol}(\{\theta\})\).

\subsection*{Zero mass}
\label{sec:appH-zero-mass}

If \(M_\theta(u)=0\), the following expressions are not defined without
additional choices:
\[
\begin{gathered}
 M_\theta(u)^{-1}\operatorname{MA}_\theta(u),\\
  \text{the maximality characterization in \Cref{thm:appH-p-projection}},\qquad
 \mathcal E^1(X,\theta,P_\theta[u]).
\end{gathered}
\]
The perturbative inverse systems of \cref{part:complex} were designed precisely to
avoid assigning fictitious values to these expressions.

\subsection*{Decreasing limits}
\label{sec:appH-decreasing-limits}

From \(u_j\searrow u\) one may infer pointwise and \(L^1\) convergence under
the usual noncollapse hypothesis. One may not infer
\[
                         M_\theta(u_j)\to M_\theta(u)
\]
or weak convergence of \(\operatorname{MA}_\theta(u_j)\) unless a
full-mass, capacity, or strong-energy theorem is invoked.

\section{Source map}
\label{sec:appH-source-map}

The detailed treatments in the supplied book \cite{XiaBook} are organized
as follows:
\begin{center}
\begin{tabular}{@{}p{0.43\textwidth}p{0.47\textwidth}@{}}
\toprule
topic used here & chapter of the supplied book\\
\midrule
psh functions, Lelong numbers, multiplier ideals & Plurisubharmonic functions\\
Bedford--Taylor and nonpluripolar products & Non-pluripolar products\\
\(P\)- and \(\mathcal I\)-envelopes & The envelope operators\\
finite-energy geodesics and ray metrics & Geodesic rays in the space of potentials\\
\(P\)- and \(\mathcal I\)-preorders, \(d_S\) & Comparison of singularities\\
\(\mathcal I\)-good classes and mixed volumes & \(\mathcal I\)-good singularities\\
restriction and trace & The trace operator\\
test curves and Legendre correspondence & Test curves\\
\bottomrule
\end{tabular}
\end{center}

The foundational construction of global nonpluripolar products is
\cite{BEGZ10}; relative finite-energy metric geometry is
\cite[\S4.3]{XiaBook}, extending the K\"ahler metric of \cite{Dar15};
the distinct metric on singularity types is \cite{DDNLmetric}; and the test-curve correspondence is
\cite{RWN14}. This appendix follows the conventions of the supplied book,
even when another source uses \(\ddc=\rmi\partial\bar\partial/\pi\), the
weight \(-\log|s|_h\) instead of \(-\log|s|_h^2\), or probability-normalized
Monge--Amp\`ere measures.
